\medterm{Eagle Syndrome} A rare cause of throat, face, jaw, or neck pain related to an unusually long bony projection near the skull base or a hardened nearby ligament. Pain may worsen with swallowing or turning the head.

\medterm{Eales Disease} An uncommon inflammatory disorder of small retinal veins, usually affecting otherwise healthy young adults. Blocked vessels can cause abnormal new vessels and bleeding inside the eye, leading to floaters or reduced vision.

\medterm{Ear} The organ for hearing and balance. The outer ear collects sound, the middle ear transmits vibrations through the eardrum and tiny bones, and the inner ear converts sound and movement into signals sent to the brain.

\medterm{Earache} Pain felt in the ear or nearby area. It may arise inside the ear or be referred from the teeth, jaw, throat, or neck; identifying the source requires an examination when pain persists or is severe.

\synonyms
Otalgia

\medterm{Ear Barotrauma} Ear injury caused by unequal pressure between the middle ear and surrounding air or water, often during flying or diving. It may cause pain, fullness, blocked hearing, ringing, dizziness, bleeding, or a torn eardrum.

\medterm{Ear Canal} The passage between the outer ear and eardrum. Its skin and earwax protect the middle ear, but infection, injury, foreign objects, or wax buildup can cause pain, discharge, or reduced hearing.

\medterm{Ear Cartilage} Flexible connective tissue that forms much of the outer ear and gives it shape. Injury or infection can cause pain, swelling, deformity, or damage that permanently changes the ear.

\medterm{Ear Disease} Any disorder of the outer, middle, or inner ear. It may cause pain, discharge, hearing loss, ringing, dizziness, infection, or balance problems; diagnosis requires examination of the ear and sometimes hearing or imaging tests.

\medterm{Ear Diseases} Disorders affecting the outer ear, ear canal, eardrum, middle-ear bones, pressure-equalizing tube, cochlea, or balance organs. They include infection, injury, blockage, inflammation, hearing disorders, and growths.

\medterm{Ear, Diseases Of} Disorders of the outer, middle, or inner ear. Symptoms, examination findings, causes, and treatment depend on the affected structure and may include pain, discharge, hearing loss, ringing, dizziness, or imbalance.

\medterm{Ear Drainage Culture} A laboratory test of fluid draining from the ear to identify bacteria or fungi. Results help guide treatment but must be interpreted with the examination because normal organisms, contamination, or colonization can affect the result.

\medterm{Eardrum} A thin membrane separating the ear canal from the middle ear. Sound makes it vibrate, and the vibration moves the tiny middle-ear bones; infection, pressure, or injury can tear it or reduce hearing.

\medterm{Eardrum Repair} Treatment to close a hole or repair damage in the eardrum, often using a tissue graft or other surgical technique. It may improve hearing and prevent infection, but healing is not always complete.

\medterm{Eardrum, Ruptured} Alternate name; see \textit{Perforated Eardrum}.

\medterm{Ear Examination} Inspection of the outer ear and ear canal and viewing of the eardrum, usually with an otoscope. Hearing, balance, or additional tests may be added when symptoms suggest middle- or inner-ear disease.

\medterm{Ear Foreign Body} An object lodged in the ear canal, most often in a child. It may cause pain, reduced hearing, bleeding, discharge, or injury; removal should be done with direct visualization and suitable instruments.

\synonyms
Auditory Canal Foreign Body; Aural Foreign Body

\medterm{Ear Infection} Infection or inflammation affecting part of the ear, including the outer ear, ear canal, middle ear, or inner ear. It may cause ear pain, fever, discharge, reduced hearing, dizziness, or irritability in children.

\medterm{Ear Infection, Outer Ear} Alternate name; see \textit{Swimmer's Ear}.

\medterm{Earliest Signs Of Increased Intracranial Pressure} Early warning signs may include worsening headache, vomiting, blurred or double vision, drowsiness, slowed thinking, or behavior change. Causes include bleeding, swelling, infection, a growth, or excess brain fluid; progressive symptoms need urgent assessment.

\medterm{Earlobes} The soft lower parts of the outer ears, made mainly of skin and fatty or connective tissue without cartilage. They may be affected by piercing injury, infection, cysts, allergic reactions, or tears.

\medterm{Early Afterdepolarization} An extra abnormal electrical signal during the heart’s recovery phase after each beat. It is more likely when the QT interval is prolonged and can trigger dangerous abnormal rhythms.
\synonyms
EAD; Early After-Depolarization

\medterm{Early Childhood Tooth Decay} Early-childhood tooth decay is bacterial acid destruction of primary teeth, promoted by frequent sugars, poor brushing, enamel defects, and prolonged bottle or nursing exposure.

\medterm{Early Decelerations} A gradual temporary slowing of the fetal heart rate that begins and ends with a uterine contraction, usually from pressure on the baby’s head. When isolated, it is generally reassuring, but the complete tracing and clinical situation matter.

\medterm{Early Detection} Finding a disease at an early stage, before it causes major complications or spreads. It may result from screening or prompt evaluation of symptoms and can improve treatment options for some conditions.

\medterm{Early Diagnosis} Identifying a disease before it has progressed or caused major complications. Earlier diagnosis may provide more treatment choices, but its benefit depends on the disease and the accuracy and usefulness of available tests.

\medterm{Early Endosome} A small membrane-bound compartment inside a cell that receives material brought in from the cell surface. It sorts that material for return to the surface, delivery elsewhere in the cell, or breakdown.

\medterm{Early Gastric Cancer} Stomach cancer limited to the inner lining or the layer just beneath it, whether or not nearby lymph nodes are involved. It may be found early because of screening or testing for digestive symptoms.

\synonyms
EGC

\medterm{Early Infantile Epileptic Encephalopathy} A group of severe seizure disorders beginning in infancy that can disrupt brain development. Causes may be genetic, structural, metabolic, or unknown; frequent seizures and developmental concerns require specialist evaluation.

\medterm{Early Intervention} Support provided to infants and young children with developmental delays or disabilities. It may include physical, occupational, speech, behavioral, educational, and family services to improve development, function, and participation.

\medterm{Early Puberty} Alternate name; see \textit{Precocious Puberty}.

\medterm{Early Repolarization Pattern} An electrocardiogram pattern with a small rise or notch near the end of the QRS wave, often in the lower or side leads. It is often harmless, but its meaning depends on symptoms, family history, and the rest of the tracing.
\synonyms
Early Repolarization ECG Pattern; BER

\medterm{Early Repolarization Syndrome} A rare heart-rhythm condition associated with a characteristic ECG pattern and, in some people, an increased risk of dangerous ventricular arrhythmias. Risk depends on the pattern, symptoms, and personal or family history.
\synonyms
ERS; Malignant Early Repolarization

\medterm{Early Satiety} Feeling full sooner than expected while eating, sometimes after only a small amount of food. It may result from stomach or digestive disease and should be assessed if persistent or associated with weight loss.

\medterm{Early-Stage Prostate Cancer} Prostate cancer that remains confined to the prostate or nearby tissues and has not spread to distant organs. The term includes cancers with different grades and risks of progression.

\synonyms
Early Prostatic Carcinoma; Localized Prostate Cancer

\medterm{Ear Microbiome} The community of microorganisms living on or in the ear, especially the ear canal. They interact with skin and earwax and may influence infection risk; moisture, cleaning, hearing devices, inflammation, and illness can change the community.

\medterm{Ear Neoplasm} An abnormal growth in the outer, middle, or inner ear that may be benign or cancerous. It can cause a persistent lump, pain, discharge, hearing loss, dizziness, or facial weakness and may require imaging or biopsy.

\medterm{Ear, Nose, And Throat} The medical specialty caring for disorders of the ears, nose, throat, voice, head, and neck. Specialists diagnose and treat problems such as hearing loss, sinus disease, swallowing or voice disorders, and head and neck tumors.

\medterm{Ear Speculum} A small cone-shaped tip fitted to an otoscope to direct light into the external ear canal and permit examination of the canal and tympanic membrane. Disposable and reusable types are available in different sizes.

\medterm{Ear Tag} A small, skin-covered projection on or near the outer ear, usually present from birth. It is generally harmless, but examination may be appropriate if it causes problems or occurs with other birth findings.

\medterm{Ear Thermometer} A device that estimates body temperature by measuring heat in the ear canal near the eardrum. Correct positioning is important; earwax, a small ear canal, or an ear infection can affect the reading.

\medterm{Ear Tube Insertion} A procedure placing a small ventilation tube through the eardrum to allow air into the middle ear and fluid to drain. It may help repeated infections or persistent fluid but can cause drainage, blockage, scarring, or a persistent opening.

\medterm{Ear Wax} A mixture of gland secretions and shed skin that protects and cleans the ear canal. A buildup can cause reduced hearing, fullness, ringing, itching, or discomfort; cotton swabs and inserted objects can worsen blockage or injure the ear.

\synonyms
Cerumen; Earwax

\medterm{Earwax} A waxy mixture made in the outer ear canal that lubricates and protects the skin and traps dust and microorganisms. It usually clears naturally, but buildup can block hearing or cause discomfort.

\medterm{Earwax Blockage} A buildup of wax that partly or completely blocks the ear canal. It may cause reduced hearing, fullness, ringing, itching, pain, or cough; safe removal may require a clinician.

\medterm{Earwax Removal} Removal of excess or impacted earwax by a clinician after examining the ear. Instruments, gentle suction, or irrigation may be used when safe; inserting objects at home can cause injury or push wax deeper.

\medterm{Eating Disorder} A group of mental-health conditions involving harmful eating patterns that cause distress, medical risk, or impaired functioning. Some involve intense concern about food, weight, or body shape, but avoidant/restrictive food intake disorder (ARFID) does not require body-image concerns. Examples include anorexia nervosa, bulimia nervosa, binge-eating disorder, and ARFID.

\medterm{Eating Disorders} Alternate name; see \textit{Eating Disorder}.

\medterm{Eating Disorders, Anorexia} Alternate name; see \textit{Anorexia Nervosa}.

\medterm{Eating Disorders, Binge Eating} Alternate name; see \textit{Binge-Eating Disorder}.

\medterm{Eating Disorders, Bulimia} Alternate name; see \textit{Bulimia Nervosa}.

\medterm{Eaton-Lambert Syndrome} An uncommon condition in which the immune system interferes with nerve signals to muscles. It causes weakness, often in the hips and shoulders, reduced reflexes, dry mouth, and other automatic-body problems. It can be linked to small-cell lung cancer.
\synonyms
Lambert-Eaton Myasthenic Syndrome; LEMS

\medterm{EBCT} Electron-beam computed tomography, an older rapid CT method that was used especially to measure calcium in the coronary arteries. Newer CT technology is now more commonly used for heart and vessel imaging.

\synonyms
Electron-Beam Computed Tomography; Ultrafast CT

\medterm{Ebola} Alternate name; see \textit{Ebola virus disease}.

\medterm{Ebola Hemorrhagic Fever} Alternate name; see \textit{Ebola Virus Disease}.

\medterm{Ebola Virus} A virus that causes Ebola disease. It spreads mainly through direct contact with infected blood or body fluids, usually after symptoms begin; illness may include fever, vomiting, diarrhea, bleeding, shock, and organ failure.

\medterm{Ebolavirus} A group of viruses that includes those causing Ebola disease. Infection can produce severe fever, vomiting, diarrhea, bleeding, shock, and organ damage and spreads through contact with infected body fluids.

\medterm{Ebola Virus Disease} A severe, often fatal viral disease caused by Ebola viruses and affecting humans and other primates. After an incubation period of about 2--21 days, it may begin with fever, tiredness, muscle pain, headache, or sore throat and progress to vomiting, diarrhea, rash, bleeding, shock, or organ failure. It spreads mainly when blood or other body fluids from an infected person or animal, or contaminated objects such as clothing or bedding, enter another person through broken skin or the eyes, nose, or mouth. People generally cannot spread the disease before symptoms begin, but can spread it after becoming ill.

\medterm{Ebstein Anomaly} A birth defect in which the tricuspid valve is positioned too low in the right ventricle and does not close normally. It can cause leakage of blood, enlargement of the right heart, blue coloring, heart failure, or abnormal rhythms.

\medterm{Ebstein's Anomaly} A birth defect in which the tricuspid valve is displaced downward and leaks. Effects range from no symptoms to blue coloring, heart failure, enlargement of the right heart, or abnormal heart rhythms.

\medterm{EBUS-TBNA} A procedure that uses a flexible breathing-tube camera and ultrasound to locate lymph nodes or masses beside the airways, then obtains tissue with a needle. It helps diagnose and stage lung cancer, sarcoidosis, lymphoma, and tuberculosis.

\synonyms
Endobronchial Ultrasound Transbronchial Needle Aspiration; Convex Probe EBUS

\medterm{Ecbolics} Medicines that stimulate contractions of the uterus. They may be used to start or strengthen labor or control heavy bleeding after birth; dosing requires medical supervision because excessive contractions can harm the parent or baby.

\medterm{Eccentric} Positioned away from the center. In exercise, an eccentric muscle action occurs when a muscle produces force while lengthening, such as lowering a weight; it can improve strength but may cause soreness or injury.

\medterm{Eccentric Action} Muscle contraction that produces force while the muscle lengthens, such as the controlled lowering of an object. It helps control movement and build strength but may cause muscle soreness when new or intense.

\synonyms
Eccentric Contraction

\medterm{Eccentrocyte Count} The number or proportion of red blood cells whose hemoglobin is pushed toward one side of the cell. This smear finding may indicate oxidative injury or a hemolytic disorder but must be interpreted with other blood results.

\medterm{Ecchymosis} A flat purple, blue, or brown skin mark caused by blood leaking into tissues, commonly after an injury. Easy or unexplained bruising can also result from medicines, low platelets, clotting disorders, fragile vessels, or illness.

\medterm{Eccrine Gland} A small sweat gland found throughout the skin that releases watery sweat onto the surface. It helps cool the body and is especially numerous on the palms, soles, and forehead.

\medterm{Eccrine Glands} Sweat glands distributed through the skin that release watery sweat directly onto the surface. They help control body temperature; disorders can cause excessive sweating, reduced sweating, blockage, or tumors.

\medterm{Eccrine Porocarcinoma} A rare skin cancer arising from cells of sweat-gland ducts. It may appear as a slowly enlarging lump or patch and can recur or spread; biopsy and complete removal are often needed.

\medterm{Eccrine Poroma} A usually harmless growth arising from a sweat-gland duct, often appearing as a flesh-colored, pink, or reddish lump on a hand or foot. It may bleed or be tender and can be removed if needed.

\medterm{ECG} A recording of the heart’s electrical activity made with skin electrodes. It can show heart rate, rhythm, conduction, reduced blood flow, prior heart damage, or other cardiac problems.

\synonyms
Electrocardiogram; EKG

\medterm{ECG/EKG Machine} A device that records the heart’s electrical activity through electrodes placed on the skin. A standard 12-lead recording helps assess rhythm, conduction, heart strain, reduced blood flow, and prior heart injury.

\medterm{Echinacea} A group of plants used in some herbal products for colds or claimed immune support. Evidence of benefit is inconsistent; allergic reactions and interactions with medicines can occur, especially in people with plant allergies.

\medterm{Echinocandins} A group of antifungal medicines that stop fungi building their protective cell wall. They are given mainly by injection for serious Candida infections and selected Aspergillus infections; liver reactions and infusion effects can occur.

\medterm{Echinococcal Cysts} Fluid-filled cysts caused by the larval stage of Echinococcus tapeworms, most often in the liver or lungs. They may remain silent or cause pressure, pain, breathing symptoms, or a severe allergic reaction if they rupture.

\medterm{Echinococcosis} A parasitic infection caused by larval Echinococcus tapeworms. Cystic disease usually forms cysts in the liver or lungs; alveolar disease grows invasively in the liver and may spread, requiring specialist treatment.

\medterm{Echinococcus} A group of tapeworms that can cause disease when people swallow eggs from contaminated soil, food, water, or contact with infected animals. Larvae can form cysts or invasive lesions, especially in the liver or lungs.

\medterm{Echinococcus Granulosus} A tapeworm whose larvae cause cystic echinococcosis. Dogs commonly carry the adult worm, and people become infected by swallowing eggs from contaminated material; cysts most often develop in the liver or lungs.

\medterm{Echinococcus Multilocularis} A tapeworm species whose larvae cause alveolar echinococcosis. The infection mainly affects the liver, where it grows slowly and invasively and may spread to other organs.

\medterm{Echinococcus Vogeli} A tapeworm species that causes polycystic echinococcosis, a disease producing multiple parasitic cysts mainly in the liver and sometimes other abdominal organs.

\medterm{Echinocytes} Red blood cells with many small, evenly spaced projections around their edges. They may reflect kidney disease or an inherited blood disorder, but can also be created by improper storage or preparation of the blood sample.

\synonyms
Burr Cells

\medterm{Echo} A common abbreviation for echocardiography, an ultrasound test that examines the heart’s chambers, valves, pumping action, and blood flow. The full term is preferable in patient instructions.

\medterm{Echocardiogram} An ultrasound picture of the heart that shows its chambers, valves, pumping movement, and blood flow without ionizing radiation. It helps evaluate heart failure, valve disease, congenital defects, clots, and other heart problems.

\medterm{Echocardiography} Ultrasound imaging of the heart. It shows the chambers, valves, pumping action, and blood flow; the probe usually scans through the chest and may be passed into the esophagus for closer views.

\medterm{Echoencephalography} Ultrasound imaging of the brain, historically used mainly in infants to look for enlarged fluid spaces or displacement of brain structures. CT or MRI is now usually preferred for detailed brain imaging.

\medterm{Echogenic Bowel} An ultrasound finding in which a fetus’s bowel appears unusually bright, sometimes as bright as bone. It may be a normal variation or may be associated with infection, chromosome conditions, cystic fibrosis, bleeding, or other fetal problems.

\medterm{Echogenic Intracardiac Focus} A small bright spot, usually in a muscle inside the fetal heart, seen on ultrasound. When isolated and all other findings and screening are reassuring, it is commonly a harmless variation; its significance depends on the full evaluation.

\medterm{Echography} Medical ultrasonography, which uses high-frequency sound waves to create images of tissues, organs, blood flow, or a developing fetus.

\medterm{Echolalia} Repeating words or phrases spoken by someone else, either immediately or after a delay. It can occur in autism, developmental conditions, brain disorders, or some psychiatric illnesses and must be interpreted with other symptoms.

\medterm{Echophenomena} Automatic, non-purposeful imitation or repetition of actions, words, or sounds observed in others without explicit awareness or intent. Subtypes include echolalia (vocal repetition) and echopraxia (imitation of movements), commonly observed in Tourette syndrome, catatonia, and certain neurodevelopmental or neurodegenerative disorders.

\medterm{Echopraxia} Involuntary imitation or repetition of another person’s movements. It may occur in some neurologic, developmental, or psychiatric conditions and is interpreted with other symptoms.

\medterm{Echovirus} A group of enteroviruses spread mainly by the fecal-oral route and sometimes by respiratory secretions. Infection is often mild or silent but can cause rash, respiratory illness, meningitis, encephalitis, heart inflammation, or severe newborn infection.

\medterm{E-Cigarettes And E-Hookahs} Electronic devices that heat a liquid to produce an inhaled aerosol. Many contain nicotine, which can cause dependence; the aerosol may contain irritants and toxic substances and is not harmless water vapor.

\medterm{Eclampsia} A new seizure during pregnancy or after delivery related to preeclampsia or pregnancy-related high blood pressure. It is an emergency requiring protection of breathing, seizure treatment, blood-pressure control, and urgent obstetric care.

\medterm{Eclampsia in Pregnancy} Alternate name; see \textit{Eclampsia}.

\medterm{ECMO} Alternate name; see \textit{Extracorporeal Membrane Oxygenation}.

\medterm{ECOG Performance Status} A scale from 0 to 5 describing how much a person’s illness limits daily activity. It is used to estimate prognosis and decide whether someone may tolerate certain cancer treatments; 0 means fully active and 5 means death.

\medterm{E. Coli} A group of bacteria commonly found in the intestines of people and animals. Most are harmless, but some strains cause diarrhea, urinary infection, or serious illness such as bloodstream infection or kidney injury.

\medterm{E Coli Enteritis} Infection and inflammation of the intestine caused by disease-producing strains of Escherichia coli. It may cause watery or bloody diarrhea, cramps, vomiting, and dehydration; some strains can injure the kidneys.

\medterm{E Coli Infection} Infection caused by Escherichia coli bacteria. Depending on the strain and site, it may cause diarrhea, urinary infection, wound or abdominal infection, newborn meningitis, or bloodstream infection; some strains cause kidney injury.

\medterm{E. Coli O157:H7} A strain of Escherichia coli that produces Shiga toxin and can cause severe cramps and bloody diarrhea. It may injure the kidneys, especially in children; urgent medical advice is needed for bloody diarrhea or reduced urination.

\medterm{Econazole} A topical antifungal that damages the fungal cell membrane and treats ringworm, athlete’s foot, jock itch, and some yeast skin infections. It may cause local irritation, burning, itching, or allergic contact dermatitis.

\medterm{Economy Class Syndrome} Alternate informal name; see \textit{Deep Vein Thrombosis}.

\medterm{Ecstasy} A street name for MDMA, an illegal stimulant and hallucinogenic drug. It can cause overheating, dehydration or dangerously low sodium, fast heartbeat, seizures, confusion, liver injury, and death.

\medterm{Ectasia} Abnormal widening of a blood vessel, duct, tube, or organ. Its importance depends on the structure involved; it may cause pressure, poor flow, bleeding, blockage, or rupture.

\medterm{Ecthyma} A bacterial skin infection that extends deeper than a superficial sore and causes painful, punched-out ulcers covered by thick crusts. It commonly affects the legs and may need wound care and antibiotic treatment.

\medterm{Ecthyma Gangrenosum} A serious skin infection causing rapidly developing dark or black ulcers from tissue death, often due to Pseudomonas bloodstream infection. It occurs mainly in people with severely weakened immunity or very low white-cell counts and requires emergency treatment.

\medterm{Ectoderm} The outer layer of cells in an early embryo. It develops into the skin, hair, nails, nervous system, eyes, ears, and other sense-related structures.

\medterm{Ectodermal Dysplasia} A group of inherited disorders affecting tissues such as hair, teeth, nails, skin, and sweat glands. Features may include sparse hair, abnormal teeth, reduced sweating, overheating, dry skin, and repeated infections.

\medterm{Ectodermal Dysplasias} Inherited disorders affecting structures such as hair, teeth, nails, skin, and sweat glands. Depending on the type, they may cause abnormal teeth or hair, reduced sweating, overheating, dry skin, and repeated infections.

\medterm{Ectomesenchymoma} A rare cancer containing both nerve-related and connective-tissue parts, usually occurring in children. It may appear as a growing mass or cause pain; diagnosis requires tissue examination and specialist cancer treatment.

\medterm{Ectoparasite} A parasite that lives on the outer surface of a host, such as a louse, flea, tick, or mite. It may cause itching, bites, skin disease, anemia, or transmission of infections.

\medterm{Ectopia Cordis} A rare birth defect in which the heart lies partly or completely outside the chest because the chest wall did not form normally. It is usually life-threatening and requires immediate specialist care.

\medterm{Ectopia Lentis} Displacement of the eye’s natural lens from its normal position. Injury or inherited connective-tissue disorders may cause it, leading to blurred or double vision, unequal focusing, or increased eye pressure.

\medterm{Ectopic} Located or occurring in an abnormal position. For example, an ectopic pregnancy implants outside the main cavity of the uterus, most often in a fallopian tube.

\medterm{Ectopic ACTH Syndrome} Excess cortisol production caused by a tumor outside the pituitary making adrenocorticotropic hormone (ACTH). It can cause weight gain, muscle weakness, high blood pressure, diabetes, low potassium, infections, and other features of cortisol excess.

\medterm{Ectopic Beat} An extra or early heartbeat that starts outside the heart’s usual pacemaker. It may feel like a skipped beat, flutter, or thump and can occur in healthy people or with heart disease.

\medterm{Ectopic Cushing Syndrome} Excess cortisol caused by a tumor outside the pituitary producing ACTH or, rarely, a related hormone. It may cause high blood pressure, diabetes, low potassium, muscle weakness, infections, and rapid physical changes.

\medterm{Ectopic Heartbeat} A heartbeat that starts outside the heart’s usual pacemaker. Early beats may feel like a skipped beat, flutter, or sudden thump and can occur in healthy people or with structural heart disease.

\medterm{Ectopic Heartbeats} Extra or early beats arising outside the heart’s usual pacemaker. They may feel like a skipped beat, thump, or flutter and are often harmless, but frequent episodes, fainting, chest pain, or breathlessness need assessment.

\medterm{Ectopic Kidney} A kidney located somewhere other than its usual position in the upper back of the abdomen. It may work normally or be associated with abnormal drainage, urine reflux, stones, or infection.

\medterm{Ectopic Pacemaker} An excitable cardiac focus or group of pacemaker cells located outside the sinoatrial node (in the atria, atrioventricular junction, or ventricles) that prematurely initiates cardiac depolarizations and rhythm, frequently producing ectopic beats or tachycardia.

\synonyms
Ectopic Focus; Latent Pacemaker

\medterm{Ectopic Pregnancy} A pregnancy implanted outside the main cavity of the uterus, most often in a fallopian tube. It may cause missed periods, abdominal pain, or vaginal bleeding; rupture can cause life-threatening internal bleeding and requires emergency care.

\synonyms
Tubal Pregnancy; Extrauterine Gestation

\medterm{Ectopic Thyroid} Thyroid tissue located outside its usual position in the lower front of the neck, most often at the base of the tongue. It may be the only functioning thyroid tissue and can cause swallowing, voice, or breathing problems.

\medterm{Ectopic Thyroid Tissue} Thyroid tissue that developed outside the usual thyroid location because it did not move normally during development. It is often at the tongue base and may be the only working thyroid tissue, sometimes causing swallowing or airway symptoms.

\medterm{Ectopic Tissue} Tissue located in an unusual position, usually because it developed there, failed to move normally during growth, or became implanted elsewhere. Effects depend on the tissue and whether it causes pressure, bleeding, blockage, or abnormal function.

\medterm{Ectopic Ureter} A ureter that does not enter the bladder in its normal position. It may drain into the urethra, vagina, or another abnormal site and can cause urine leakage, infection, blockage, or kidney damage.

\medterm{Ectopion} Outward turning or eversion of a body structure, most often the lower eyelid or the cervix. The effect depends on the site and may include irritation, exposure, bleeding, or abnormal discharge.

\medterm{Ectrodactyly} A birth difference in which one or more central fingers or toes are missing, leaving a deep cleft in the hand or foot. It may occur alone or with other limb or genetic differences and can affect function.

\medterm{Ectromelia} A birth defect in which one or more limbs are absent or severely underdeveloped. In laboratory medicine, the same word names a mouse poxvirus infection; these are separate uses and should not be confused.

\medterm{Ectropion} Outward turning of an eyelid, usually the lower lid, so the eye may not close or drain tears normally. It can cause tearing, dryness, irritation, and corneal injury; aging, weakness, scarring, or paralysis may contribute.

\medterm{Eculizumab} An antibody medicine that blocks part of the immune system called complement. It treats selected blood, kidney, and nerve disorders but greatly increases the risk of meningococcal and other serious infections, so vaccination and monitoring are essential.

\medterm{Eczema} A group of skin conditions that cause itchy, dry, inflamed, red, scaly, cracked, or sometimes oozing skin. The skin barrier is weakened, so the skin loses moisture and becomes easier to irritate. Atopic dermatitis is the most common type. Symptoms may worsen because of irritants, dry weather, heat, sweating, stress, or skin infection. Eczema is not contagious.

\medterm{Eczema Herpeticum} A rapidly spreading herpes simplex infection in skin affected by eczema or another barrier disorder. It causes many similar painful blisters or punched-out sores, often with fever, and needs urgent antiviral treatment because the eyes and other organs may be affected.

\medterm{Eczematous Dermatitis} A pattern of inflammatory skin disease characterized by itching, redness, scaling, dryness, oozing, or vesicles. Causes include atopic, allergic contact, irritant contact, and other forms of dermatitis.

\medterm{ED} Erectile dysfunction, difficulty getting or keeping an erection sufficient for sexual activity. It may reflect blood-vessel, nerve, hormone, medicine, psychological, or relationship factors and can signal broader health problems.

\synonyms
Erectile Dysfunction; Impotence

\medterm{Edema} Swelling caused by extra fluid in body tissues. It may leave a dent when pressed or feel firm. Common causes include heart, kidney, or liver disease, poor vein drainage, blocked lymph flow, and some medicines.

\medterm{Edema, Pulmonary} Alternate index label; see \textit{Pulmonary Edema}.

\medterm{Edema Trunk} Abnormal accumulation of fluid in the soft tissues of the chest or abdomen, producing swelling from heart, kidney, liver, lymphatic, venous, inflammatory, or medication-related causes.

\medterm{Edentulism} Complete or partial absence of natural teeth. It can impair mastication, speech,
nutrition, facial support, and quality of life and may be managed with removable, fixed,
or implant-supported prostheses.

\medterm{Edentulous} Having some or all natural teeth missing. Tooth loss may result from decay, gum disease, injury, extraction, or a developmental condition and can affect chewing, speech, nutrition, jawbone support, and appearance; dentures or implants may restore function.

\medterm{Edentulousness} Complete or partial absence of natural teeth. It may result from dental disease, trauma, developmental abnormalities, or tooth extraction and can impair chewing, speech, and oral function.

\medterm{Edetic Acid} Another name for EDTA, a chemical that binds certain metals. It is used in blood-collection tubes, laboratories, and selected treatments for heavy-metal poisoning.

\medterm{EDTA} A chemical that binds metals. It is used in some blood-collection tubes and, in selected cases, as a supervised treatment for heavy-metal poisoning; unnecessary use can harm the kidneys and lower essential minerals.

\medterm{Edwardsiella} A group of bacteria associated with water and animals. Some species, especially E. tarda, can rarely cause diarrhea, wound infection, bloodstream infection, or other illness in people.

\medterm{Edwardsiella Tarda} A bacterium found in fish, reptiles, and water that can rarely infect people. It may cause diarrhea, wound infection, gallbladder or abdominal infection, or bloodstream infection.

\medterm{Edwards Syndrome} A chromosome condition caused by an extra chromosome 18 in some or all cells. It commonly causes poor growth, developmental impairment, feeding difficulty, heart defects, and other birth abnormalities; medical needs and survival vary widely.

\synonyms
Trisomy 18; Trisomy E

\medterm{EEG} A test that records the brain’s electrical activity through sensors placed on the scalp. It helps evaluate seizures, unexplained episodes, altered awareness, certain sleep disorders, and some brain illnesses.

\synonyms
Electroencephalogram

\medterm{EFAST Examination} A rapid bedside ultrasound used after serious injury to look for internal bleeding, fluid around the heart, and air or blood around the lungs. A normal scan does not exclude every internal injury, so repeat assessment or CT may be needed.

\medterm{Efavirenz} An antiretroviral medicine that blocks HIV reverse transcriptase and is used in combination treatment for HIV infection. It can cause vivid dreams, dizziness, mood or nervous-system effects, rash, and important medicine interactions.

\medterm{Effacement} Thinning and shortening of the cervix during labor, measured as a percentage from 0 to 100. Effacement helps the cervix open and usually occurs along with dilation as contractions progress.

\medterm{Effective Dose} An estimate that combines radiation exposure to different organs into one number. It helps compare the approximate risk of imaging procedures but cannot predict an individual person’s exact risk.

\medterm{Effective Half-Life} The time needed for the amount of a radioactive substance in the body or a tissue to fall by half through both natural decay and removal by the body. It helps estimate how long radiation exposure lasts.

\medterm{Effectiveness} How well a treatment, test, or public-health intervention achieves its intended result in usual real-world practice. It differs from efficacy, which describes performance under ideal, controlled study conditions.

\medterm{Effective Reproduction Number} The average number of people one infected person infects at a particular time, taking account of immunity, behavior, and control measures. Continued spread generally occurs when the number is above 1.

\synonyms
Rt

\medterm{Effector} A cell, molecule, organ, or mechanism that carries out a response after receiving a signal. For example, an immune cell may destroy an infected cell after activation.

\medterm{Efferent} Carrying signals away from a central structure. In the nervous system, efferent pathways carry movement or automatic-control signals from the brain or spinal cord to muscles, glands, or organs.

\medterm{Efferent Arteriole} The small blood vessel that carries blood away from the kidney’s filtering unit. It helps control pressure in the filter and leads to the small vessels around the kidney tubules.

\medterm{Efferent Ductules} A series of 10 to 20 delicate, ciliated convoluted tubules that penetrate the tunica albuginea of the testis to transport spermatozoa from the rete testis into the head (caput) of the epididymis, reabsorbing the majority of testicular luminal fluid.

\synonyms
Ductuli Efferentes

\medterm{Efferent Lymphatic Vessel} A lymphatic vessel that carries lymph out of a lymph node.

\medterm{Efferent Nerve} A nerve pathway carrying signals away from the brain or spinal cord to muscles, glands, or organs. Damage can cause weakness, impaired automatic function, or loss of control in the affected area.

\medterm{Efferent Neuron} A nerve cell that carries movement or automatic-control signals from the brain or spinal cord to muscles, glands, or organs. Damage can cause weakness or impaired control of body functions.

\medterm{Efferent Pathway} A nerve route carrying signals away from the brain or spinal cord to muscles, glands, or organs. It controls movement and automatic functions such as heart rate, digestion, sweating, and bladder activity.

\medterm{Effervescent Tablets} Tablets dissolved in water before use, releasing bubbles as they form a drinkable solution. Some contain substantial sodium or other ingredients, so people with heart, kidney, or blood-pressure problems should check the label.

\medterm{Efficacy} How well a treatment produces its intended result under ideal, carefully controlled study conditions. It differs from effectiveness, which describes how well the treatment works in ordinary clinical practice.

\medterm{Effort Angina} A classic presentation of stable ischemic heart disease characterized by substernal chest discomfort, tightness, or pressure predictably provoked by physical exertion, emotional stress, or cold exposure, and consistently relieved within minutes by rest or sublingual nitroglycerin.

\synonyms
Classic Angina; Exertional Angina

\medterm{Effort Syncope} Transient loss of consciousness occurring abruptly during or immediately following physical exertion, serving as a cardinal red-flag symptom of severe cardiac outflow obstruction (such as critical aortic stenosis or hypertrophic obstructive cardiomyopathy) or exertion-induced ventricular arrhythmia.

\medterm{Effusion} Abnormal buildup of fluid in a body space, such as around a joint, lung, or heart. It may cause swelling, pain, stiffness, or breathing difficulty; the cause may be injury, infection, inflammation, heart failure, or cancer.

\medterm{Effusion-Constriction Syndrome} A rare condition in which fluid around the heart and a stiff heart covering occur together. The fluid and stiffness can prevent the heart from filling normally and may continue to cause pressure problems even after fluid is drained.

\synonyms
Effusive-Constrictive Pericarditis

\medterm{Effusion, Pleural} Alternate name; see \textit{Pleural Effusion}.

\medterm{EGD} An examination of the esophagus, stomach, and duodenum using a flexible camera passed through the mouth. It can investigate pain, swallowing problems, bleeding, or anemia and can obtain samples or treat selected lesions.

\medterm{EGFR} A protein on the cell surface that receives growth signals. Changes that overactivate EGFR can help some cancers grow; testing the tumor may guide use of treatments that block this pathway.

\medterm{EGFR Inhibition} Cancer treatment that blocks growth signals sent through the epidermal growth factor receptor. It may slow selected colorectal, lung, or head-and-neck cancers when the tumor has suitable features; rash, diarrhea, and other effects may occur.

\medterm{Egg} A nutrient-rich food containing protein, fat, vitamins, and minerals, but also a common allergen. Raw or undercooked eggs can transmit Salmonella; thorough cooking and individual allergy or dietary needs guide safe use.

\medterm{Egg Allergy} An immune reaction to egg proteins that may cause hives, swelling, vomiting, stomach pain, wheezing, or anaphylaxis. It is more common in children; reactions vary in severity and require an individualized avoidance and emergency plan.

\synonyms
Egg Hypersensitivity; Hen's Egg Allergy

\medterm{Egg-Cell} A mature female reproductive cell, also called an oocyte, that can unite with sperm during fertilization. It contains half the usual genetic material and can develop into an embryo after fertilization.

\medterm{Eggerthella} A group of bacteria that can live in the human intestine without causing disease. Rarely, especially in people with serious illness, they can cause bloodstream infection, abscesses, or other opportunistic infections.

\medterm{Egg Shell} The hard calcium-containing outer covering of an egg. In medical imaging, “eggshell calcification” describes a thin rim of calcium around a lesion or lymph node.

\medterm{Eggshell Nail} A nail that becomes unusually thin, soft, flexible, and sometimes pale or concave. Possible causes include nutritional deficiency, chronic illness, aging, or repeated trauma; persistent changes should be medically assessed.

\medterm{Egophony} A change in voice sounds heard through a stethoscope in which a spoken “E” sounds like an “A.” It may occur when lung tissue becomes dense or filled with fluid, such as in pneumonia.

\medterm{Ehlers-Danlos Syndrome} A group of inherited connective-tissue disorders that can cause unusually flexible joints, stretchy or fragile skin, easy bruising, and poor wound healing. Some types also weaken blood vessels or internal organs and require specialist care.
\synonyms
EDS; Hypermobility Syndrome

\medterm{EHR} Electronic health record: a digital collection of a person’s health information, such as diagnoses, medicines, allergies, test results, and clinical notes, that can be shared among authorized healthcare professionals.

\medterm{Ehrlichia} A group of bacteria spread mainly by ticks that infect certain white blood cells and cause ehrlichiosis. Infection may produce fever, headache, muscle aches, low blood counts, and serious illness if untreated.

\medterm{Ehrlichia And Anaplasma} Tick-borne bacteria that infect white blood cells. They can cause fever, headache, muscle aches, weakness, low blood counts, and abnormal liver tests; severe disease is more likely in older or immunocompromised people.

\medterm{Ehrlichiosis} A tick-borne infection caused by Ehrlichia bacteria. It commonly causes fever, headache, muscle aches, weakness, low white-cell or platelet counts, and raised liver enzymes; untreated infection can damage the brain, lungs, or other organs.

\medterm{Ehrlichiosis And Anaplasmosis} Tick-borne bacterial infections that may cause fever, headache, muscle aches, weakness, low blood counts, and abnormal liver tests. Prompt medical assessment and appropriate antibiotics reduce the risk of severe complications.

\medterm{Eicosanoid} A signaling substance made from certain fats that helps regulate inflammation, blood-vessel width, platelet activity, pain, and muscle contraction. Prostaglandins, leukotrienes, and thromboxanes are examples.

\medterm{Eicosanoid Degradation} Breakdown and inactivation of eicosanoids, the body’s short-lived fatty-acid signals. This limits how long and how strongly they affect inflammation, blood vessels, platelets, airways, and other muscles.

\medterm{Eicosanoid Metabolism} The body’s production, modification, and breakdown of prostaglandins, leukotrienes, thromboxanes, and related fatty-acid signals that influence inflammation, blood vessels, platelets, pain, and airway function.

\medterm{Eicosanoid Receptor} A protein on or inside a cell that responds to prostaglandins, leukotrienes, thromboxanes, or related lipid signals. Activation can change inflammation, blood-vessel tone, platelet activity, pain, or smooth-muscle contraction.

\medterm{Eicosapentaenoic Acid} The omega-3 fat EPA, found in oily fish and some supplements. It forms part of cell membranes and can be changed into signaling substances involved in inflammation and blood clotting.

\medterm{Eicosatetraenoic Acid} A 20-carbon fatty acid with four double bonds. Arachidonic acid is the best-known form and is used by cells to make prostaglandins, leukotrienes, thromboxanes, and other signaling substances.

\medterm{Eighth Cranial Nerve} The vestibulocochlear nerve, which carries hearing and balance signals from the inner ear to the brain. Damage can cause hearing loss, ringing, dizziness, imbalance, or abnormal eye movements.

\synonyms
Vestibulocochlear Nerve; Auditory Nerve

\medterm{Eikenella} A group of bacteria found in the mouth. Some species can cause deep wound, bone, joint, or bloodstream infection, especially after human bites or injuries involving a person’s teeth.

\medterm{Eikenella Corrodens} A slow-growing bacterium normally found in the mouth that can cause infection after human bites or clenched-fist injuries. It may produce abscesses, bone or joint infection, or rarely bloodstream infection.

\medterm{Eisenmenger Complex} Alternate form; see \textit{Eisenmenger syndrome}.

\medterm{Eisenmenger's Syndrome} Alternate form; see \textit{Eisenmenger syndrome}.

\medterm{Eisenmenger syndrome} A form of advanced pulmonary vascular disease caused by a longstanding congenital heart defect that initially produces excessive blood flow to the lungs. Progressive elevation of pulmonary vascular resistance leads to right-to-left or bidirectional shunting, reduced blood oxygen, and cyanosis.

\medterm{Ejaculation} The release of semen through the penis, usually during sexual climax. It involves coordinated contractions of reproductive ducts, glands, pelvic muscles, and nerves; pain, blood, or major change in ejaculation needs assessment.

\medterm{Ejaculatory Duct} One of two small tubes formed where the vas deferens joins a seminal-vesicle duct. Each passes through the prostate and opens into the urethra, carrying semen during ejaculation.

\medterm{Ejaculatory Ducts} Two small tubes within the prostate that carry semen from the vas deferens and seminal vesicles into the urethra. Blockage or inflammation can contribute to painful ejaculation, infertility, or reduced semen volume.

\medterm{Ejaculatory Dysfunction} Difficulty ejaculating, such as delayed ejaculation, absent ejaculation, painful ejaculation, or ejaculation into the bladder instead of out through the penis. Causes include medicines, nerve or pelvic disorders, diabetes, prostate procedures, and psychological factors.

\medterm{Ejection Click} A sharp heart sound heard just after the heart begins to contract, caused by opening of an abnormal aortic or pulmonary valve. It may occur with a bicuspid valve or valve narrowing and is assessed with other findings.

\medterm{Ejection Fraction} The percentage of blood pumped out of a heart chamber with each beat. It is used to describe the heart’s pumping strength, but a normal value does not rule out heart failure.

\medterm{Ejection Murmur} A midsystolic crescendo-decrescendo heart murmur that begins shortly after the first heart sound (S1), peaks in midsystole, and ends before the second heart sound (S2), produced by turbulent blood flow across an obstructed or narrowed ventricular outflow tract or semilunar valve, as seen in aortic stenosis and pulmonic stenosis.
\synonyms
Systolic Ejection Murmur; Midsystolic Murmur

\medterm{Ekbom Syndrome} An ambiguous historical term that may refer to Willis--Ekbom disease, now usually called restless legs syndrome, or to delusional infestation. The specific modern diagnosis should be used instead.

\medterm{EKG} A recording of the heart’s electrical activity made through electrodes on the skin. It can show heart rate, rhythm, conduction problems, reduced blood flow, prior heart injury, or signs of heart enlargement.

\synonyms
Electrocardiogram; ECG

\medterm{Elastance} A physical and physiological measure of the stiffness and tendency of a hollow organ or vascular structure to recoil toward its original resting dimensions after being distended by pressure, defined mathematically as the change in pressure divided by the change in volume (the inverse of compliance).

\medterm{Elastic Cartilage} Flexible supporting tissue containing many elastic fibers. It helps structures such as the outer ear, epiglottis, and pressure-equalizing ear tube keep their shape while bending.

\medterm{Elastic Fiber} A thin connective-tissue fiber that allows skin, lungs, blood vessels, and other structures to stretch and recoil. Damage or loss of elastic fibers can reduce flexibility or weaken tissues.

\medterm{Elasticity} The ability of a tissue or material to return toward its original shape after stretching or compression. It contributes to the function of skin, lungs, blood vessels, muscles, and ligaments.

\medterm{Elastics} Elastic bands used in medical or dental care, such as orthodontic bands, compression supports, wound devices, or surgical ties. Their purpose and safe use depend on the specific treatment.

\medterm{Elastic Skin} Alternate name; see \textit{Ehlers-Danlos Syndrome}.

\medterm{Elastic Tissue} Connective tissue containing fibers that let structures stretch and return toward their original shape. It helps arteries pulse, lungs expand, skin recoil, and some ligaments support movement.

\medterm{Elastin} A flexible protein in connective tissue that lets skin, lungs, blood vessels, and other structures stretch and recoil. Aging, inflammation, and chronic sun exposure can damage elastin.

\medterm{Elastofibroma} A harmless, slowly growing mass made of fatty and elastic tissue, most often beneath the lower shoulder blade in older adults. It may cause swelling or movement-related discomfort but often needs no treatment.

\medterm{Elastography} An imaging method that estimates how stiff tissue is. It is commonly used to assess liver scarring and may help evaluate breast, thyroid, prostate, or other abnormal areas without tissue removal.

\medterm{Elastomer} A flexible rubber-like material that can stretch and return toward its original shape. It is used in some medical tubing, seals, gloves, devices, and implants; biocompatibility depends on the material.

\medterm{Elastosis} Abnormal buildup, breakdown, or alteration of elastic fibers in tissue. It is commonly associated with long-term sun exposure and can occur in some degenerative, inherited, or lung disorders.

\medterm{Elastosis Perforans Serpiginosa} A rare perforating dermatosis characterized by transepidermal elimination of abnormal elastic fibers, presenting clinically as small, keratotic, umbilicated papules arranged in a serpiginous or annular configuration on the neck, face, or upper extremities, often associated with systemic disorders such as Down syndrome, Ehlers-Danlos syndrome, and pseudoxanthoma elasticum.
\synonyms
EPS; Lutz Disease

\medterm{Elbow} The joint connecting the upper arm and forearm, formed by the humerus, radius, and ulna. It allows bending and straightening and works with nearby joints to rotate the forearm.

\medterm{Elbow Extensor} A muscle that straightens the elbow, mainly the triceps and anconeus. These muscles help push, reach, lift, throw, and control lowering of the arm.

\medterm{Elbow Flexors} Muscles that bend the elbow, mainly the biceps, brachialis, and brachioradialis. They help lift the forearm and turn the palm upward; injury, pain, nerve problems, or tendon damage can weaken movement.

\medterm{Elbow Joint} A joint complex between the humerus, radius, and ulna that allows the arm to bend and straighten and the forearm to rotate. Cartilage, ligaments, muscles, and fluid-filled coverings provide stability and movement.

\medterm{Elbow Pain} Pain in or around the elbow caused by tendon injury, arthritis, bursitis, nerve pressure, fracture, instability, infection, or referred pain from the neck or shoulder. Swelling, deformity, weakness, or injury needs assessment.

\synonyms
Elbow Arthralgia; Cubital Pain

\medterm{Elbow Replacement} Surgery that removes damaged parts of the elbow and replaces them with artificial components to relieve pain and improve movement. It may be used for severe arthritis or injury; infection, loosening, and nerve or bone problems can occur.

\medterm{Elder Abuse} An intentional act or failure to provide care that harms an older person. It includes physical, emotional, sexual, or financial abuse and neglect; warning signs require sensitive assessment and safeguarding according to local law.

\medterm{Elective} Planned in advance rather than performed immediately for an emergency. An elective procedure may still be medically necessary; its timing follows assessment of benefits, risks, alternatives, and the person’s preferences.

\medterm{Elective Abortion} A planned ending of a pregnancy using medicines or a procedure, usually before an immediate medical emergency develops. The method depends on pregnancy duration, health, personal choice, and local clinical and legal requirements.

\medterm{Elective Surgery} Surgery scheduled in advance because it is not needed immediately to prevent serious harm. It may still be important or medically necessary, and timing follows discussion of benefits, risks, alternatives, and preferences.

\medterm{Electrical Alternans} A beat-to-beat change in the size or direction of ECG complexes. It can occur when the heart moves within a large fluid collection around it and may indicate cardiac tamponade, an emergency.

\medterm{Electrical Cardioversion} A controlled electrical shock delivered at a specific point in the heartbeat to restore a normal rhythm, often for atrial fibrillation or flutter. Sedation, clot risk, heart stability, and follow-up treatment must be considered.

\medterm{Electrical Conductance} The ability of a material or tissue to carry electrical current. It depends on how easily charged particles move and is important in nerve, muscle, heart, and medical-device measurements.

\medterm{Electrical Injury} Damage caused by electrical current passing through the body. Severity depends on voltage, current, duration, pathway, and skin resistance; effects may include burns, heart-rhythm problems, muscle breakdown, nerve injury, falls, or breathing failure.

\medterm{Electrical Stimulation} Use of controlled electrical currents to activate nerves or muscles. It may help selected pain, weakness, movement, or rehabilitation problems; benefits and risks depend on the device, body site, and underlying condition.

\medterm{Electrical Synapse} A direct connection between cells through tiny channels that let charged particles and electrical signals pass rapidly from one cell to another. It helps coordinate some brain, heart, and other tissue activity.

\medterm{Electric Axis of the Heart} The net average direction and magnitude of the electrical electromotive forces traveling through the ventricular myocardium during ventricular depolarization, represented on the frontal plane electrocardiogram and normally oriented between -30 degrees and +90 degrees.

\synonyms
Mean QRS Axis

\medterm{Electric Impedance} The opposition a material or tissue offers to alternating electrical current, measured in ohms. Medical devices use impedance changes to assess body composition, monitor signals, or detect contact and function.

\medterm{Electric Organ} A specialized organ in some fish that produces electrical discharges for defense, hunting, navigation, or communication. It is mainly a biological term rather than a human medical structure.

\medterm{Electric Shock} Injury caused when electrical current passes through the body. It can cause burns, muscle damage, falls, abnormal heart rhythms, breathing failure, nerve injury, or internal damage even when skin burns are small.

\medterm{Electrocardiogram} A recording of the heart’s electrical signals from skin electrodes. It can help identify abnormal rhythms, conduction problems, reduced blood flow, heart attack, chamber strain, and some effects of medicines or electrolyte changes.

\synonyms
ECG; EKG

\medterm{Electrocardiographic Gating} An imaging synchronization technique used in cardiac computed tomography, magnetic resonance imaging, and nuclear cardiology that synchronizes image acquisition with specific phases of the patient's cardiac cycle (prospective or retrospective gating) to eliminate cardiac motion artifacts.

\synonyms
ECG Gating

\medterm{Electrocardiography} The process of recording the heart’s electrical activity through skin electrodes, usually as a 12-lead ECG. It assesses rate, rhythm, conduction, chamber strain, reduced blood flow, prior injury, and some medicine or electrolyte effects.

\medterm{Electrocauterization} A procedure using controlled electrical energy to heat tissue, destroy selected tissue, or seal bleeding vessels. It can cause burns, scarring, smoke exposure, or injury to nearby tissue and requires appropriate precautions.

\medterm{Electrocautery} Use of heat produced by electrical current to cut tissue, destroy unwanted tissue, or stop bleeding during a procedure. Possible complications include burns, scarring, smoke inhalation, and damage to nearby structures.

\medterm{Electrocoagulation} A surgical technique using electrical energy to heat tissue, seal blood vessels, or destroy selected tissue. It is used to control bleeding and may cause burns, scarring, or unintended injury.

\synonyms
Surgical Diathermy; Electrocautery

\medterm{Electroconvulsive Therapy} A treatment that causes a brief controlled seizure while the person is under general anesthesia. It can rapidly treat severe depression, catatonia, or mania; temporary confusion and memory problems are common risks.

\medterm{Electrocution} Serious injury or death caused by electrical current passing through the body. It may produce burns, abnormal heart rhythms, muscle or nerve damage, breathing failure, falls, or hidden internal injury and requires emergency evaluation.

\medterm{Electrode} A conductive contact that carries electrical signals or current to or from the body, a device, or a laboratory sample. Electrodes are used in tests such as ECG and EEG and in selected treatments.

\medterm{Electrodessication} A procedure using electrical energy to dry and destroy selected superficial tissue or control bleeding. It may treat some skin lesions, often with scraping, and can cause pain, scarring, infection, or pigment change.

\medterm{Electroencephalogram} A recording of brain electrical activity through scalp electrodes. It helps evaluate seizures, prolonged altered awareness, encephalopathy, selected sleep disorders, and some intensive-care conditions; a normal result does not exclude epilepsy.

\medterm{Electroencephalography} The process of recording brain electrical activity with electrodes on the scalp. It can help identify abnormal patterns associated with seizures, sleep disorders, or diffuse brain dysfunction.

\medterm{Electrogastrogram} A recording of the stomach’s electrical activity, usually through electrodes on the abdominal skin. It measures stomach rhythm and is used mainly in specialist assessment of suspected motility disorders.

\medterm{Electrolysis} Hair removal using electrical energy delivered through a fine probe into individual hair follicles. Repeated treatment may permanently reduce hair growth, but pain, irritation, pigment changes, scarring, or infection can occur.

\medterm{Electrolyte} A mineral that carries an electrical charge when dissolved in body fluids, such as sodium, potassium, calcium, magnesium, chloride, or phosphate. Electrolytes help control nerves, muscles, heart rhythm, fluid balance, and acid-base balance.

\medterm{Electrolyte Balance} Maintenance of appropriate levels of minerals such as sodium, potassium, calcium, magnesium, chloride, and phosphate in body fluids. The kidneys, hormones, intestines, and lungs regulate them; imbalance can affect nerves, muscles, heart rhythm, and hydration.

\medterm{Electrolyte Imbalance} An abnormal level of minerals such as sodium, potassium, calcium, magnesium, chloride, or phosphate in the blood. It may cause weakness, cramps, confusion, seizures, or abnormal heart rhythms; causes include fluid loss, kidney disease, medicines, and hormone disorders.

\medterm{Electrolytes} Electrically charged minerals in body fluids, including sodium, potassium, chloride, bicarbonate, calcium, magnesium, and phosphate. They support nerve conduction, muscle contraction, fluid balance, and acid-base regulation.

\medterm{Electromagnetic Navigation Bronchoscopy} A computer-guided form of bronchoscopy that helps a flexible camera reach small areas near the outer lung. It can guide a biopsy of a lesion not visible with ordinary bronchoscopy, but bleeding, infection, or a collapsed lung can occur.

\medterm{Electromechanical Dissociation} A cardiac arrest in which the monitor shows electrical activity but the heart does not pump enough blood to produce a pulse. It requires immediate CPR and treatment of the underlying cause.

\synonyms
EMD; Pulseless Electrical Activity; PEA

\medterm{Electromyography} A test that records the electrical activity of skeletal muscles, often together with nerve-conduction studies, to evaluate peripheral nerve, nerve-root, neuromuscular-junction, and muscle disorders.

\medterm{Electron-Beam Computed Tomography} A rapid CT technique that uses an electronically steered X-ray beam to image the heart and measure calcium deposits in coronary arteries. It has largely been replaced by newer CT methods.
\synonyms
EBCT

\medterm{Electronic Cigarette} A battery-powered device that heats liquid into an inhaled aerosol, often containing nicotine and other chemicals. It can cause dependence, airway irritation, lung injury, or poisoning from liquid nicotine and is not harmless.

\medterm{Electronic Fetal Monitoring} Continuous or intermittent recording of the fetal heart rate and uterine contractions during pregnancy or labor. It helps identify patterns suggesting that the baby may not be receiving enough oxygen, but results require clinical interpretation.

\medterm{Electron Microscope} A microscope that uses electrons rather than visible light to produce highly magnified images of cells, tissues, microorganisms, or materials.

\medterm{Electron Transport Chain} A series of protein systems inside mitochondria that transfer energy from nutrients to oxygen and use it to make ATP, the cell’s main energy source. Problems in this process can impair organs that need substantial energy, such as the brain and muscles.

\medterm{Electronystagmography} A test that records eye movements with sensors around the eyes, often while the balance system is stimulated. It helps evaluate involuntary eye movements and inner-ear or balance-nerve disorders; video testing is now often preferred.

\medterm{Electrophoresis} A laboratory method that separates charged molecules by how they move through a gel or liquid in an electric field. It can analyze proteins, hemoglobin, DNA, and other substances.

\medterm{Electrophysiologic Testing} Testing that records or deliberately stimulates the heart’s electrical system, usually with catheters, to investigate abnormal rhythms and guide treatment. It is invasive and carries risks such as bleeding, infection, or rhythm disturbance.

\medterm{Electrophysiology} The study of electrical activity in living cells and tissues. Clinical electrophysiology includes tests and treatments involving heart rhythm, nerves, muscles, and the connections between them.

\medterm{Electrophysiology Study} An invasive test in which thin electrode catheters are placed in the heart to map electrical pathways and trigger or measure abnormal rhythms. It can guide ablation or other treatment; bleeding, infection, and rhythm complications are possible.

\medterm{Electroretinography} A test that records the retina’s electrical response to flashes or patterned light. It helps evaluate inherited or acquired retinal disorders when examination and imaging do not fully explain visual loss.

\medterm{Electrosurgery} Use of electrical energy during a procedure to cut tissue, destroy tissue, or control bleeding. Risks include burns, scarring, smoke exposure, unintended tissue injury, and interference with implanted electronic devices.

\medterm{Electrosurgical Unit} A generator that delivers high-frequency electrical energy through an electrode to cut tissue or stop bleeding. Safe use requires attention to current pathways, burns, nearby structures, and possible interference with implanted electronic devices.

\medterm{Element} A pure substance made of atoms with the same number of protons. Some elements, such as iron and iodine, are essential nutrients; others can be toxic, and health effects depend on dose and exposure.

\medterm{Elephant Ear Poisoning} Irritation or poisoning after eating or handling elephant-ear or taro plants, which contain sharp calcium-oxalate crystals. It can cause intense burning and swelling of the mouth and throat and may obstruct breathing.

\medterm{Elephantiasis} Severe, long-lasting swelling and thickening of skin, usually in a limb or the scrotum, caused by blocked lymph drainage. Lymphatic filariasis is one cause; other causes include repeated infection, surgery, cancer, or injury.

\medterm{Elevated Blood Pressure} Blood pressure readings that are repeatedly above the healthy range. Persistent elevation increases the risk of heart disease, stroke, kidney disease, and blood-vessel problems, although a single high reading does not establish a diagnosis.

\medterm{Elevated Hemidiaphragm} An imaging finding in which one side of the diaphragm sits higher than usual. Causes include reduced lung volume, nerve or muscle weakness, abdominal disease, or a normal variation; significance depends on symptoms and other findings.

\medterm{Elevated Liver Enzymes} Higher-than-expected blood levels of enzymes released by stressed or injured liver cells. Causes include fatty liver, medicines, alcohol, viral hepatitis, blocked bile flow, and muscle injury; the pattern and degree guide further evaluation.

\synonyms
Transaminitis; Elevated Aminotransferases

\medterm{Elevation} Upward movement or position of a body part, such as raising the arm, shoulder, jaw, or foot. The term can also describe raising an injured limb to reduce swelling.

\medterm{Elevator} A surgical or dental instrument used to lift, separate, or loosen tissue, bone, or a tooth from nearby structures. The exact design and use depend on the procedure.

\medterm{Eligible} Meeting the requirements for a treatment, clinical trial, screening program, benefit, or healthcare service. Eligibility does not guarantee that an intervention is appropriate or that it will be available.

\medterm{Elimination} Removal of waste products or medicines from the body, mainly through the kidneys, liver, lungs, skin, or intestines. Reduced elimination can make a medicine remain longer in the body and increase its effects or toxicity.

\medterm{Elimination Diets} Short-term eating plans that remove suspected foods and later reintroduce them in a planned way to identify food-related symptoms. They should be supervised when many foods are excluded or nutrition could be compromised.
\synonyms
Exclusion Diet; Diagnostic Elimination Diet

\medterm{Elimination Disorders} A group of childhood-onset disorders involving inappropriate urination or defecation, including enuresis and encopresis. These disorders are diagnosed in a developmentally inappropriate context and are distinguished from incontinence caused by a known medical, neurologic, structural, infectious, or substance-related condition.

\medterm{Elimination Half-Life} The time required for the amount of a medicine in the blood to fall by half. It helps guide dosing intervals and predicts accumulation and clearance; kidney or liver disease, age, and interactions can change it.

\medterm{Elimination Patterns} The timing, frequency, amount, and appearance of urination and bowel movements. Changes can provide clues about hydration, kidney or bowel function, continence, infection, bleeding, or other illness.

\medterm{ELISA Blood Test} A blood test that uses an enzyme-linked reaction to detect or measure an antibody, antigen, hormone, protein, or other substance. Results depend on the test design, timing, sample, and likelihood of disease before testing.

\medterm{Elixir} A sweetened liquid preparation, historically sometimes containing alcohol, used to carry a medicine. The exact ingredients and strength determine its medical effects, interactions, and poisoning risk.

\medterm{Elizabethkingia} A group of environmental bacteria that can rarely cause severe infection, especially in newborns, older adults, or people with weakened immunity. Illness may involve the blood, lungs, brain, or other organs.

\medterm{Elliptocyte Count} The number or proportion of oval or elongated red blood cells seen on a blood smear. Increased elliptocytes may occur in an inherited red-cell disorder or some anemias and must be interpreted with other results.

\medterm{Elliptocytosis} A condition in which many red blood cells are oval or elongated rather than round. Inherited forms may cause red cells to break down early and produce anemia, while some people have no symptoms.

\synonyms
Ovalocytosis

\medterm{Ellis-Van Creveld Syndrome} An inherited skeletal condition causing short limbs, extra fingers or toes, a narrow chest, and abnormal teeth or nails. Congenital heart disease is common and may affect breathing, growth, or physical function.

\medterm{Ellsworth-Howard Test} A specialized test of how the kidneys respond to parathyroid hormone. It measures changes in urine after the hormone is given and may help distinguish pseudohypoparathyroidism from other causes of low parathyroid hormone effects.
\synonyms
PTH Responsiveness Test; Parathyroid Hormone Infusion Test

\medterm{Elongin} A protein complex that helps RNA polymerase II continue copying genetic information into RNA. It contributes to gene regulation and cell function and is mainly a molecular-biology term.

\medterm{Elotuzumab} An antibody medicine that helps the immune system recognize and attack myeloma cells. It is used with other medicines for selected multiple-myeloma patients and can cause infusion reactions, infection, fatigue, or low blood counts.

\medterm{Eltrombopag} An oral medicine that stimulates the thrombopoietin receptor in bone marrow and increases platelet production. It is used for selected low-platelet conditions and can affect the liver, clotting risk, or blood counts.

\medterm{Eltrombopag Olamine} The salt form of eltrombopag, an oral medicine that increases platelet production in selected blood disorders. Liver tests, blood counts, medicine interactions, and clotting risk may require monitoring.

\medterm{Emaciation} Extreme loss of body weight, muscle, and fat caused by severe malnutrition, chronic disease, cancer, infection, or other illness. It can weaken immunity and function and requires assessment of causes and nutritional needs.

\medterm{Embolectomy} Removal of a traveling blood clot or other blockage from a blood vessel with surgery or a catheter procedure to restore blood flow. It may be urgent when a limb or organ is threatened.

\synonyms
Arterial Embolectomy; Surgical Embolectomy

\medterm{Emboli} Particles or material carried through the bloodstream that can lodge in a smaller vessel and block blood flow. Emboli may be clots, fat, air, infected material, or rarely tumor fragments.

\medterm{Embolic Stroke} A stroke caused by a clot or other material traveling to and blocking an artery that supplies the brain. Emergency treatment may restore blood flow and limit permanent brain injury.

\synonyms
Brain Embolism
Cerebral Embolism

\medterm{Embolic Stroke of Undetermined Source} A non-lacunar ischemic stroke detected on neuroimaging without proximal arterial stenosis exceeding 50 percent, a major cardioembolic source (such as atrial fibrillation), or other established specific cause despite comprehensive diagnostic evaluation.

\synonyms
ESUS

\medterm{Embolism} Blockage of a blood vessel by material that traveled through the bloodstream from another site. The material is most often a clot but may be fat, air, infected material, or rarely tumor tissue.

\medterm{Embolism, Pulmonary} Alternate name; see \textit{Pulmonary Embolism}.

\medterm{Embolization} Deliberate blockage of a blood vessel or abnormal vessel connection using a catheter and materials such as coils, particles, or plugs. It can control bleeding, reduce tumor blood flow, or treat selected aneurysms and malformations.

\medterm{Embolization Therapy} An image-guided procedure that blocks selected blood vessels with coils, particles, plugs, or other materials. It can control bleeding, reduce blood supply to a tumor, treat aneurysms, or close abnormal vessel connections.

\medterm{Embolus} Material traveling through the bloodstream that may lodge in a smaller vessel and block it. It may be a blood clot, fat, air bubble, infected material, or rarely tumor tissue.

\medterm{Embryo} A developing human from fertilization through the end of the eighth week after fertilization. From the ninth week until birth, the developing human is called a fetus.

\medterm{Embryogenesis} The process by which an embryo develops after fertilization through cell division, movement, specialization, and formation of organs and body structures.

\medterm{Embryology} The study of development from fertilization through the embryonic and fetal stages, including how organs and body structures form, grow, and become specialized.

\medterm{Embryoma} A rare tumor made of primitive cells resembling early embryonic tissue. The term can refer to different tumor types, so diagnosis depends on the organ involved and microscopic examination.

\medterm{Embryonal Carcinoma} An aggressive germ-cell tumor composed of primitive malignant cells, occurring in the testis, ovary, or other sites and often treated with chemotherapy and surgery.

\medterm{Embryonal Rhabdomyosarcoma} A malignant skeletal-muscle-lineage tumor that commonly affects young children and may arise in the head and neck, urinary tract, or reproductive tract.

\medterm{Embryonal Tumor Of The Central Nervous System} A highly cellular malignant pediatric CNS tumor with primitive or embryonal morphology. Diagnosis depends on integrated histologic and molecular classification, and treatment usually combines surgery, radiation when appropriate, and chemotherapy.

\medterm{Embryonal Tumors} Cancers that arise from cells resembling developing embryonic tissue, often occurring in children and commonly affecting the nervous system.

\medterm{Embryo Neoplasm} A tumor arising from primitive cells resembling developing embryonic tissue, usually in a child. Diagnosis and treatment depend on the tissue of origin, microscopic and molecular features, location, and spread.

\medterm{Embryonic} Relating to an embryo, the developing human or animal during the early period after fertilization. In humans, the embryonic period ends at about eight weeks after fertilization, followed by the fetal period.

\medterm{Embryonic Cell} A cell of an early embryo capable of division and differentiation into specialized tissues; embryonic stem cells have broad developmental potential.

\medterm{Embryonic Induction} A developmental process in which one group of embryonic cells signals nearby cells to develop into a particular tissue or structure.

\medterm{Embryonic Mosaic} A developing embryo containing two or more groups of cells with different genetic makeups, produced by a change after fertilization. Effects depend on the cells involved and the proportion of each cell group.

\medterm{Embryonic Stage} The early developmental period after fertilization during which the basic body plan and major organs begin to form. In human pregnancy it lasts through about eight weeks after fertilization.

\medterm{Embryo Transfer} A fertility procedure in which one or more laboratory-grown embryos are placed into the uterus through a thin catheter. It may be done soon after fertilization or after freezing and thawing; implantation and pregnancy are not guaranteed.

\synonyms
ET; Intrauterine Embryo Transfer

\medterm{Emergence Agitation} Temporary confusion, restlessness, fear, or inconsolable crying while waking from anesthesia. It is more common in children and usually resolves with reassurance and monitoring, but pain, low oxygen, or another complication must be excluded.

\medterm{Emergency Airway Puncture} An emergency procedure that creates an airway through the neck when severe upper-airway blockage prevents breathing and a breathing tube cannot be placed. It is a temporary, high-risk lifesaving measure performed only by trained professionals.

\medterm{Emergency Care} Immediate assessment and treatment for a serious or time-sensitive illness or injury. Clinicians first protect breathing, circulation, brain function, and control of bleeding, pain, or infection, then arrange continued treatment or transfer.

\medterm{Emergency Code} A hospital alert that summons a designated response to an urgent event, such as cardiac arrest, fire, or security threat. The meaning of each code varies by institution and should be confirmed locally.

\medterm{Emergency Contraception} Birth control used after unprotected sex or contraceptive failure to reduce the chance of pregnancy. Options include an oral medicine or copper intrauterine device; effectiveness depends mainly on the method and how soon it is used.

\medterm{Emergency Department} The hospital service that evaluates and treats sudden illness and injury. Staff prioritize people by seriousness, stabilize life-threatening problems, diagnose causes, provide initial treatment, and arrange discharge, admission, or transfer.

\medterm{Emergency Medicine} The medical specialty focused on rapid assessment, stabilization, diagnosis, and initial treatment of sudden illness and injury at any age. It includes resuscitation, urgent procedures, poisoning care, and decisions about admission or transfer.

\medterm{Emergency Procedures} Planned actions for responding safely to urgent events such as cardiac arrest, airway blockage, contamination, equipment failure, or radiation exposure while protecting patients, staff, and the public.

\medterm{Emerging Infectious Disease} An infection that is newly recognized, has recently increased, or is spreading into a new population or region. Causes may include animal-to-human transmission, genetic change, travel, climate, or loss of effective control.

\medterm{Emery-Dreifuss Muscular Dystrophy} A rare inherited muscle disease causing early tightness around the elbows, ankles, and neck followed by slowly progressive weakness. Heart conduction problems, abnormal rhythms, and cardiomyopathy may occur, so regular heart monitoring is essential.

\medterm{Emesis} Forceful emptying of stomach contents through the mouth, commonly called vomiting. It may result from infection, medicines, poisoning, pregnancy, blockage, brain disease, or other illness and can cause dehydration.

\medterm{Emetic} A substance or medicine that causes vomiting. Emetics are generally not used for poisoning because vomiting can cause aspiration, worsen burns or injury, and delay safer treatment.

\medterm{Emetics} Substances that induce vomiting. They are generally avoided for poisoning because vomit can enter the lungs or cause additional injury, and vomiting does not reliably remove a dangerous substance.

\medterm{EMG} A test that records electrical activity in muscles, usually with nerve-conduction testing. It helps investigate weakness, numbness, nerve injury, trapped nerves, muscle disease, and disorders of the nerve-muscle connection.

\synonyms
Electromyography

\medterm{Eminentia} An anatomical term for a raised or prominent part of a body structure. It may provide attachment for a muscle or ligament, form part of a joint, or serve as a landmark during examination or surgery.

\medterm{Emission} The discharge or release of a substance, usually a fluid. In male reproduction, emission is the movement of semen into the urethra before ejaculation.

\medterm{Emmenagogue} A substance traditionally claimed to start or increase menstrual bleeding. Evidence and safety vary, and some products can cause poisoning, miscarriage, or interactions; missed or abnormal periods need medical evaluation rather than self-treatment.

\medterm{Emmetropia} The usual focusing state of the eye in which incoming parallel light focuses on the retina without extra effort. Nearsightedness, farsightedness, or astigmatism causes light to focus away from the retina.

\medterm{Emodin} A plant chemical studied for laxative, anti-inflammatory, and anticancer effects. Evidence and safety are not established for routine treatment, and concentrated products may cause diarrhea, dehydration, or medicine interactions.

\medterm{Emollient} A substance applied to skin to soften and moisturize it, reduce water loss, and strengthen the protective barrier. Emollients are commonly used for dry skin and eczema; greasy products can make floors or clothing slippery.

\medterm{Emotion} A feeling state involving subjective experience, brain and body responses, behavior, and interpretation of events. Emotions are normal, but persistent or extreme changes may occur with mental or physical illness.

\medterm{Emotional Adjustment} The process of adapting thoughts, feelings, and behavior to illness, loss, stress, or a major life change. Difficulty adjusting may cause persistent distress or impaired functioning and may benefit from support or treatment.

\medterm{Emotional Dependency} Reliance on another person for reassurance, self-worth, or emotional stability to a degree that limits independence or causes distress. It is not a diagnosis by itself and may reflect anxiety, low self-esteem, trauma, or relationship difficulties.

\medterm{Emotional Disorders} Mental-health conditions involving persistent or severe changes in mood or emotional responses that cause distress or interfere with daily life, relationships, school, or work. Examples include depressive, anxiety, and trauma-related disorders.

\medterm{Emotional Dysregulation Disorder} Alternate name; see \textit{Borderline Personality Disorder}.

\medterm{Emotional Eating} Eating in response to stress, boredom, sadness, anger, or other feelings rather than mainly physical hunger. Occasional emotional eating is common, but frequent episodes can affect health and may benefit from behavioral or psychological support.

\medterm{Empathy} The ability to understand another person’s feelings and point of view. It supports respectful communication and compassionate care but does not require agreeing with the person or experiencing the same feelings.

\medterm{Emphysema} A type of chronic obstructive pulmonary disease in which air sacs are permanently damaged and enlarged, reducing air exchange. Smoking is a common cause; breathlessness, cough, and reduced exercise capacity may develop gradually.

\medterm{Emphysematous Aortitis} A rare, life-threatening infection of the aortic wall that produces gas within the artery, usually caused by organisms such as \textit{Clostridium septicum}. It may cause severe back or abdominal pain, sepsis, bleeding, or an abnormal connection with the intestine.

\medterm{Emphysematous Cholecystitis} A severe gallbladder infection in which gas-forming bacteria produce gas in the gallbladder wall or cavity. It is more common in people with diabetes and can rapidly cause tissue death, perforation, or sepsis.

\medterm{Emphysematous Hepatitis} A rare, life-threatening gas-forming infection of liver tissue, usually caused by bacteria in people with diabetes or weakened immunity. It can cause fever, abdominal pain, sepsis, and liver failure.

\medterm{Emphysematous Prostatitis} A rare, severe prostate infection in which gas-forming bacteria produce gas within prostate tissue. It may cause fever, pelvic or perineal pain, difficulty urinating, urinary retention, and sepsis.

\medterm{Emphysematous Pyelonephritis} A severe kidney infection in which gas-forming bacteria destroy kidney tissue. It occurs mainly in people with diabetes and can cause fever, flank pain, shock, and sepsis; urgent drainage, antibiotics, and sometimes surgery are needed.

\medterm{Empirical} Based on observation, measurement, or experience rather than theory alone. Empirical treatment begins before the exact cause is confirmed and is revised when test results or the clinical course provide more information.

\medterm{Empirical Treatment} Treatment started on the basis of symptoms and likely causes before a diagnosis is confirmed. It should be reassessed when test results or the response to treatment becomes available.

\medterm{Empiric Antibiotic Therapy} Starting an antibiotic before the exact germ and its medicine sensitivity are known. The first choice is based on the symptoms, likely source of infection, local resistance patterns, and the person’s health; it may be changed when test results return.

\medterm{Empiric Risk} An estimate of risk based on observed data from a population or clinical setting. Its usefulness depends on how similar the studied population, exposure, and outcome are to the person or situation being assessed.

\medterm{Employee} A person who works for an employer under an agreement. In occupational health, assessment may consider workplace exposures, safety, physical or mental function, and ability to perform duties safely.

\medterm{Empty Calories} A nontechnical term for foods or drinks that provide energy but little protein, fiber, vitamins, or minerals, such as many sugary drinks and highly processed snacks. Frequent intake can displace more nutritious foods.

\medterm{Empty Can Test} A clinical examination manoeuvre used to evaluate the integrity of the supraspinatus muscle and tendon. The patient's arms are abducted ninety degrees in the scapular plane with full internal rotation (thumbs pointing downward), and the examiner applies downward pressure against active resistance; weakness or pain indicates supraspinatus tendinopathy or tear.

\synonyms
Jobe Test; Supraspinatus Test

\medterm{Empty Nose Syndrome} A poorly understood, controversial condition that can follow extensive surgery on the nasal turbinates, in which the nose feels blocked even though the airway is open. People may describe a dry, suffocating or breathless feeling, inability to sense airflow, crusting, headache, fatigue, and poor concentration, with little relief from decongestants. Findings of obstruction may be absent, so diagnosis is clinical; management is supportive and includes moisturizing, saline care, and consultation with clinicians experienced with the condition.

\medterm{Empty Sella Syndrome} A condition in which fluid presses the pituitary gland flat within its bony space. It is often found incidentally and causes no symptoms, but some people have headaches, vision problems, or reduced pituitary-hormone production.

\medterm{Empyema} A collection of pus in a normally sterile body space, most often infected fluid around a lung. It can cause fever, chest pain, cough, breathlessness, and sepsis and usually requires antibiotics plus drainage.

\medterm{Empyema Necessitans} Extension of infected pus from the space around a lung through the chest wall, forming a painful swelling. It is a serious complication requiring imaging, antibiotics, and drainage.

\medterm{Empyema Necessitatis} The direct extension and burrowing of an intrapleural empyema through the parietal pleura and chest wall tissues into the subcutaneous tissue and skin, creating an extrapleural abscess, classically caused by Mycobacterium tuberculosis, Actinomyces israelii, or neglected bacterial empyema.
\synonyms
Empyema Necessitatis; Pleurocutaneous Fistula

\medterm{EMR} Electronic medical record: a digital record of health information maintained within a particular healthcare organization or practice. It may include diagnoses, notes, medicines, allergies, test results, and treatment plans.

\medterm{Emtricitabine} An antiretroviral medicine that blocks HIV reverse transcriptase and is used with other medicines to treat or prevent HIV infection. Kidney function and hepatitis B status are important because stopping it can worsen hepatitis B.

\medterm{Emulsion} A mixture of two liquids that normally do not blend, such as oil and water, with one dispersed in the other. Emulsions are used in some medicines, nutrition products, cosmetics, and laboratory preparations.

\medterm{Enalapril} A medicine that blocks formation of angiotensin II, relaxing blood vessels and reducing strain on the heart. It treats high blood pressure and heart failure but can cause cough, high potassium, kidney problems, or rarely facial swelling.

\medterm{Enalaprilat} The active injectable form related to enalapril that blocks angiotensin II formation and lowers blood pressure. It is used in selected settings and can cause low blood pressure, high potassium, kidney problems, or facial swelling.

\medterm{Enamel} The hard, mineral-rich outer covering of a tooth crown. It protects the sensitive inner tooth but cannot repair itself after major damage; acids and bacteria can cause erosion and decay.

\medterm{Enamel Hypoplasia} Incomplete formation of tooth enamel while a tooth is developing, leaving it thin, grooved, pitted, or discolored. Affected teeth are more vulnerable to sensitivity, wear, and decay.

\medterm{Enamel Organ} The developing tissue of a tooth that produces enamel and helps shape the crown. It is present during tooth formation and later contributes to the protective covering around the developing tooth.

\medterm{Enamel Pearl} A small lump of enamel on a tooth root, often near where the roots divide. It can make cleaning difficult and interfere with gum attachment, increasing the risk of localized gum disease.

\medterm{Enanthem} A rash, spot, ulcer, or other lesion on a mucous membrane, especially inside the mouth. It may occur with infection, inflammation, an allergic or medicine reaction, or a systemic illness.

\medterm{Enanthema} A rash or lesion on a moist membrane lining the mouth, throat, genitals, or another body surface. Infection, inflammation, medicines, or systemic disease may cause it.

\medterm{Enarthrosis} A ball-and-socket joint that allows movement in many directions, including bending, straightening, rotation, and movement away from or toward the body. The hip and shoulder are examples.

\medterm{Encapsulated} Enclosed by a capsule or clearly defined layer of tissue. This may describe a cyst, growth, infection, or organ; being encapsulated does not by itself prove that a lesion is harmless or cancerous.

\medterm{Encapsulated Peritoneal Sclerosis} A rare serious complication of long-term peritoneal dialysis in which scar tissue forms around the intestines. It can block the bowel and cause repeated pain, vomiting, and poor nutrition.

\synonyms
Sclerosing Encapsulating Peritonitis; EPS

\medterm{Encapsulating Peritoneal Sclerosis} A rare condition in which thick scar tissue surrounds the intestines and can cause repeated blockage, pain, vomiting, and poor nutrition. It is most associated with long-term peritoneal dialysis but can have other causes.

\medterm{Encephalitis} Inflammation of brain tissue, often caused by an infection or an abnormal immune reaction. It may cause fever, headache, confusion, seizures, behavior change, weakness, or movement problems and requires urgent medical care.

\medterm{Encephalitis Lethargica} A rare inflammatory brain disorder described mainly during the early twentieth century, causing sleepiness, abnormal movements, eye-movement problems, and other neurologic symptoms. Its cause remains uncertain.

\medterm{Encephalitis Virus} A virus capable of infecting the brain and causing encephalitis. It may spread through mosquitoes, ticks, animals, respiratory secretions, food, or other routes depending on the virus.

\medterm{Encephalitozoon} A group of tiny spore-forming parasites that can infect the intestine, eyes, kidneys, or several organs. Serious disease occurs mainly in people with weakened immunity.

\medterm{Encephalocele} A birth defect in which the coverings of the brain and sometimes brain tissue protrude through an opening in the skull. It may cause a visible head swelling and can be associated with seizures, fluid buildup, developmental problems, or other brain abnormalities.

\medterm{Encephalomalacia} Softening or loss of brain tissue after injury, stroke, bleeding, infection, or reduced blood flow. The resulting permanent damage may cause seizures, weakness, movement problems, memory changes, or other neurologic effects.

\medterm{Encephalomyelitis} Inflammation affecting both the brain and spinal cord, caused by infection or an abnormal immune response. It may cause confusion, seizures, weakness, numbness, balance problems, or bladder changes and requires urgent neurologic assessment.

\medterm{Encephalopathy} A disorder that broadly affects brain function, causing confusion, drowsiness, personality change, difficulty thinking, or coma. Causes include infection, toxins, organ failure, medicines, metabolic problems, seizures, and reduced oxygen.

\medterm{Encephalotrigeminal Angiomatosis} A birth condition involving abnormal blood vessels in the skin of the face and the coverings of the brain. It may cause a port-wine stain, seizures, weakness, learning problems, or glaucoma, but symptoms vary.

\synonyms
Sturge-Weber Syndrome; SWS

\medterm{Enchondroma} A usually harmless cartilage growth inside a bone, often in the hands or long bones and found by chance. Pain, enlargement, fracture, or unusual imaging features may require evaluation for a more serious tumor.

\medterm{Enchondromatosis} A disorder in which multiple cartilage growths develop inside bones. It can cause uneven growth, limb deformity, pain, or fractures; a small number of lesions may become cancerous and require monitoring.

\medterm{Encoding} The process by which sensory or other information is changed into a form the brain can store in memory. Attention, sleep, stress, illness, and brain injury can affect encoding.

\medterm{Encopresis} Repeated passage of stool in inappropriate places by a child at an age or developmental level when bowel control is expected. It is often related to long-term constipation and retained stool but may have other medical, developmental, or emotional contributors.

\medterm{Endangiitis} Inflammation of the inner lining of a blood vessel. Swelling or scarring may narrow the vessel and reduce blood flow; the effects depend on the cause and the vessels involved.

\medterm{Endarterectomy} Surgery that removes fatty plaque and diseased inner lining from an artery to improve blood flow, most often a carotid artery in the neck. Possible complications include bleeding, stroke, heart attack, or nerve injury.

\medterm{Endarteritis} Inflammation of the inner lining of an artery that can narrow the vessel and reduce blood flow. Causes include infection, autoimmune disease, injury, and other vascular disorders.

\medterm{Endemic} Regular presence of a disease or infectious agent in a particular population or region at an expected baseline level. Endemic does not mean that every person is infected or that cases cannot increase.

\medterm{Endemic Diseases} Diseases that remain present at a predictable level in a particular population or region. Their frequency is shaped by local environment, immunity, behavior, exposures, healthcare access, and prevention measures.

\medterm{Endemic Fungi} Fungi naturally found in particular geographic regions whose spores can cause infection when inhaled. Disease may affect the lungs and sometimes spread to other organs, especially in people with weakened immunity.

\medterm{Endemic Goiter} Thyroid enlargement occurring commonly in a population, traditionally because of insufficient iodine intake. It may occur with normal, low, or high thyroid function depending on the cause.

\medterm{Endergonic} Describing a chemical reaction that requires an input of energy. In cells, energy from food or another reaction drives processes such as building proteins, DNA, or other complex molecules.

\medterm{Ending Pregnancy With Medicines} Ending a pregnancy with prescribed medicines that cause the uterus to empty. The appropriate method depends on pregnancy duration, health, symptoms, and local clinical requirements; follow-up confirms completion and checks for complications.

\medterm{Endobronchial Foreign Body} An object lodged in a bronchus, partly or completely blocking an airway. It may cause sudden coughing, one-sided wheezing, collapse of part of the lung, or pneumonia and often requires urgent bronchoscopy for removal.

\medterm{Endobronchial Tumor} A growth inside a bronchus that can narrow the airway and cause cough, wheezing, coughing blood, lung collapse, or repeated pneumonia. It may be benign or cancerous; diagnosis usually requires imaging and tissue sampling.

\medterm{Endobronchial Ultrasound} A procedure combining bronchoscopy with ultrasound to examine structures beside the airways and guide needle sampling of lymph nodes or masses. It is commonly used to diagnose and stage lung cancer and other chest diseases.

\medterm{Endobronchial Valve} A one-way valve placed through a bronchoscope into a selected airway to reduce air trapping in severe emphysema. It lets air and mucus leave but limits new air entering; pneumothorax and infection are possible complications.

\medterm{Endocannabinoid} A signaling molecule made naturally by the body that activates cannabinoid receptors. Endocannabinoids help regulate pain, appetite, mood, memory, movement, and immune activity.

\medterm{Endocardial Cushion Defect} A birth defect in which the central wall of the heart and the valves between its chambers do not form normally. It can allow abnormal blood flow and cause heart failure, high pressure in lung vessels, poor growth, or breathing difficulty.

\medterm{Endocardial Fibroelastosis} Thickening of the inner lining of the heart from excess fibrous and elastic tissue, usually in infants or children. It can stiffen the heart or weaken its pumping and may lead to heart failure.

\medterm{Endocarditis} Infection or inflammation of the heart’s inner lining, usually involving a valve. It can cause fever, fatigue, a new or changed murmur, stroke or other emboli, valve destruction, or heart failure.

\synonyms
Infective Endocarditis; Valvular Endocarditis

\medterm{Endocarditis Prophylaxis} Preventive antibiotics given before selected dental procedures to people at highest risk of serious infective endocarditis complications. They are not needed for everyone; eligibility depends on the heart condition and procedure.

\medterm{Endocarditis Surgical Indications} Reasons surgery may be needed for endocarditis include heart failure from valve damage, uncontrolled infection or abscess, infection of a prosthetic valve, fungal infection, or a high risk of emboli. Decisions require a specialist heart team.

\medterm{Endocardium} The thin inner lining of the heart chambers and valves. It provides a smooth surface for blood flow and continues into the inner lining of blood vessels; inflammation or infection can damage it.

\medterm{Endocervical} Relating to the inner canal or lining of the cervix, the lower part of the uterus that connects the uterine cavity with the vagina.

\medterm{Endocervical Culture} A laboratory test that attempts to grow bacteria or fungi from a sample collected from the inner cervix. It may help investigate cervicitis or selected infections, although molecular tests are preferred for many sexually transmitted infections.

\medterm{Endocervical Gram Stain} A microscope test that stains cells and organisms from the inner cervix. It may show inflammation or some bacteria but cannot by itself reliably diagnose every sexually transmitted infection.

\medterm{Endocrine} Relating to hormones released into the bloodstream by glands or specialized cells. Endocrine organs include the thyroid, pituitary, adrenal glands, pancreas, and reproductive glands.

\medterm{Endocrine Disease} A disorder of a hormone-producing gland or hormone action. It may affect growth, metabolism, reproduction, stress responses, blood pressure, or fluid and electrolyte balance.

\medterm{Endocrine Disorders} Diseases in which hormone-producing glands make too much or too little hormone, or the body does not respond to hormones properly. They may affect growth, metabolism, reproduction, mood, blood pressure, and fluid balance. Examples include thyrotoxicosis, hypothyroidism, hypogonadism, Cushing syndrome, and Addison's disease.

\medterm{Endocrine Disrupting Chemicals} Chemicals that can interfere with hormone production, transport, breakdown, or action. Some exposures are associated with reproductive, developmental, thyroid, metabolic, or cancer outcomes, but risk depends on the substance, dose, and timing.

\medterm{Endocrine Disruptor} A chemical that interferes with hormone systems. Possible effects include altered development, reproduction, thyroid function, or metabolism; health risk depends on the chemical, exposure level, duration, and timing.

\medterm{Endocrine Finding} A symptom, examination result, imaging result, or laboratory result that may suggest abnormal hormone production, hormone action, or an endocrine-gland disorder.

\medterm{Endocrine Function} The production, release, transport, and action of hormones that regulate metabolism, growth, reproduction, stress responses, blood pressure, and fluid or electrolyte balance.

\medterm{Endocrine Gland} An organ that releases hormones directly into the bloodstream. Its chemical messages regulate growth, metabolism, reproduction, stress responses, blood pressure, and other body functions.

\medterm{Endocrine Glands} Ductless organs that release hormones into the bloodstream to regulate distant tissues. They include the pituitary, thyroid, adrenal glands, pancreas, and reproductive glands.

\medterm{Endocrine Hypertension} High blood pressure caused by excess or abnormal hormone activity. Causes include primary aldosteronism, excess cortisol, or excess adrenaline-like hormones from an adrenal or related tumor.

\medterm{Endocrine Neoplasm} A benign or cancerous growth arising in a hormone-producing gland or cell. It may release too much hormone or cause symptoms by pressing on nearby structures; evaluation depends on its location and behavior.

\medterm{Endocrine Pancreas} The hormone-producing part of the pancreas, made mainly of pancreatic islets. Its cells release insulin, glucagon, somatostatin, and other hormones that help regulate blood glucose and digestion.

\medterm{Endocrine System} The network of glands and hormone-producing cells and tissues that release chemical messengers into the blood. It includes structures such as the hypothalamus, pituitary, thyroid, parathyroids, adrenal glands, pancreatic islets, and gonads; its hormones regulate growth, metabolism, reproduction, stress responses, sleep, blood pressure, fluid balance, and other aspects of homeostasis.

\medterm{Endocrine Testing} Blood, urine, saliva, stimulation, suppression, and sometimes imaging tests used to assess hormone production and action. Results depend on timing, medicines, age, sex, illness, and the test method.

\medterm{Endocrine Therapy} Treatment that lowers, blocks, or changes hormone signals to control hormone-sensitive cancer. It may be used for breast, prostate, or other cancers; effects depend on the medicine and can include hot flashes, sexual changes, or bone loss.

\medterm{Endocrine Tissues} Cells or groups of cells that produce and release hormones into the bloodstream or nearby tissues. They are found in glands and in organs such as the pancreas, ovaries, testes, kidneys, and intestine.

\medterm{Endocrinologist} A physician who diagnoses and treats disorders of hormones, endocrine glands, and metabolism, including diabetes, thyroid disease, adrenal disease, osteoporosis, and reproductive hormone problems.

\medterm{Endocrinologists} Physicians specializing in hormone, endocrine-gland, and metabolic disorders, including diabetes, thyroid disease, adrenal disease, osteoporosis, and reproductive hormone problems.

\medterm{Endocrinology} The medical specialty concerned with hormones, endocrine glands, and metabolism. It diagnoses and treats conditions affecting diabetes, thyroid function, growth, reproduction, bones, adrenal glands, and fluid or electrolyte balance.

\medterm{Endocrinopathy} A disorder of an endocrine gland or hormone system that causes too much, too little, or abnormal hormone action. Effects depend on the hormone and may involve growth, metabolism, reproduction, mood, or fluid balance.

\medterm{Endocytosis} The process by which a cell folds its outer membrane around material and brings it inside in a small sac. Cells use it to absorb nutrients, remove particles, and control surface signals.

\medterm{Endoderm} The innermost of the three early embryonic cell layers. It develops into the linings of the digestive and respiratory systems and contributes to organs such as the liver, pancreas, thyroid, and bladder.

\medterm{Endodontics} The dental specialty concerned with the soft tissue inside teeth, root canals, and tissues around tooth roots. Endodontists diagnose and treat tooth infection, inflammation, injury, and related pain.

\medterm{Endodontist} A dentist with advanced training in diagnosing and treating disease or injury of the tooth’s inner tissue, root canals, and tissues around tooth roots. They commonly perform complex root-canal treatment.

\medterm{End-Of-Life Care} Care for a person nearing death that focuses on comfort, symptom relief, dignity, communication, and personal goals. It also supports family and caregivers and may include hospice or palliative services.

\medterm{Endogenous} Originating within the body or produced by its own cells or tissues, rather than introduced from outside. Endogenous hormones, chemicals, infections, or substances may be contrasted with external sources.

\medterm{Endogenous Opiates} Opioid-like substances naturally made by the body, including endorphins, enkephalins, and dynorphins. They act on opioid receptors and influence pain, stress, mood, reward, movement, and other nervous-system functions.

\medterm{Endogenous Opioids} Opioid-like peptides made naturally by the body, including endorphins, enkephalins, and dynorphins. They activate opioid receptors to influence pain, stress, mood, reward, movement, and hormone regulation.

\synonyms
Endorphins and Enkephalins; Endogenous Opioid Peptides

\medterm{Endogenous Retrovirus} A piece of viral DNA inherited in the human genome from an ancient virus that infected a reproductive cell. Most such sequences are inactive, although some can influence gene control or normal development.

\medterm{Endoleak} Continued blood flow into the weakened part of an artery after endovascular aneurysm repair. Some types indicate an inadequate seal or device problem and may need urgent repair; others are monitored according to sac growth and risk.

\medterm{Endolimax} A group of amoebas that can live in the intestine. They are usually harmless colonizers rather than causes of disease, and a stool finding should be interpreted with symptoms and other test results.

\medterm{Endolimax Nana} An amoeba that can live harmlessly in the intestine and appear in stool tests. It usually does not cause illness; its presence may indicate exposure to contaminated food or water and should not alone explain digestive symptoms.

\medterm{Endolymph} Fluid inside the membranous part of the inner ear that helps sensory organs detect sound and head movement. Abnormal pressure or composition can contribute to dizziness, hearing changes, or ringing.

\medterm{Endolymphatic Hydrops} Pathological distension and increased volume of the endolymphatic fluid compartment of the membranous inner ear (scala media and vestibular apparatus) resulting from impaired fluid resorption in the endolymphatic sac, constituting the underlying pathophysiologic hallmark of Ménière disease.
\synonyms
Meniere Hydrops; Inner Ear Hydrops

\medterm{Endometrial} Relating to the endometrium, the inner lining of the uterus. This lining changes during the menstrual cycle and can support implantation of a pregnancy; abnormal growth may cause bleeding or affect fertility.

\medterm{Endometrial Ablation} A procedure that destroys much of the uterine lining to reduce heavy menstrual bleeding. It is not contraception, pregnancy afterward can be dangerous, and it is generally intended for people who do not plan future pregnancy.

\medterm{Endometrial Adenocarcinoma} A gland-forming cancer of the uterine lining. It often causes abnormal bleeding, especially after menopause; treatment depends on the cancer’s type, grade, stage, molecular features, and overall health.

\medterm{Endometrial Biopsy} Removal of a small sample of the uterine lining for laboratory examination, often to investigate abnormal bleeding or suspected overgrowth or cancer. Cramping and light bleeding are common; occasionally the sample is insufficient for diagnosis.

\medterm{Endometrial Cancer} Cancer arising from the inner lining of the uterus. Abnormal bleeding—especially after menopause—is a common warning sign; evaluation usually includes examination, imaging, and a sample of the uterine lining.

\synonyms
Uterine Endometrial Neoplasm; Endometrial Carcinoma

\medterm{Endometrial Carcinoma} Cancer arising from the uterine lining, often causing abnormal bleeding, particularly after menopause. Diagnosis requires examination of uterine-lining tissue, and treatment depends on type and spread.

\medterm{Endometrial Disorder} Any condition affecting the uterine lining, including inflammation, polyps, abnormal overgrowth, endometriosis-related changes, or cancer. It may cause heavy or irregular bleeding, pelvic pain, infertility, or no symptoms.

\medterm{Endometrial Hyperplasia} Abnormal overgrowth and thickening of the uterine lining, often from too much estrogen without enough progesterone. It may cause heavy, irregular, or postmenopausal bleeding; atypical changes can precede cancer and need evaluation.

\medterm{Endometrial Implants} Alternate name; see \textit{Endometriosis}.

\medterm{Endometrial Intraepithelial Neoplasia} A precancerous overgrowth of abnormal glands in the uterine lining. It can exist with or progress to endometrial cancer, so treatment may involve removal of the uterus or hormone treatment when preserving fertility is important.

\synonyms
EIN; Atypical Endometrial Hyperplasia

\medterm{Endometrial Neoplasm} An abnormal growth of the uterine lining that may be benign, precancerous, or cancerous. Symptoms, tissue examination, imaging, and—if cancer is present—stage determine treatment.

\medterm{Endometrial Polyp} A localized growth of the uterine lining that projects into the uterine cavity. It may cause irregular or heavy bleeding, infertility, or no symptoms; some polyps contain precancerous or cancerous cells.

\medterm{Endometrial Polyps} Localized growths of the uterine lining projecting into the uterine cavity. They may cause irregular or heavy bleeding, infertility, or no symptoms, and some require removal for tissue examination.

\medterm{Endometrial Receptivity Analysis} A test that examines activity of genes in a sample of the uterine lining to estimate when it may best support embryo implantation. Its ability to improve pregnancy rates is uncertain, so results need specialist interpretation.

\synonyms
ERA; Endometrial Receptivity Assay

\medterm{Endometrial Stromal Nodule} A rare, usually harmless, well-defined growth of connective-tissue cells in the uterus. It differs from cancer by not invading nearby tissue or blood vessels, but microscopic examination is needed for reliable classification.

\medterm{Endometrial Stromal Sarcoma} A rare cancer arising from connective-tissue cells of the uterine lining. Some forms grow slowly and others aggressively; abnormal bleeding or pelvic symptoms may occur, and treatment depends on grade, spread, hormone receptors, and molecular findings.

\synonyms
ESS

\medterm{Endometrioid Carcinoma Of The Vulva} A rare gland-forming cancer of the vulvar skin or nearby tissue that resembles cancer of the uterine lining. Diagnosis requires tissue examination and confirmation that it did not spread from the uterus or ovary.

\medterm{Endometrioid Ovarian Carcinoma} A type of ovarian cancer whose cells resemble glands of the uterine lining. It may be associated with endometriosis; symptoms, treatment, and outlook depend on stage, grade, and tumor features.

\medterm{Endometrioma} An ovarian cyst caused by endometriosis in which old blood collects. It may cause pelvic pain, painful periods, pain during sex, or infertility, although some people have no symptoms.

\medterm{Endometriosis} A chronic inflammatory condition in which tissue resembling the lining of the uterus grows outside the uterus, most often in the pelvis, ovaries, fallopian tubes, or tissue around the pelvic organs. It can also affect the bowel or bladder. This tissue responds to hormones and can cause repeated inflammation, bleeding, scar tissue, adhesions, or ovarian cysts. Symptoms may include severe period pain, ongoing pelvic pain, heavy or irregular bleeding, pain during sex or bowel movements, painful urination, bloating, nausea, tiredness, or infertility. Some people have no symptoms.

\synonyms
Endometrioma; Pelvic Endometriosis

\medterm{Endometritis} Inflammation or infection of the uterine lining, often after childbirth, miscarriage, abortion, or a procedure entering the uterus. It may cause fever, pelvic pain, tenderness, abnormal discharge, or bleeding and needs treatment.

\medterm{Endometrium} The inner lining of the uterus. It thickens each month under hormone control to prepare for pregnancy and sheds during menstruation; abnormal thickening, inflammation, polyps, or cancer can cause bleeding or fertility problems.

\medterm{Endomyocardial Biopsy} A procedure that removes tiny samples from the heart muscle through a catheter in a vein. The samples can help check for rejection after a heart transplant or investigate some heart-muscle diseases.

\synonyms
EMB

\medterm{Endomyocardial Fibrosis} Scarring and thickening of the inner heart lining and nearby muscle that makes the heart stiff and reduces filling. It can cause breathlessness, swelling, abnormal rhythms, or heart failure.

\medterm{Endomysium} A delicate layer of reticular connective tissue that individually invests each single skeletal muscle fiber, containing an intricate network of capillaries and nerve terminals essential for metabolic exchange and muscle contraction.

\medterm{Endoneurium} A delicate layer of loose connective tissue and collagen fibrils that directly surrounds each individual myelinated or unmyelinated axon and its accompanying Schwann cell sheath within a peripheral nerve fascicle.

\medterm{Endoparasite} A parasite that lives inside its host, such as a worm in the intestine or a microscopic parasite in blood or tissue. It may cause no symptoms or lead to inflammation, anemia, organ damage, or infection.

\medterm{Endopeptidase} An enzyme that cuts proteins or peptides at internal points rather than removing amino acids only from an end. Digestive enzymes such as pepsin and trypsin are examples.

\medterm{Endophthalmitis} A severe infection or inflammation inside the eye, usually after eye surgery, an injury, or spread through the bloodstream. It can cause severe pain, redness, swelling, and rapid vision loss and requires emergency eye treatment.

\medterm{Endoplasmic Reticulum} A network of tiny membrane-lined spaces inside cells. It helps make and transport proteins and fats, stores calcium, and processes substances needed for cell structure and function.

\medterm{Endoprosthesis} An artificial device placed inside the body to replace, support, or restore a body structure. Examples include joint replacements, vascular grafts or stents, and some devices used to keep an airway open.

\medterm{Endoribonuclease} An enzyme that cuts RNA at internal chemical bonds. It helps process RNA into usable forms or break down RNA that is damaged, unnecessary, or produced by viruses.

\medterm{Endorphin} A natural chemical made by the body that acts on opioid receptors in the brain and other tissues. It can reduce the feeling of pain and influence stress, mood, and other body responses.

\medterm{Endorphins} Natural opioid-like chemicals made by the body that reduce pain and influence mood. They are released during stress, exercise, injury, and other situations and act through opioid receptors in the nervous system.

\synonyms
Endogenous Opioid Peptides; Beta-Endorphin

\medterm{Endoscope} A flexible or rigid instrument with a light and camera used to view an internal body passage or cavity. It can sometimes collect tissue, stop bleeding, remove growths, or perform other treatments.

\medterm{Endoscopic Endonasal Surgery} A minimally invasive operation performed through the nostrils and nasal passages using an endoscope, commonly to reach selected skull-base, pituitary, sinus, or nearby lesions. Risks include bleeding, infection, cerebrospinal-fluid leak, and injury to nearby structures.

\medterm{Endoscopic Mucosal Resection} An endoscopic procedure that removes a superficial growth from the digestive tract after lifting it away from deeper tissue. The removed tissue is examined to determine whether it was completely removed and whether cancer is present.

\medterm{Endoscopic Retrograde Cholangiopancreatography} A procedure combining an endoscope and X-ray imaging to examine and treat the bile and pancreatic ducts. It can remove stones, open narrowed ducts, place stents, or obtain samples; pancreatitis, bleeding, and infection are possible risks.

\medterm{Endoscopic Submucosal Dissection} An endoscopic procedure that removes a larger or flatter superficial growth from the digestive tract in one piece by dissecting beneath it. It allows detailed tissue examination but carries risks of bleeding and perforation.

\medterm{Endoscopic Third Ventriculostomy} A brain procedure used to treat some types of hydrocephalus, in which fluid builds up because its normal path is blocked. A small camera is used to make an opening in one of the brain's fluid spaces so the fluid can drain. It may reduce the need for a shunt, but not everyone is a suitable candidate and the opening can close later.
\synonyms
ETV; Third Ventriculostomy; Endoscopic Ventriculostomy

\medterm{Endoscopic Thoracic Sympathectomy} A minimally invasive operation that interrupts selected nerves in the chest, usually for severe sweating of the hands. Possible complications include increased sweating elsewhere, drooping eyelid, collapsed lung, nerve injury, and heart or blood-pressure changes.

\medterm{Endoscopic Ultrasound} A procedure combining an endoscope with ultrasound to view the digestive-tract wall and nearby organs. It can guide needle sampling of lymph nodes, pancreas, or other masses and help assess the depth or spread of disease.

\medterm{Endoscopy} Examination of an internal body passage or cavity with a camera-equipped instrument. Depending on the site, it can directly inspect tissue, collect samples, remove growths, widen narrow areas, or stop bleeding.

\medterm{Endosome} A small membrane-bound compartment inside a cell that receives material taken in from the cell surface and sorts it for recycling, transport, or breakdown.

\medterm{Endosteum} A thin membrane lining the inner surfaces of bones, including the marrow cavity and small internal channels. It contains cells involved in bone formation, remodeling, and repair.

\medterm{Endostosis} Abnormal formation of bone on or within an existing bone. It may be a localized developmental finding or part of a broader bone disorder; the significance depends on its location, extent, and symptoms.

\medterm{Endothelial} Relating to the thin layer of cells lining blood and lymphatic vessels and some internal surfaces. These cells help control blood flow, vessel leakiness, inflammation, clotting, and new-vessel growth.

\medterm{Endothelial Cell} A cell lining a blood or lymphatic vessel. It regulates movement between blood and tissues and helps control vessel widening, inflammation, clotting, and formation of new blood vessels.

\medterm{Endothelial Nitric Oxide Synthase} A calcium- and calmodulin-dependent enzyme constitutively expressed in vascular endothelial cells that synthesizes nitric oxide from L-arginine, mediating basal vasodilation, inhibition of leukocyte adhesion, and suppression of platelet aggregation.

\synonyms
eNOS; NOS3

\medterm{Endothelin} A small protein made mainly by cells lining blood vessels that strongly narrows vessels and affects blood pressure, circulation, inflammation, and tissue growth. Excess signaling contributes to some cases of pulmonary arterial hypertension.

\medterm{Endothelin Receptor Antagonists} Medicines that block signals that narrow blood vessels, especially in the lungs. They are used for selected pulmonary arterial hypertension; liver injury, anemia, swelling, and serious fetal harm are important risks.

\medterm{Endothelium} The thin layer of specialized cells lining blood and lymphatic vessels. It controls exchange between blood and tissues and helps regulate vessel diameter, inflammation, clotting, and new-vessel growth.

\medterm{Endothelium-Derived Relaxing Factor} A substance released by vessel-lining cells that relaxes blood vessels. It is largely identified with nitric oxide, which also helps reduce platelet sticking and vessel-wall inflammation.

\synonyms
EDRF; Nitric Oxide

\medterm{Endotoxin} A component of the outer membrane of certain bacteria, especially Gram-negative bacteria. When released into the body, it can trigger fever, inflammation, low blood pressure, organ injury, and life-threatening sepsis.

\medterm{Endotracheal Intubation} Placement of a breathing tube through the mouth or nose into the windpipe to keep the airway open, deliver ventilator support, give anesthetic gases, or protect the lungs. Incorrect placement can cause severe harm and must be confirmed.

\medterm{Endotracheal Tube} A breathing tube passed through the mouth or nose into the windpipe to maintain an airway and support ventilation. An inflatable cuff may seal the airway; irritation, injury, infection, or incorrect placement can occur.

\medterm{Endourologist} A urologist specializing in minimally invasive procedures performed through the urinary tract or small access sites. Common work includes treating urinary stones, narrowing, blockage, and selected tumors.

\medterm{Endovaginal Ultrasound} Ultrasound imaging performed with a covered probe placed gently in the vagina to examine the uterus, uterine lining, ovaries, cervix, and early pregnancy. It uses sound waves rather than radiation and may cause temporary discomfort.

\medterm{Endovascular} Located inside a blood vessel or performed from within it, usually with a catheter passed through a small skin puncture. Endovascular procedures can treat aneurysms, clots, narrowing, bleeding, or abnormal vessels.

\medterm{Endovascular Aneurysm Repair} A catheter procedure that places a fabric-covered graft inside an aortic aneurysm to reinforce the artery and reduce pressure on its weakened wall. It usually recovers faster than open surgery but requires lifelong imaging for leakage or device problems.

\synonyms
EVAR; Endovascular Stent-Graft Repair; Endovascular Aortic Repair

\medterm{Endovascular Coiling} A minimally invasive, catheter-based procedure performed to treat brain aneurysms from within the blood vessels without opening the skull. Guided by real-time X-ray fluoroscopy, an interventional neurosurgeon or neuroradiologist threads a long, micro-thin catheter through an artery in the groin (femoral artery) or wrist (radial artery) all the way up into the blood vessels of the brain and directly inside the bulging aneurysm sac. Soft, ultra-fine platinum coils (each as thin as a human hair) are deployed through the catheter into the aneurysm, where they pack together and induce natural blood clot formation (thrombosis). The resulting clot seals the aneurysm from the inside, preventing blood from entering and drastically reducing the risk of aneurysm rupture and fatal intracranial bleeding.
\synonyms
Endovascular Aneurysm Coiling; Guglielmi Detachable Coiling; GDC; Aneurysm Embolization Coils

\medterm{Endovascular Embolization} A catheter procedure that deliberately blocks a selected blood vessel using coils, particles, plugs, or another material. It may control bleeding or treat an aneurysm, malformation, or tumor blood supply; non-target blockage is a possible risk.

\medterm{Endovascular Graft} A fabric-covered artificial tube placed inside a diseased blood vessel through a catheter. It can reinforce a weak artery or exclude an aneurysm; suitability and follow-up depend on anatomy, symptoms, and imaging.

\medterm{Endovascular Repair} A minimally invasive procedure that uses a catheter-delivered graft, balloon, stent, or other device to repair a blood vessel from inside. Bleeding, clotting, infection, leakage, or need for further treatment can occur.

\medterm{Endovenous Laser Ablation} A treatment for some varicose veins in which a thin laser fiber is placed inside the vein through a small needle opening. Heat closes the treated vein, and blood flows through healthier veins. It is usually done with local numbing medicine; bruising, pain, clots, nerve irritation, or return of the vein can occur.

\synonyms
EVLA; Endovenous Laser Therapy; EVLT; Endovenous Thermal Ablation

\medterm{End Point} A predefined outcome or measurement used to judge a clinical study, treatment, diagnostic test, or disease process. It may be a symptom, laboratory result, event, survival measure, or patient-reported outcome.

\medterm{End Stage} The advanced phase of a disease in which normal function is severely impaired. Meaning, criteria, expected course, and treatment options depend on the disease, such as advanced kidney, liver, lung, or heart failure.

\medterm{End-Stage Kidney Disease} Advanced chronic kidney disease in which the kidneys can no longer adequately remove waste or control fluid, electrolytes, and acid levels. Treatment options include dialysis, kidney transplantation, or supportive care based on health and preferences.

\medterm{End-Stage Renal Disease} The most advanced stage of chronic kidney disease, usually when kidney filtration is below 15 or dialysis is needed. Waste and fluid can accumulate, causing nausea, itching, confusion, swelling, anemia, or heart complications.

\synonyms
ESRD

\medterm{Endurance} The ability to continue physical or mental activity without becoming tired quickly. It depends on heart and lung function, muscle conditioning, energy, sleep, and overall health and can be assessed with activity or exercise tests.

\medterm{Enema} Introduction of liquid into the rectum through the anus to soften stool, stimulate bowel emptying, deliver selected treatment, or assist a test. Overuse can cause irritation, dehydration, or electrolyte problems.

\medterm{Energy} The capacity needed for body functions, movement, growth, and temperature control, obtained mainly from food. Persistent low energy may reflect poor sleep, inadequate nutrition, anemia, hormone disease, mood disorder, infection, or another illness.

\medterm{Energy Conservation} Strategies that reduce unnecessary effort and fatigue during daily activities, such as planning, pacing, prioritizing, resting, sitting for tasks, changing position, and using helpful equipment.

\medterm{Energy Expenditure} The energy the body uses for basic functions, movement, digestion, growth, pregnancy, and temperature control. It varies with body size, muscle mass, age, activity, illness, and other conditions.

\medterm{Energy Increased} A noticeable rise in energy or activity. It may be normal, but unusually high energy with little need for sleep, racing thoughts, impulsive behavior, or irritability can occur with stimulant use, thyroid disease, mania, or hypomania.

\medterm{Energy Intake} The amount of energy obtained from food and drink, usually measured in calories. Appropriate intake depends on age, body size, activity, illness, pregnancy, and nutritional goals.

\medterm{Energy Metabolism} The chemical processes that convert nutrients into usable energy for cells. They include breaking down carbohydrates, fats, and proteins and storing or using the resulting energy for body functions.

\medterm{Energy Resolution} The ability of a radiation detector to distinguish radiation particles with different energies. Better resolution allows more accurate identification and separation of radioactive substances during imaging or laboratory measurement.

\medterm{Energy Window} The range of radiation energies accepted by a detector or gamma camera to create an image. Choosing it appropriately improves image quality by accepting desired signals and reducing scattered radiation.

\medterm{Enfuvirtide} An injectable HIV medicine that blocks fusion of HIV with human immune cells. It is used with other antiretroviral medicines for resistant infection; injection-site reactions and allergic reactions can occur.

\medterm{Engagement} Descent of the widest part of the baby’s presenting head or body into the parent’s pelvis. It may occur before labor or during labor and indicates that the baby has entered the birth passage.

\medterm{Enhanced Elimination} Medical methods used to remove an absorbed poison more quickly, such as dialysis or changing urine chemistry. They are used only for selected toxins because benefits, timing, and risks depend on the substance and the person’s condition.

\medterm{Enhanced Pause} An unusually long pause in breathing or another measured body signal. Its significance depends on the test and setting and may require interpretation with symptoms and the complete recording.

\medterm{Enhanced Recovery After Surgery} A coordinated care plan before, during, and after surgery that supports recovery through education, appropriate nutrition and fasting, pain control, early movement, early eating, and prevention of complications.

\medterm{Enhancement Artifact} An imaging effect in which tissue behind a fluid-filled or low-density structure appears brighter than it really is, especially on ultrasound. It is a technical artifact, not necessarily a sign of increased blood flow or tissue density.

\medterm{Enhancer} A stretch of DNA that increases activity of a gene when regulatory proteins bind to it. It may be far from the gene and can work in either direction, helping control when and where the gene is active.

\medterm{Enkephalin} A natural opioid-like peptide made by the body that activates mainly delta and mu opioid receptors. It helps regulate pain, stress, mood, movement, and other nervous-system functions.

\synonyms
Leu-Enkephalin; Met-Enkephalin

\medterm{Enlarged Adenoids} Overgrowth of immune tissue behind the nose that can block nasal breathing and cause mouth breathing, snoring, repeated ear-fluid problems, or sleep-disordered breathing. Symptoms and age determine whether treatment is needed.

\medterm{Enlarged Fontanelle} A soft spot on an infant's skull that is significantly larger than normal for age, which may be associated with conditions such as congenital hypothyroidism, rickets, skeletal dysplasias, or increased intracranial pressure.

\synonyms
Large Fontanelle; Delayed Fontanelle Closure

\medterm{Enlarged Heart} An increase in heart size caused by conditions such as high blood pressure, valve disease, heart-muscle disease, anemia, or athletic adaptation. It may cause no symptoms or lead to breathlessness, swelling, rhythm problems, or heart failure.
\synonyms
Cardiomegaly

\medterm{Enlarged Liver} Alternate name; see \textit{Hepatomegaly}.

\synonyms
Hepatomegaly; Hepatic Enlargement

\medterm{Enlarged Prostate} Alternate name; see \textit{Benign Prostatic Hyperplasia}.

\medterm{Enlarged Spleen} Increase in spleen size caused by infection, liver or blood disease, inflammation, cancer, or congestion. It may cause fullness or pain and can increase rupture risk; evaluation considers symptoms, blood counts, and the underlying cause.

\synonyms
Splenomegaly; Splenic Enlargement

\medterm{Enophthalmos} Abnormal backward sinking of an eyeball within the eye socket. It may follow injury, loss of tissue around the eye, dehydration, aging, or a nerve disorder and can cause unequal eye appearance or double vision.

\medterm{Enoxaparin} A low-molecular-weight heparin that reduces clotting by enhancing antithrombin activity, especially against factor Xa. It prevents or treats blood clots but can cause bleeding, low platelets, or injection-site bruising.

\medterm{Enrichment} Adding nutrients to food after processing has removed some of them. Enriched flour, for example, commonly has certain B vitamins and iron added back to help prevent nutritional deficiencies.

\medterm{Entamoeba} A group of single-celled amoebas that live in the intestine or other tissues. Some species are harmless, while Entamoeba histolytica can cause diarrhea, intestinal ulcers, or liver abscesses.

\medterm{Entamoeba Coli} A usually harmless amoeba found in the human intestine and spread by fecal contamination. Its presence in stool indicates exposure to contaminated food or water but usually does not cause illness or require treatment.

\medterm{Entamoeba Hartmanni} A small, usually harmless intestinal amoeba spread through fecal contamination. It generally causes no disease; identification in stool mainly indicates possible exposure to contaminated food, water, or hands.

\medterm{Entamoeba Histolytica} A parasite spread when contaminated food, water, or hands are swallowed. It can cause diarrhea, abdominal pain, and bloody stools and may invade the intestine or spread to the liver, causing an abscess.

\medterm{Entamoeba Polecki} An intestinal amoeba that can infect people after fecal contamination, often through contact with animals or contaminated food and water. It is usually harmless but may occasionally be associated with diarrhea.

\medterm{Enteral} Relating to or delivered through the digestive tract. Enteral nutrition provides food and fluids by mouth or through a feeding tube into the stomach or intestine.

\medterm{Enteral Formula} A liquid nutrition preparation given by mouth or feeding tube. It may provide complete nutrition or selected nutrients and is chosen according to digestion, absorption, fluid needs, kidney or liver disease, and metabolic requirements.

\medterm{Enteral Nutrition} Delivery of food, fluids, and nutrients through the digestive tract, by mouth or feeding tube. It may be used when swallowing or eating is unsafe or inadequate; aspiration, diarrhea, blockage, and tube displacement are possible complications.

\medterm{Enteric} Relating to the intestine or digestive tract. The term can also describe organisms, infections, medicines, or coatings associated with the intestines.

\medterm{Enteric Bacteria} Bacteria that live in or affect the intestinal tract. Some are harmless normal residents, while others cause diarrhea or infections outside the intestine, such as urinary, wound, or bloodstream infection.

\medterm{Enteric-Coated} Covered with a protective layer that resists stomach acid and dissolves in the small intestine. This can protect the stomach or protect a medicine from stomach acid; such tablets should not usually be crushed.

\medterm{Enteric Duplication Cyst} A birth-related cyst attached to or within the digestive tract, most often the small intestine. It may cause pain, blockage, bleeding, infection, or no symptoms and can require surgery.

\medterm{Enteric Fever} A serious infection caused by \textit{Salmonella} Typhi or \textit{Salmonella} Paratyphi, usually spread through food or water contaminated with human waste. It can cause gradually rising fever, headache, stomach pain, weakness, diarrhea or constipation, and sometimes a rash. Untreated illness can cause bleeding or a hole in the intestine; testing and antibiotic treatment are needed.

\synonyms
Typhoid Fever; Paratyphoid Fever; Salmonella Enteric Infection

\medterm{Enteric Nervous System} The network of nerve cells within the digestive tract that helps control movement, secretion, blood flow, and communication with the brain and spinal cord.

\medterm{Enteritis} Inflammation of the small intestine caused by infection, inflammatory disease, radiation, reduced blood flow, medicines, or other injury. It may cause abdominal pain, diarrhea, vomiting, fever, bleeding, or dehydration.

\medterm{Enterobacter} A group of Gram-negative bacteria that may live in the intestine but can cause healthcare-associated urinary, bloodstream, wound, respiratory, or abdominal infection, especially in vulnerable people.

\medterm{Enterobacter Aerogenes} A former name for Klebsiella aerogenes, a Gram-negative bacterium that can cause urinary, lung, wound, abdominal, or bloodstream infection, especially in healthcare settings. It may resist several antibiotics.

\medterm{Enterobacter Cloacae} A Gram-negative bacterium that can cause healthcare-associated urinary, bloodstream, lung, wound, or abdominal infection. It may resist several antibiotics, so laboratory susceptibility testing guides treatment.

\medterm{Enterobacter Hormaechei} A bacterium in the Enterobacter cloacae group that can cause opportunistic healthcare-associated infection of blood, urine, lungs, wounds, or devices. Testing for antibiotic resistance is often important.

\medterm{Enterobacteriaceae} A family of Gram-negative bacteria that includes Escherichia, Klebsiella, Enterobacter, Salmonella, and Shigella. Some normally live in the intestine, while others cause intestinal, urinary, wound, or bloodstream infection.

\medterm{Enterobacteriaceae Carbapenem Resistance} Resistance in certain Gram-negative bacteria to carbapenem antibiotics, often caused by enzymes that destroy these medicines. Such infections can be difficult to treat and require laboratory testing, infection control, and specialist guidance.

\medterm{Enterobiasis} Infection with pinworms, small intestinal worms whose eggs spread easily between people. It commonly causes intense itching around the anus at night, disturbed sleep, and sometimes irritation in the genital area.

\medterm{Enterobius} A group of pinworms that live in the intestine. Enterobius vermicularis causes the common pinworm infection, which spreads when eggs are swallowed and often causes nighttime anal itching.

\medterm{Enterobius Vermicularis} A small intestinal worm commonly called the pinworm or threadworm. Adult females lay eggs around the anus, usually at night, causing itching and allowing infection to spread through contaminated hands or surfaces.

\synonyms
Pinworm; Threadworm; Enterobius

\medterm{Enterocele} A pelvic organ prolapse in which a pouch of abdominal lining and sometimes small intestine bulges into the space between the vagina and rectum. It may cause pelvic pressure, a vaginal bulge, back or rectal discomfort, or no symptoms.

\synonyms
Small Bowel Prolapse; Douglas Pouch Hernia

\medterm{Enterochromaffin Cells} Specialized enteroendocrine cells located throughout the mucosal epithelium of the gastrointestinal tract that synthesize, store, and release serotonin (5-hydroxytryptamine) and substance P to modulate gastrointestinal motility, secretion, and visceral pain perception.
\synonyms
EC Cells; Kultschitzky Cells

\medterm{Enterochromaffin-Like Cells} Specialized neuroendocrine cells situated within the oxyntic mucosa of the gastric fundus and body that synthesize and secrete histamine in response to gastrin and acetylcholine stimulation, directly stimulating parietal cell H+/K+-ATPase proton pumps.

\synonyms
ECL Cells

\medterm{Enteroclysis} An imaging test in which contrast material is delivered through a tube into the small intestine to look for inflammation, narrowing, tumors, or blockage. CT or MRI enterography is now more commonly used in many settings.

\medterm{Enterococcaceae} A family of bacteria that includes Enterococcus species. These bacteria may normally live in the intestine but can cause urinary, wound, bloodstream, or heart-valve infections when they enter protected body sites.

\medterm{Enterococcal Infection} An infection caused by Enterococcus bacteria, which normally live in the intestine but can cause urinary-tract, wound, abdominal, bloodstream, or heart-valve infections. Enterococcal infections are more common after hospitalization, surgery, antibiotic exposure, or use of medical devices. Some strains resist multiple antibiotics, so culture and susceptibility testing guide treatment.

\synonyms
Enterococcosis; Enterococcal Sepsis

\medterm{Enterococcus} A group of bacteria normally found in the intestine that can cause urinary, wound, abdominal, bloodstream, or heart-valve infection. Some strains resist several antibiotics and require laboratory susceptibility testing.

\medterm{Enterococcus Durans} A bacterium that may live in the bowel but can rarely cause infection in seriously ill or immunocompromised people. It may affect blood, urine, wounds, or other sites, and testing guides antibiotic treatment.

\medterm{Enterococcus Faecalis} A common intestinal bacterium that can cause urinary, wound, abdominal, bloodstream, or heart-valve infection. Disease is more likely after surgery, hospitalization, medical-device use, or when immunity is weakened.

\medterm{Enterococcus Faecium} An Enterococcus species that can cause healthcare-associated urinary, wound, abdominal, bloodstream, or heart-valve infection. Some strains resist vancomycin and other antibiotics, making laboratory testing essential.

\medterm{Enterocolitis} Inflammation of the small intestine and colon caused by infection, immune disease, reduced blood flow, toxins, or other injury. It may cause abdominal pain, diarrhea, fever, vomiting, bleeding, dehydration, or severe illness.

\medterm{Enterocyte} A cell lining the small intestine that digests and absorbs nutrients, water, and electrolytes. It also helps form the intestinal barrier and participates in immune defense.

\medterm{Enterocytozoon} A group of tiny spore-forming parasites that can cause persistent diarrhea and intestinal disease, especially in people with weakened immunity. Some species can also affect other organs.

\medterm{Enteroendocrine Cell} A specialized cell in the intestinal lining that senses food and releases hormones regulating digestion, appetite, gallbladder and pancreatic activity, and blood-glucose control.

\medterm{Enterohepatic Circulation} Recycling of a substance from the liver into bile and the intestine, then back to the liver through the bloodstream. This can prolong the effects of some medicines and natural body chemicals.

\medterm{Enteromonas Hominis} A usually harmless microorganism that may live in the intestine and appear in stool samples. It generally does not cause disease, so its presence should be interpreted with symptoms and other test results.

\medterm{Enteropathic} Relating to disease of the intestine or caused by intestinal disease. The term may describe arthritis, skin changes, protein loss, or other problems occurring with inflammatory bowel disease or impaired intestinal absorption.

\medterm{Enteropathic Arthritis} Inflammatory joint disease associated with inflammatory bowel disease such as Crohn disease or ulcerative colitis. It may affect the spine, joints between the spine and pelvis, or joints of the arms and legs.

\medterm{Enteropathy} Any disease or disorder affecting the intestine. Causes include infection, inflammation, gluten sensitivity, poor blood flow, toxins, medicines, inherited conditions, or impaired absorption.

\medterm{Enteropathy-Associated T-Cell Lymphoma} A rare, fast-growing cancer of immune cells in the intestine. It can be linked to long-standing celiac disease and may cause severe abdominal pain, weight loss, poor nutrient absorption, blockage, or a hole in the bowel.

\synonyms
EATL; Enteropathy-Type Intestinal T-Cell Lymphoma

\medterm{Enteroscopy} Examination of the small intestine with a camera-equipped instrument. It can identify disease, obtain tissue samples, treat bleeding, or remove selected growths beyond the reach of routine upper endoscopy.

\medterm{Enterospasm} An involuntary tightening or contraction of the intestine that may cause cramping, abdominal pain, bloating, or altered movement of intestinal contents.

\medterm{Enterostomy} A surgically created opening that brings part of the intestine to the abdominal surface so stool or intestinal contents can leave the body or a feeding solution can enter.

\medterm{Enterostomy Procedure} Surgery that creates an opening from the intestine to the abdominal wall, such as an ileostomy or colostomy. It may divert stool or support feeding; bleeding, infection, blockage, and skin problems can occur.

\medterm{Enterotoxin} A toxin made by some microorganisms that acts on the intestinal lining and can cause vomiting, cramps, watery diarrhea, dehydration, or other intestinal illness.

\medterm{Enterovesical Fistula} An abnormal connection between the intestine and urinary bladder. It can cause repeated urinary infections, air or stool in the urine, abdominal pain, or urine leaking into the intestine and may result from inflammation, cancer, or surgery.

\medterm{Enterovirulent E. Coli} A group of Escherichia coli strains that cause intestinal disease by making toxins, attaching to the bowel, or invading it. Illness may include watery or bloody diarrhea, cramps, vomiting, or dehydration.

\medterm{Enterovirus} A group of viruses spread mainly by respiratory secretions or contact with infected stool. Many infections are mild or cause no symptoms, but some cause hand-foot-and-mouth disease, rash, meningitis, heart inflammation, or severe illness in newborns.

\medterm{Enterovirus D68} A virus that usually causes a respiratory illness but can rarely cause severe pneumonia or sudden weakness of the limbs, especially in children. It spreads mainly through respiratory secretions and contaminated surfaces.

\medterm{Enterovirus Infection} Alternate name; see \textit{Enterovirus}.

\medterm{Enthesis} The place where a tendon, ligament, or joint capsule attaches to bone. Repeated stress or inflammation there can cause localized pain and stiffness and occurs in several inflammatory joint diseases.

\medterm{Enthesitis} Inflammation where a tendon, ligament, or joint capsule attaches to bone. It can cause localized pain, tenderness, and stiffness and is common in spondyloarthritis, including psoriatic and inflammatory-bowel-disease-related arthritis.

\medterm{Enthesitis And Enthesopathy} Enthesitis is inflammation where a tendon, ligament, or fascia attaches to bone. Enthesopathy includes inflammation, wear, injury, or other disease at that site and may cause localized pain and stiffness.

\synonyms
Enthesopathy; Insertion Tendinopathy

\medterm{Enthesitis-Related Arthritis} A type of arthritis that begins in childhood and causes inflammation where tendons or ligaments attach to bone. It often affects the legs, feet, or lower back and may also cause swelling in joints or pain where the heel tendon attaches.

\synonyms
ERA; Juvenile Spondyloarthritis

\medterm{Entorhinal Cortex} A region of the brain that links the hippocampus with other memory and navigation areas. It helps process spatial information and new memories and is affected early in many cases of Alzheimer disease.

\medterm{ENT Physician} An otolaryngologist, a physician who diagnoses and treats disorders of the ear, nose, throat, and related head and neck structures.

\medterm{Entrance Wound} The wound where a projectile enters the body. Its appearance may help estimate the direction and distance of the shot, but findings vary with the weapon, clothing, body site, and circumstances and require forensic examination.

\medterm{Entropion} Inward turning of an eyelid, usually the lower lid, so the eyelashes rub against the eye. It can cause irritation, tearing, pain, corneal scratches, infection, and reduced vision.

\medterm{Enucleation} The surgical removal of an entire organ or well-defined encapsulated tumor in one piece without cutting into surrounding healthy tissue. In ophthalmology, it specifically refers to the complete surgical removal of the eyeball (globe) and optic nerve stump while preserving the eyelids, orbital fat, and extraocular muscles. It is performed for intraocular malignancies (e.g., retinoblastoma, uveal melanoma), severe irreparable ocular trauma, endophthalmitis, or a blind, painfully disfigured eye. An orbital implant is placed and attached to the eye muscles to support a prosthetic eye (artificial eye). In general surgery, it describes the clean ``shelling out'' of an intact benign tumor (such as a uterine fibroid, insulinoma, or cyst).
\synonyms{Ocular Enucleation; Eye Enucleation; Surgical Enucleation}

\medterm{Enucleation Procedure} Alternate name; see \textit{Enucleation}.

\medterm{Enuresis} Repeated involuntary urination, usually during sleep in a child expected to remain dry. It may be primary or begin again after dryness; constipation, infection, diabetes, sleep problems, stress, or urinary disorders may contribute.

\medterm{Environment} The physical, chemical, biological, and social surroundings in which a person lives or works. Environmental factors such as air, water, housing, noise, climate, and pollutants can influence health.

\medterm{Environmental Exposure} Contact with a chemical, pollutant, organism, physical hazard, or other factor in the surroundings. Health effects depend on the substance, dose, route, duration, timing, and individual susceptibility.

\medterm{Environmental Health} The field concerned with how surroundings affect human health. It includes air and water quality, food safety, housing, workplace hazards, climate, radiation, pollution, and prevention of harmful exposures.

\medterm{Environmental Illness} Illness caused or worsened by an environmental exposure. Evaluation considers the timing, dose, setting, symptoms, other possible causes, and whether symptoms improve when exposure stops.

\medterm{Enzyme} A protein, or occasionally an RNA molecule, that speeds a specific chemical reaction without being used up. Enzymes are essential for digestion, energy production, DNA processes, and many other cell functions.

\medterm{Enzyme Activity} The rate at which an enzyme carries out a chemical reaction under specified conditions. It is affected by temperature, acidity, inhibitors, helper molecules, substrate amount, and enzyme concentration.

\medterm{Enzyme Deficiency} Reduced or absent activity of an enzyme, usually inherited or sometimes acquired. It can cause buildup of substances, lack of needed products, and characteristic metabolic or cellular disease.

\medterm{Enzyme Induction} Increased production or activity of an enzyme, often after exposure to a medicine or chemical. If liver drug-metabolizing enzymes are induced, some medicines may be broken down faster and become less effective.

\medterm{Enzyme Inhibition} Reduced activity of an enzyme, sometimes caused by a medicine or chemical. If a drug-metabolizing enzyme is inhibited, another medicine may accumulate and cause stronger effects or toxicity.

\medterm{Enzyme Kinetics} The study of how quickly enzyme reactions occur and how their rate changes with substrate amount, temperature, acidity, inhibitors, helper molecules, and enzyme concentration.

\medterm{Enzyme-Linked Immunosorbent Assay} A laboratory test that uses antibodies and an enzyme-generated color signal to detect or measure a target such as an antibody, antigen, hormone, or infection marker.

\medterm{Enzyme Markers} Enzymes measured in blood or another body fluid to help assess tissue injury or organ function. Abnormal levels can suggest liver, bile-duct, muscle, heart, bone, or pancreatic disease but are not a diagnosis by themselves.

\medterm{Enzyme Precursors} Inactive or weakly active proteins that become working enzymes after being cut or chemically changed. This allows the body to activate digestive and other enzymes at the correct place and time.

\medterm{Enzyme Therapy} Treatment that replaces, supplements, blocks, or uses an enzyme to correct a deficiency or alter a disease process. The purpose, dose, benefits, and risks depend on the specific enzyme and condition.

\medterm{Eosinophil} A type of white blood cell involved in allergic inflammation and defense against some parasites. Eosinophils normally make up a small fraction of blood cells; high levels can occur with allergy, infection, medicines, or other disease.

\medterm{Eosinophilia} An abnormally high number of eosinophils in the blood, commonly caused by allergy, parasites, medicines, autoimmune disease, or some blood disorders. Severe or persistent elevation can injure the heart, lungs, skin, nerves, or other organs.

\medterm{Eosinophilic Asthma} Asthma in which eosinophils, a type of white blood cell, contribute to airway inflammation. It may cause recurring wheezing, cough, chest tightness, or breathlessness and can be difficult to control in some people.

\medterm{Eosinophilic Bronchitis} A cause of long-lasting cough in which eosinophils inflame the airways without the variable narrowing typical of asthma. The cough may resemble asthma or reflux and needs assessment to identify the cause.

\medterm{Eosinophilic Esophagitis} A chronic immune-related disease in which eosinophils inflame the esophagus. It may cause difficulty swallowing, food becoming stuck, chest or upper-abdominal pain, reflux-like symptoms, and narrowing of the esophagus.

\medterm{Eosinophilic Fasciitis} A rare inflammatory disease causing painful swelling and hardening of the deeper skin and tissue around muscles, usually in the arms or legs. Blood eosinophils may be high, and joint stiffness or limited movement can develop.

\synonyms
Shulman Syndrome; Diffuse Fasciitis with Eosinophilia

\medterm{Eosinophilic Gastroenteritis} A rare immune-related disease in which eosinophils inflame the stomach or intestines. Symptoms may include abdominal pain, nausea, vomiting, diarrhea, bleeding, poor absorption, or weight loss.

\medterm{Eosinophilic Gastrointestinal Disorders} Conditions in which eosinophils, a type of white blood cell, build up in one or more parts of the digestive tract and cause inflammation. Symptoms may include pain, vomiting, swallowing difficulty, diarrhea, or poor growth; diagnosis usually requires endoscopy and biopsy.

\medterm{Eosinophilic Granuloma} A localized form of Langerhans-cell disease in which abnormal immune cells collect, often in bone or skin. Bone lesions can cause pain, swelling, or fractures and may require imaging and tissue examination.

\medterm{Eosinophilic Granulomatosis With Polyangiitis} A rare disease in which inflamed blood vessels occur with asthma and high levels of eosinophils, a type of white blood cell. It can damage the lungs, skin, nerves, heart, kidneys, or digestive organs.

\synonyms
Allergic Granulomatosis
Allergic Granulomatosis and Angiitis

\medterm{Eosinophils} White blood cells that help control allergic inflammation and defend against some parasites. High levels may occur with allergies, asthma, infection, medicines, autoimmune disease, or blood disorders; low levels usually have little significance.

\synonyms
Eosinophil Granulocytes

\medterm{Ependyma} A specialized layer of cells lining the brain’s fluid-filled ventricles and the spinal cord’s central canal. It helps separate nervous tissue from cerebrospinal fluid and supports movement of that fluid.

\medterm{Ependymal Cell} A supporting cell lining the brain ventricles and spinal cord central canal. It helps maintain the cerebrospinal-fluid boundary and, in some areas, helps move or produce the fluid.

\medterm{Ependymal Tumor} A growth arising from cells lining the brain ventricles or spinal canal. It may block cerebrospinal-fluid flow and cause headache, vomiting, balance problems, weakness, or hydrocephalus; treatment depends on its type and extent.

\medterm{Ependymoblastoma} A rare, aggressive brain tumor of childhood made from very immature cells. It may cause headache, vomiting, seizures, or changes in movement or development; tissue and genetic testing guide treatment and outlook.

\medterm{Ependymoma} A tumor arising from cells lining the brain ventricles or spinal canal. It may block fluid flow and cause headache, vomiting, seizures, balance problems, weakness, or hydrocephalus, especially in children.

\medterm{Ephedrine} A stimulant that activates the sympathetic nervous system and can widen airways and raise blood pressure and heart rate. It has limited medical uses and may cause palpitations, insomnia, anxiety, high blood pressure, or misuse.

\medterm{Ephedrinum} Another name for ephedrine, a stimulant that can increase heart rate, blood pressure, alertness, and airway diameter. Excessive use may cause dangerous heart-rhythm, blood-pressure, or nervous-system effects.

\medterm{Ephelides} Small, flat freckles caused by increased pigment production, often inherited. They become darker with sunlight and usually fade when exposure decreases; sun protection helps prevent further pigment change and skin damage.

\medterm{Ephrins} Proteins on the surface of cells that help guide cell position, nerve growth, blood-vessel development, and tissue organization. Abnormal ephrin signaling can contribute to some diseases and cancers.

\medterm{Epiblast} A layer of cells in the early embryo that develops into the three main cell layers of the body and ultimately forms most tissues and organs.

\medterm{Epicanthal Folds} Skin folds of the upper eyelid that partly cover the inner corner of the eye. They are normal in many people but can also occur with some chromosome or developmental conditions.

\medterm{Epicardium} The thin outer layer of the heart wall and the inner layer of its surrounding protective sac. It contains connective tissue, fat, blood vessels, and nerves and helps reduce friction as the heart beats.

\medterm{Epichlorohydrin} An industrial chemical used to make resins and other materials. Exposure can irritate or burn the skin, eyes, and airways; swallowing or substantial exposure can cause serious poisoning and requires urgent medical advice.

\medterm{Epicondyle} A bony projection beside a rounded joint surface where muscles and ligaments attach. Repeated strain or injury around an elbow epicondyle can cause localized pain and weakness.

\medterm{Epicondylitis} Pain and tenderness where a tendon attaches near an epicondyle, usually from repeated strain rather than infection. At the elbow it can weaken gripping and wrist movement and is commonly called tennis or golfer’s elbow.

\synonyms
Tennis Elbow; Golfer’s Elbow

\medterm{Epicondylitis, Lateral} Alternate name; see \textit{Tennis Elbow}.

\medterm{Epidemic} An increase in disease cases above the number normally expected in a population, place, or period. An epidemic may result from infection, environmental exposure, behavior, or another shared cause.

\medterm{Epidemic Keratoconjunctivitis} A very contagious viral infection of the eye. It can cause red eyes, pain, watering, swollen eyelids, light sensitivity, blurred vision, and swollen glands near the ear; symptoms may last for weeks.

\synonyms
EKC; Adenoviral Keratoconjunctivitis; Shipyard Eye

\medterm{Epidemiological Study} A study of health and disease in groups of people. It may measure how often a condition occurs, compare possible causes or exposures, or examine outcomes over time to help guide prevention and healthcare planning.

\synonyms
Epidemiologic Study

\medterm{Epidemiologic Method} A structured way to measure how often a health problem occurs and to study whether exposures, behaviors, or other factors are associated with it in a population. Results guide prevention and policy.

\medterm{Epidemiologic Studies} Research that examines the distribution, causes, risk factors, prevention, or outcomes of health events in populations.

\medterm{Epidemiologist} A scientist or public-health professional who studies the distribution, causes, risk factors, and prevention of health conditions in populations.

\medterm{Epidemiology} The study of how often diseases and other health events occur, who is affected, and which factors influence them. It uses population data to identify causes, guide prevention, investigate outbreaks, and evaluate healthcare.

\medterm{Epidermal} Relating to the epidermis, the outermost layer of skin. An epidermal problem affects this protective surface layer rather than the deeper skin tissues.

\medterm{Epidermal Acanthosis} Thickening of the epidermis caused by increased growth of its cells, seen in conditions such as psoriasis, chronic rubbing, and some inflammatory or genetic skin disorders.

\medterm{Epidermal Cells} Cells forming the outer layer of skin, chiefly keratinocytes, which create a protective barrier; melanocytes, immune cells, and touch-associated cells are also present.

\medterm{Epidermal Growth Factor} A signaling protein that binds the epidermal-growth-factor receptor and promotes cell growth, survival, and tissue repair; abnormal signaling can contribute to cancer.

\medterm{Epidermal Inclusion Cyst} A common, usually harmless lump beneath the skin filled with keratin, a protein found in skin. It often has a small central opening and may become swollen, painful, infected, or recur if its wall remains.

\medterm{Epidermal Layers} The epidermis has several layers of cells that make and protect the skin. From deepest to most superficial they are basal, spinous, granular, clear, and horny layers; the clear layer is mainly in the palms and soles.

\medterm{Epidermal Nevus} A usually benign, often present-from-birth or childhood skin lesion caused by an overgrowth of epidermal cells. It commonly appears as a localized, linear, or warty thickened patch and may occur with other developmental abnormalities.

\medterm{Epidermis} The outer layer of skin that forms a protective barrier against injury, germs, and water loss. It contains cells that make keratin and pigment, as well as immune and touch-sensing cells, and it is replaced continuously.

\medterm{Epidermodysplasia Verruciformis} A rare inherited disorder of skin immunity associated with susceptibility to certain human papillomavirus infections, widespread wart-like lesions, and increased risk of some skin cancers.

\medterm{Epidermoid Cyst} Alternate index label; see \textit{Epidermoid Cysts}.

\medterm{Epidermoid Cysts} Slow-growing, usually harmless sacs beneath the skin filled with keratin. They often appear as movable lumps and may become painful, red, swollen, or infected if they rupture.
\synonyms
Epidermal Inclusion Cysts; Infundibular Cysts

\medterm{Epidermolysis Bullosa} A group of inherited disorders in which the skin blisters or tears after minor rubbing or injury. Severity varies; some forms heal without scars, while others cause scarring, tight joints, poor nutrition, infection, or increased skin-cancer risk.

\medterm{Epidermolysis Bullosa Acquisita} An autoimmune blistering disease in which the body attacks a protein that helps hold skin layers together. The skin becomes fragile and blisters easily after minor injury; scars may form as it heals.

\medterm{Epidermophyton} A group of fungi that grow in keratin, a tough protein in skin, hair, and nails. They can cause ringworm infections, including itchy, scaly skin disease and fungal infection of the nails, but do not usually infect living hair shafts.

\medterm{Epidermophyton Floccosum} A dermatophyte fungus that commonly causes tinea of the body, groin, and feet and can infect nails.

\medterm{Epididymal Cyst} A usually benign, fluid-filled cyst arising in the epididymis. It may be painless and may require treatment only if it causes significant symptoms or complications.

\medterm{Epididymal Neoplasm} A benign or malignant growth arising in the epididymis; most palpable epididymal masses are benign, but persistent or enlarging lesions require evaluation.

\medterm{Epididymides} The paired coiled ducts on the posterior surface of the testes where sperm mature and are stored before entering the vasa deferentia.

\synonyms
Epididymal Ducts; Convoluted Seminiferous Ducts

\medterm{Epididymis} A tightly coiled tube along the back of each testicle where sperm mature, gain the ability to move, and are stored before passing into the vas deferens. It has a head, body, and tail.

\medterm{Epididymitis} Inflammation of the epididymis, usually caused by infection. It commonly causes gradually developing pain and swelling on one side of the scrotum. Causes include sexually transmitted infections, urinary bacteria, urine-flow problems, and less commonly injury or medicines.

\medterm{Epididymitis And Orchitis} Alternate name for epididymo-orchitis, inflammation of the epididymis and testicle.

\medterm{Epididymo-Orchitis} Inflammation of the epididymis and testicle, usually caused by an infection. It can cause one-sided scrotal pain, swelling, redness, painful urination, and fever. The cause may be a sexually transmitted infection or bacteria from the urinary tract, depending on the person and situation. Sudden severe pain needs emergency assessment to rule out testicular torsion; treatment depends on the cause.

\synonyms
Epididymo-orchitis; Testicular and Epididymal Inflammation; Orchiepididymitis

\medterm{Epidural} Pain relief or regional anesthesia given through a small catheter placed in the lower back, near the spinal nerves. It is commonly used during labor and some operations; possible effects include temporary low blood pressure, itching, headache, or incomplete relief.

\medterm{Epidural Abscess} A pocket of pus between the tough outer covering of the brain or spinal cord and nearby bone. It can compress the brain, spinal cord, or nerves, causing fever, severe pain, weakness, numbness, or loss of bladder control and requires emergency treatment.

\medterm{Epidural Analgesia} Pain relief produced by delivering numbing medicine, sometimes with an opioid, through a small catheter into the space around the spinal nerves. It may be used during labor or surgery while the person remains awake. Low blood pressure, itching, headache, infection, or incomplete relief can occur.

\medterm{Epidural Analgesia For Labor} Pain relief during labor provided through a small catheter placed in the lower back. Numbing medicine, sometimes with an opioid, is given gradually to reduce pain while allowing the laboring person to remain awake and participate in birth. Temporary low blood pressure, itching, headache, or incomplete relief can occur.

\medterm{Epidural Anesthesia} Regional anesthesia produced by injecting numbing medicine, sometimes with an opioid, into the space around the spinal nerves. It blocks pain from part of the body while the person stays awake; low blood pressure, headache, infection, or nerve injury can rarely occur.

\synonyms
Epidural Anaesthesia

\medterm{Epidural Blood Patch} A procedure that injects a small amount of the person’s own blood into the space around the spinal nerves to seal a leak after spinal anesthesia or an accidental puncture. It usually treats a severe positional headache; back pain, infection, or another puncture can occur.

\medterm{Epidural Hematoma} A collection of blood between the skull and the brain’s tough outer covering, usually after a head injury. It can rapidly compress the brain, causing headache, vomiting, confusion, weakness, loss of consciousness, or death and requires emergency treatment.

\medterm{Epidural Neoplasm} A tumor in the space outside the dura surrounding the spinal cord or brain; spinal epidural tumors can compress nerves or the spinal cord and cause pain, weakness, or sensory loss.

\medterm{Epidural Space} The space between the tough outer covering of the spinal canal and the canal wall. Medicines injected there can reduce pain during childbirth, surgery, or some pain treatments.

\medterm{Epidural Steroid Injection} A minimally invasive outpatient spine procedure used to relieve severe radiating back, neck, arm, or leg pain (such as sciatica or lumbar radiculopathy) caused by an inflamed or compressed spinal nerve. Under live X-ray guidance (fluoroscopy) and local anesthesia, a pain specialist or spine clinician carefully directs a thin needle into the epidural space surrounding the spinal nerves and dural sac, injects a tiny amount of contrast dye to verify accurate placement, and administers a powerful long-acting corticosteroid combined with a local anesthetic. The medication coats the irritated nerve roots, flushing away inflammatory chemicals and dramatically reducing swelling caused by herniated discs, spinal stenosis, degenerative disc disease, or bone spurs. Epidural steroid injections can be performed through translaminar, transforaminal, or caudal routes, offering substantial pain relief that allows patients to engage in active physical therapy and frequently avoid spinal surgery.

\synonyms
ESI; Epidural Injections for Back Pain; Translaminar Epidural Injection; Transforaminal Epidural Injection; Caudal Epidural Steroid Injection

\medterm{Epigastric} Relating to the upper middle area of the abdomen between the ribs. Pain there may arise from the stomach, esophagus, pancreas, gallbladder, intestine, or other organs.

\medterm{Epigastric Artery} An artery supplying the front abdominal wall. The superior and inferior epigastric arteries provide blood to abdominal muscles and skin and are important landmarks during some procedures.

\medterm{Epigastric Pain} Pain or discomfort in the upper middle abdomen, below the ribs and above the navel. Causes include indigestion, ulcer, reflux, gallbladder disease, pancreatitis, and other conditions; severe or persistent pain needs assessment.

\medterm{Epigastric Region} The upper middle section of the abdomen between the rib margins and above the navel. It overlies parts of the stomach, duodenum, pancreas, and liver; pain there can have several digestive or vascular causes.

\medterm{Epigastrium} The upper middle part of the abdomen, below the breastbone and between the ribs. It overlies parts of the stomach, liver, pancreas, and small intestine; pain there has many possible causes.

\medterm{Epigenetics} Changes in how genes are switched on or off without changing the DNA sequence itself. Chemical tags and DNA-packaging changes can be influenced by development, environment, and illness and may sometimes be reversible.

\medterm{Epigenome} The collection of chemical tags and DNA-packaging features that help control which genes are active in a cell without changing the DNA sequence.

\medterm{Epigenomics} The study of chemical and structural changes that control which genes are active across the genome without changing the DNA sequence. These changes can influence development, disease, and treatment response.

\medterm{Epiglottic Vallecula} A shallow mucosal depression situated in the glossoepiglottic space between the base of the tongue and the anterior surface of the epiglottis, bounded by median and lateral glossoepiglottic folds, serving as the anatomical seat for the tip of a curved Macintosh laryngoscope blade during endotracheal intubation.
\synonyms
Vallecula Epiglottica; Vallecular Fossa

\medterm{Epiglottic Vallecula} A shallow mucosal depression or pocket located at the base of the tongue immediately anterior to the epiglottis, bounded medially by the median glossoepiglottic fold and laterally by lateral glossoepiglottic folds; serves as an anatomical landmark for laryngoscope blade placement during intubation.

\synonyms
Vallecula

\medterm{Epiglottis} A leaf-shaped flap of flexible tissue at the base of the tongue that helps direct food and liquid into the esophagus during swallowing and away from the airway. Swelling can rapidly obstruct breathing.

\medterm{Epiglottitis} A potentially life-threatening infection or inflammation causing swelling around the epiglottis and upper airway. Fever, severe sore throat, drooling, trouble swallowing, noisy breathing, or sitting forward requires emergency care; the throat should not be examined forcefully.

\medterm{Epilepsia Partialis Continua} Repeated or nearly continuous twitching in one part of the body, often an arm, leg, or side of the face. Awareness may remain normal, but the condition can last for hours or longer and needs urgent medical assessment.

\synonyms
EPC; Kojewnikoff Syndrome

\medterm{Epilepsy} A chronic brain disorder that causes a lasting tendency to have recurrent unprovoked seizures. A seizure is a brief episode caused by excessive, abnormal electrical activity in the brain and may change movement, sensation, awareness, behavior, or consciousness. Seizures may begin in one area of the brain (focal seizures) or involve both sides from the start (generalized seizures). Causes include genetic conditions, brain injury or malformation, stroke, infection, metabolic or immune disease, or no known cause. A single seizure, especially one triggered by a temporary illness or abnormal body chemistry, does not always mean epilepsy.

\synonyms
Seizure Disorder

\medterm{Epilepsy, Frontal Lobe} Alternate name; see \textit{Frontal Lobe Epilepsy}.

\medterm{Epilepsy Surgery} Surgery or a procedure used for selected people whose seizures continue despite appropriate medicines and arise from an area that can be safely treated. Detailed testing identifies the seizure source and estimates benefits and risks.

\medterm{Epileptic Spasms} Brief seizures, most common in infants, in which the head, body, arms, or legs suddenly bend or stiffen. They often occur in clusters and need prompt assessment because they can affect development.

\synonyms
Infantile Spasms

\medterm{Epileptiform} Resembling a seizure or the electrical pattern associated with seizure activity. It may describe movements or an EEG finding but does not by itself prove that a person has epilepsy.

\medterm{Epimysium} A substantial sheath of dense irregular collagenous connective tissue that completely envelops the entire exterior surface of a skeletal muscle, continuous with tendons, muscular fascial planes, and the periosteum of bone.

\medterm{Epinephrin} Alternate spelling; see \textit{Epinephrine}.

\medterm{Epinephrine} A hormone and neurotransmitter released during stress that increases heart activity, opens the airways, and narrows many blood vessels. As a medicine, it treats anaphylaxis and selected cardiac or breathing emergencies.

\synonyms
Adrenaline

\medterm{Epineurium} The thick, outermost fibrous connective tissue sheath that encloses an entire peripheral nerve, surrounding multiple nerve fascicles, longitudinal blood vessels (vasa nervorum), and adipose tissue to provide tensile mechanical protection.

\medterm{Epiphora} Excessive tearing caused by blocked tear drainage or irritation that triggers extra tears. Causes include dry eye, allergy, infection, eyelid position problems, and blockage of the tear duct.

\medterm{Epiphyseal Plate} The growth plate, a layer of cartilage near the end of a growing bone that allows the bone to lengthen. It gradually closes after puberty; injury can cause pain, deformity, or unequal growth.

\synonyms
Growth Plate; Physis

\medterm{Epiphyseal Plate Disorders} Conditions affecting a growing bone’s growth plate, including fractures, inflammation, infection, abnormal blood supply, premature closure, and disorders in which an epiphysis shifts or separates. They may cause pain, swelling, limping, deformity, or unequal limb growth.

\medterm{Epiphysis} The expanded end of a long bone, usually covered by smooth cartilage where it forms a joint. It contains mostly spongy bone and helps determine the shape and movement of the joint.

\medterm{Epiphysitis} Inflammation or injury involving a growth plate or the area near the end of a growing bone. It may cause localized pain, swelling, or altered growth, and the cause depends on age and the affected bone.

\medterm{Epiploic Foramen} The narrow anatomical communication aperture connecting the greater peritoneal sac with the lesser sac (omental bursa), bounded anteriorly by the hepatoduodenal ligament (containing the portal triad), posteriorly by the inferior vena cava, superiorly by the caudate lobe of the liver, and inferiorly by the first part of the duodenum.

\synonyms
Foramen of Winslow

\medterm{Epipodophyllotoxin} A group of cancer medicines derived from a plant compound that damage DNA by blocking an enzyme needed for cell division. They treat selected cancers but may cause low blood counts, nausea, hair loss, infection, or rarely a later blood cancer.

\medterm{Epiretinal Membrane} A thin layer of scar-like tissue growing on the retina, usually over the macula. It may cause blurred or distorted central vision and sometimes needs surgery if visual impairment is significant.

\medterm{Epirubicin} An anthracycline cancer medicine that damages DNA and blocks cell division. Important risks include low blood counts, infection, nausea, hair loss, tissue injury after leakage, and dose-related heart damage.

\medterm{Epirubicin Hydrochloride} The salt form of epirubicin, an anthracycline cancer medicine that damages tumor-cell DNA. It is used for selected cancers but can cause low blood counts, tissue injury after leakage, and dose-related heart damage.

\medterm{Episcleritis} Usually short-lived inflammation of the thin tissue over the white of the eye, causing a localized red patch and mild tenderness or discomfort. Severe pain, light sensitivity, or vision loss suggests deeper eye disease.

\medterm{Episiotomy} A cut made in the tissue between the vagina and anus during vaginal birth to enlarge the opening. Routine use is not recommended; selected use may help urgent or assisted delivery but can cause pain, bleeding, infection, or deeper tearing.

\medterm{Episiotomy Repair} Closing an episiotomy or related tear with stitches after birth. The area is examined to identify deeper injury, and care includes pain relief, hygiene, monitoring for infection, and follow-up of healing and pelvic-floor function.

\medterm{Episodic Ataxia} An inherited condition causing repeated attacks of poor coordination and balance. Episodes may last seconds to hours and can include dizziness, slurred speech, muscle tightening, or abnormal eye movements, with improvement between attacks.

\medterm{Epispadias} A birth defect in which the urethral opening is on the upper surface of the penis or in an unusual position in a female. It may affect urine control and is sometimes associated with bladder-exstrophy conditions.

\medterm{Epistasis} A genetic interaction in which one gene changes or masks the effect of another gene. This can alter inherited traits and complicate how a disease risk or laboratory finding is predicted from genetic testing.

\medterm{Epistaxis} Medical term; see \textit{Nosebleeds}.

\medterm{Posterior Epistaxis} Nasal hemorrhage originating from arteries situated in the posterior nasal cavity or nasopharynx (most commonly branches of the sphenopalatine artery or Woodruff plexus), typically severe, profuse, draining down the oropharynx, and requiring posterior nasal packing or arterial embolization.

\medterm{Epithalamic} Relating to the epithalamus, a small region near the center of the brain that includes the pineal gland and structures involved in sleep timing, hormonal rhythms, and responses to reward or stress.

\medterm{Epithalamus} A dorsal posterior division of the diencephalon comprising the pineal gland (epiphysis cerebri), the habenular nuclei, the stria medullaris thalami, and the posterior commissure, involved in circadian melatonin rhythmicity and limbic-midbrain connections.

\medterm{Epithelial} Relating to cells that cover the skin, line body passages and organs, or form glands. These cells protect surfaces and may absorb substances, release fluids, or help detect changes.

\medterm{Epithelial Cells} Cells that cover the skin, line body passages and organs, and form many glands. They create a protective barrier and may absorb substances, release fluids, move material, or sense changes in their surroundings.

\synonyms
Epithelial Tissue Cells

\medterm{Epithelial Cyst} A usually benign, enclosed sac lined by surface-type cells and filled with fluid, mucus, or a skin protein called keratin. It may enlarge, become painful, or become infected and sometimes requires removal.

\medterm{Epithelial Tissue} Tissue made of closely packed cells that covers the skin, lines internal spaces and tubes, and forms glands. It provides protection and may absorb, secrete, transport, or exchange substances.

\medterm{Epithelioid Cell} A rounded cell that resembles a surface-lining cell. The term often describes an activated immune cell in a granuloma, but it can also describe the appearance of some tumors.

\medterm{Epithelioid Histiocyte} An activated macrophage that has undergone structural transformation under the influence of interferon-gamma in chronic granulomatous inflammation, acquiring abundant pale eosinophilic cytoplasm, elongated slipper-shaped nuclei, and indistinct cell borders resembling epithelial cells.

\medterm{Epithelioid Sarcoma} A rare cancer of connective or soft tissue, often beginning in the skin or deeper tissues of the hands, forearms, feet, or lower limbs. It may recur or spread and requires specialist treatment.

\medterm{Epithelioid Trophoblastic Tumor} A rare pregnancy-related cancer arising from cells that normally help form the placenta. It may appear months or years after pregnancy, cause abnormal bleeding, and requires specialist cancer treatment.

\medterm{Epithelioma} An older or broad term for a tumor arising from epithelial cells; the specific diagnosis should identify whether it is benign, precancerous, or malignant.

\medterm{Epitheliopathy} Disease or injury affecting epithelial cells, which cover body surfaces and line organs. The term is nonspecific and may describe conditions of the cornea, skin, gastrointestinal tract, or other epithelial tissues.

\medterm{Epithelium} A layer or sheet of closely packed cells covering the skin, lining body spaces and tubes, or forming glands. It protects surfaces and may absorb substances, release fluids, transport material, or allow exchange.

\medterm{Epitope} The small, specific part of an antigen that is recognized by an antibody or immune-cell receptor. Recognition can trigger an immune response; epitopes are important in infections, vaccines, allergies, and antibody-based treatments.

\medterm{Epitope Mapping} The process of identifying an epitope, the small part of an antigen recognized by an antibody or specific immune-cell receptor. It helps researchers understand immune responses and develop vaccines, diagnostic tests, and antibody-based treatments.

\medterm{Eplerenone} A medicine that blocks aldosterone, a hormone that promotes salt retention and heart remodeling. It treats high blood pressure and selected heart-failure patients; kidney function and blood potassium require monitoring.

\medterm{Epley Maneuver} A sequence of head and body movements that returns displaced inner-ear particles to their usual position. It can relieve brief spinning triggered by head movement in posterior-canal positional vertigo.

\medterm{Epoetin Alfa} A laboratory-made form of erythropoietin that stimulates the bone marrow to produce red blood cells. It treats selected anemias, but can raise blood pressure and the risk of blood clots if hemoglobin rises too high.

\medterm{Epoetin Beta} A laboratory-made form of erythropoietin used for selected anemias, including some caused by chronic kidney disease. Treatment is individualized and requires monitoring because it may cause high blood pressure or blood clots.

\medterm{Eponym} A medical name derived from a person’s name or another proper name, such as the name of a disease, sign, procedure, or body structure. Eponyms may be replaced by descriptive terms to improve clarity.

\medterm{Epoprostenol} A short-acting medicine that widens blood vessels in the lungs and reduces platelet clumping. Continuous intravenous infusion treats severe pulmonary arterial hypertension; interruption can cause rapid worsening.

\medterm{EPO Test} A blood test measuring erythropoietin, a kidney-made hormone that stimulates red-blood-cell production. It may help investigate anemia or an abnormally high red-cell count alongside blood counts and other tests.

\medterm{Epoxide Hydrolase} An enzyme that changes reactive chemicals called epoxides into less reactive substances. It helps process some medicines, toxins, and natural fats and can influence their effects and the risk of tissue injury.

\medterm{Epoxy Compounds} Chemicals containing a reactive group used in adhesives, coatings, plastics, and other materials. Uncured components can irritate the skin and eyes, cause allergic skin inflammation, or irritate the airways; workplace exposure may cause asthma.

\medterm{Epsom Salts} Magnesium sulfate crystals used in some baths and, occasionally, as a laxative. Drinking them can cause diarrhea or magnesium poisoning, especially in people with kidney disease, so medical advice is important.

\medterm{Epstein-Barr Virus} A common virus spread mainly through saliva. It can cause infectious mononucleosis and remains dormant in the body; it is also linked to some lymphomas, other cancers, and immune disorders.

\medterm{Epstein-Barr Virus Antibody Test} A blood test that measures antibodies produced against Epstein-Barr virus. The pattern of antibodies, symptoms, and timing can help distinguish recent infection from past infection, but results need clinical interpretation.

\medterm{Epstein Pearls} Small, white, harmless cysts containing keratin that appear on a newborn’s gums or roof of the mouth. They usually disappear within weeks without treatment and should not be squeezed.

\medterm{Eptifibatide} An injectable antiplatelet medicine used during selected urgent coronary treatments. It prevents platelets from clumping by blocking a platelet-binding receptor, but can cause serious bleeding and requires blood-count and kidney monitoring.

\medterm{Epulis} A localized growth arising from the gum or tissues around a tooth. Most are reactive and benign, but persistent, painful, bleeding, or enlarging growths need dental examination and sometimes biopsy.

\medterm{Equilibrium Receptors} Sensory cells in the inner ear that detect head movement and position. Signals from them, the eyes, and muscles help the brain maintain balance and coordinate posture and eye movements.

\medterm{Equilin} An estrogen hormone found mainly in pregnant mares and included in some mixtures of conjugated estrogens. Estrogen treatment can increase the risk of blood clots and other adverse effects and requires individualized medical advice.

\medterm{Equine Morbillivirus} A virus related to measles that has caused severe respiratory or neurologic disease in horses and can infect other mammals. Suspected cases require veterinary investigation and infection-control measures.

\medterm{Equinovarus Deformity} A deformity in which the foot points downward and turns inward. It may be present at birth or develop from muscle, nerve, joint, or bone problems, limiting walking and requiring cause-specific treatment.

\medterm{Equipment Safety} Measures that reduce harm from medical devices by using the correct equipment, checking it before use, cleaning it properly, maintaining it, and training users. Faulty equipment should be removed from service.

\medterm{Equivalence Trial} A clinical study designed to determine whether a new treatment and an established treatment have effects within a predefined clinically acceptable range of one another. The range must be set before results are analyzed.

\medterm{Eradication} Complete removal of an infection or disease from a person or population. For an individual infection, eradication may require a full treatment course and a follow-up test to confirm that the organism is gone.

\medterm{Erbb Receptors} A family of cell-surface proteins that receive growth signals and regulate cell division, survival, and specialization. Abnormal activity in some members can promote cancers and may guide targeted treatment.

\medterm{Erbium} A chemical element used in specialized lasers and materials. Erbium lasers can precisely remove or remodel selected skin and dental tissue; treatment may cause pain, burns, pigment changes, or scarring.

\medterm{Erb Palsy} A peripheral obstetric or traumatic traction injury to the upper trunk of the brachial plexus involving spinal nerve roots C5 and C6 (and sometimes C7), causing paralysis of the deltoid, supraspinatus, infraspinatus, and biceps brachii muscles, resulting in an adducted, internally rotated shoulder and extended, pronated forearm ('waiter's tip' posture).
\synonyms
Erb-Duchenne Palsy; Upper Trunk Brachial Plexus Palsy

\medterm{Erb's Palsy} Weakness or paralysis of the shoulder and arm caused by injury to the upper nerves of the brachial plexus, often during difficult birth. The arm may hang inward with the elbow straight and palm facing backward.

\medterm{ERCP} A procedure using a flexible camera and X-ray guidance to diagnose or treat blocked, narrowed, or leaking bile or pancreatic ducts. It can remove stones or place a stent, but may cause pancreatitis, bleeding, infection, or perforation.

\medterm{Erdheim-Chester Disease} A rare non-Langerhans cell histiocytic neoplasm driven by activating kinase mutations in the MAPK pathway (predominantly BRAF V600E), characterized pathologically by lipid-laden foamy histiocyte infiltration causing symmetric diaphyseal osteosclerosis of long bones, periaortic fibrosis ('coated aorta'), and perinephric infiltration ('hairy kidney').
\synonyms
ECD; Non-Langerhans Histiocytosis

\medterm{Erectile Disorder} Alternate name; see \textit{Erectile Dysfunction}.

\medterm{Erectile Dysfunction} Repeated difficulty getting or keeping an erection firm enough for satisfactory sexual activity. It may result from blood-vessel disease, diabetes, nerve or hormonal disorders, medicines, penile conditions, anxiety, depression, or mixed causes and can signal cardiovascular risk.

\medterm{Erection} Temporary enlargement and stiffening of the penis when blood fills erectile tissue and drainage through veins decreases. It can occur with sexual stimulation, during sleep, or spontaneously and normally subsides afterward.

\medterm{Erector Spinae} A group of muscles running along both sides of the spine. They straighten and stabilize the back and help bend the spine sideways during standing, walking, and lifting.

\synonyms
Sacrospinalis Muscle; Spinalis and Longissimus Complex

\medterm{Ergocalciferol} Vitamin D2, a form of vitamin D used to prevent or treat deficiency. It helps the body absorb calcium and maintain bones, but excessive doses can cause high blood calcium, nausea, confusion, or kidney injury.

\medterm{Ergolines} A family of compounds related to chemicals made by ergot fungi. Some affect serotonin, dopamine, or adrenaline receptors and are used or studied for migraine, hormone disorders, or neurologic conditions.

\medterm{Ergometrine} A medicine that contracts the uterus and narrows some blood vessels. It may help control severe bleeding after childbirth, but can dangerously raise blood pressure and is unsuitable for some people with hypertension or heart disease.

\medterm{Ergonomics} Designing work, tools, furniture, and environments to match the body’s abilities and limits. Good ergonomics can reduce strain, awkward posture, repetitive-injury risk, fatigue, and errors.

\medterm{Ergosterol} A fat-like substance that helps maintain fungal cell membranes. Many antifungal medicines damage fungi by binding ergosterol or blocking its production; human cells use cholesterol instead.

\medterm{Ergot} A fungus that can infect grains such as rye and produce toxic chemicals. Eating contaminated grain or taking excessive ergot-derived medicines can cause severe blood-vessel narrowing, neurologic symptoms, or uterine contractions.

\medterm{Ergotamine} A medicine that narrows certain blood vessels and can stop some migraine attacks. It has many interactions and should not be overused because it can cause dangerous loss of blood flow to the limbs or organs.

\medterm{Ergotism} Poisoning caused by ergot toxins or excessive ergot-derived medicines. Severe blood-vessel narrowing can cause burning pain, numbness, cold or discolored limbs, confusion, seizures, tissue death, or dangerous uterine contractions.

\medterm{Ergots} Chemicals produced by ergot fungi, some of which are used in medicines for migraine or bleeding. They can strongly narrow blood vessels and interact with other medicines, causing serious circulation problems or poisoning.

\medterm{Eribulin Mesylate} A cancer medicine that interferes with microtubules, structures needed for cell division. It is used for selected advanced breast cancers and liposarcomas; common risks include low blood counts, nerve damage, fatigue, and nausea.

\medterm{Erlotinib Hydrochloride} A cancer medicine that blocks growth signals from the epidermal growth-factor receptor. It treats selected cancers with suitable molecular changes and may cause rash, diarrhea, liver injury, or lung inflammation.

\medterm{Erosion} Loss of only the surface layer of skin or a moist lining, without the deeper tissue loss of an ulcer. Friction, inflammation, infection, injury, or chemicals may cause it and may produce pain or bleeding.

\medterm{Erosion Lesion} A shallow area where the surface layer of skin or a moist lining has been lost. It may result from rubbing, inflammation, infection, injury, or chemicals and usually heals when the cause is removed.

\medterm{Erosive Esophagitis} Inflammation of the swallowing tube with visible breaks in its lining, usually from repeated stomach-acid reflux. It can cause painful swallowing, heartburn, or bleeding and may lead to narrowing if untreated.

\medterm{Erosive Gastritis} Irritation and inflammation of the stomach lining with shallow areas of damage. Common causes include anti-inflammatory pain medicines, alcohol, severe illness, and stress on the body; it may cause pain, vomiting, or gastrointestinal bleeding.

\medterm{Erosive Lichen Planus} A long-lasting immune-mediated inflammation that causes painful raw areas, burning, bleeding, or soreness in the mouth or genital tissues. Scarring and narrowing may develop, so examination and treatment can help protect function.

\synonyms
Vulvovaginal Gingival Syndrome; Erosive Vulvovaginal Lichen Planus

\medterm{Erratic Results} Test results that change unexpectedly or do not agree with repeated measurements. Possible causes include sample problems, device error, normal biological variation, medicines, or a test that does not fit the clinical question.

\medterm{Ertapenem} A once-daily antibiotic for selected serious infections caused by susceptible bacteria. It does not reliably treat Pseudomonas or enterococcal infections and can cause diarrhea, allergic reactions, seizures, or antibiotic-resistant organisms.

\medterm{Erucic Acid} A fatty acid found in some plant oils, including mustard and older varieties of rapeseed oil. Modern edible oils are usually bred to contain little of it; unusually high exposure may harm the heart.

\medterm{Erucic Acids} Long-chain fatty acids found mainly in certain plant oils. The amount consumed depends on the oil, and unusually high exposure may be harmful; ordinary low-erucic edible oils provide much less.

\medterm{Eructation} Release of swallowed air or gas from the stomach through the mouth, commonly called belching. Frequent troublesome belching may be related to swallowed air, reflux, food intolerance, or other digestive problems.

\synonyms
Belching; Burping

\medterm{Eruption Cyst} A soft, fluid-filled swelling over a tooth that is about to come through the gum, often blue or purple. It usually disappears when the tooth erupts, but painful or infected swelling needs dental review.

\medterm{Eruptive Fever} Fever occurring with a widespread rash. Viral infections are common causes, but bacteria, medicines, allergic reactions, and immune disorders can also be responsible; urgent assessment is needed for severe illness or breathing difficulty.

\medterm{Eruptive Xanthomas} Sudden crops of small yellow-red bumps, often on the buttocks, elbows, knees, or trunk, caused by very high blood triglycerides. They signal a serious lipid disorder requiring prompt evaluation and treatment.

\medterm{Eruptive Xanthomatosis} A sudden outbreak of many small yellow-red skin bumps caused by fat collecting in skin immune cells. It is usually linked to very high triglycerides, sometimes from uncontrolled diabetes or an inherited lipid disorder.

\medterm{Eryptosis} A controlled process by which aging or damaged red blood cells shrink and are marked for removal from the bloodstream. It helps maintain blood-cell balance but may increase in some diseases or during oxidative stress.

\medterm{Erysipelas} A sudden bacterial infection of the upper skin and lymph vessels, usually causing a painful, raised, sharply bordered red patch with fever. Prompt antibiotic treatment is important to prevent spread or complications.

\medterm{Erysipeloid} A skin infection caused by bacteria acquired through cuts while handling fish, meat, poultry, or animals. It usually causes a slowly expanding painful purple-red patch on a finger or hand and needs antibiotic treatment.

\medterm{Erysipelothrix} A group of bacteria found in animals and animal products that can enter broken skin and cause a painful localized infection. Rarely, infection spreads to the bloodstream or heart valves.

\medterm{Erythema} Redness of the skin or a moist body lining caused by increased blood flow near the surface. It may result from inflammation, infection, heat, allergy, or another illness and often fades briefly when pressed.

\medterm{Erythema Ab Igne} A net-like red-brown skin discoloration caused by repeated exposure to moderate heat, such as heating pads, hot-water bottles, or heaters. Stopping the exposure may allow improvement; persistent lesions need examination.

\medterm{Erythema Annulare Centrifugum} A reactive gyrate erythema presenting with slowly expanding, polycyclic, annular erythematous plaques characterized by central clearing and a distinctive trailing rim of fine scale along the inner border of the advancing erythematous margin.
\synonyms
EAC; Centrifugal Annular Erythema

\medterm{Erythema Chronicum Migrans} An expanding skin rash that can appear days to weeks after a tick transmits Lyme disease bacteria. It may be evenly red or ring-shaped and does not always look like a classic bull’s-eye.

\medterm{Erythema Dose} The amount of ultraviolet or other radiation needed to produce visible skin redness. The minimal erythema dose is the smallest amount that causes a defined response and varies with skin type, medicines, and exposure conditions.

\medterm{Erythema Elevatum Diutinum} A rare chronic cutaneous leukocytoclastic vasculitis presenting as persistent, symmetrical, red-violaceous to yellowish-brown papules, nodules, and plaques over the extensor surfaces of joints, associated with IgA paraproteinemia, hematologic disorders, and chronic infections.
\synonyms
EED; Chronic Extensor Leukocytoclastic Vasculitis

\medterm{Erythema Gyratum Repens} A rapidly spreading, ringed skin rash with a wood-grain pattern. It is often associated with an internal cancer, especially of the lung, so new or persistent cases require prompt medical evaluation.

\medterm{Erythema Induratum} Recurrent tender lumps caused by inflammation of the fat beneath the skin, usually on the backs of the lower legs. Tuberculosis and other infections, inflammatory diseases, or medicines may be responsible.

\medterm{Erythema Infectiosum} A common childhood infection caused by parvovirus B19, also called fifth disease. Mild fever, tiredness, or cold symptoms may come first, followed by bright red cheeks and a lacy rash on the body and limbs.

\synonyms
Fifth Disease; Slapped Cheek Disease

\medterm{Erythema Marginatum} A temporary, non-itchy, ring-shaped or wavy red rash that usually appears on the trunk and upper limbs. It is a recognized sign of acute rheumatic fever and may come and go over days or weeks.

\medterm{Erythema Migrans} An expanding rash that develops days to weeks after a tick transmits Lyme disease bacteria. It may be uniformly red or ring-shaped and is usually diagnosed from its appearance and exposure history.

\medterm{Erythema Multiforme} A short-lived immune reaction, often triggered by herpes simplex infection or, less commonly, another infection or medicine. It causes symmetrical target-shaped spots, usually on the hands, feet, arms, or face; mouth involvement may occur.

\medterm{Erythema Nodosum} Inflammation of the fat beneath the skin causing tender red or purple lumps, usually on both shins. Infections, medicines, inflammatory bowel disease, sarcoidosis, and other conditions may trigger it.

\medterm{Erythema Toxicum} A common harmless newborn rash with red spots, small bumps, or tiny pus-filled areas. It often appears during the first days of life and disappears on its own without treatment.

\medterm{Erythema Toxicum Neonatorum} A harmless newborn rash that usually appears during the first few days as red blotches, bumps, or tiny pustules on the face, trunk, or limbs. It usually spares the palms and soles and clears within one to two weeks.

\synonyms
Erythema Neonatorum; Toxic Erythema of the Newborn

\medterm{Erythrasma} A superficial bacterial infection in warm, moist skin folds that causes well-defined reddish-brown patches, often in the groin, armpits, or between toes. It usually responds to prescribed antibacterial treatment and keeping folds dry.

\medterm{Erythroblast} An immature red-blood-cell precursor made in the bone marrow. It develops through several stages before becoming a mature red blood cell; abnormal numbers may indicate altered blood-cell production.

\medterm{Erythroblastosis Fetalis} A condition in which antibodies from the pregnant person destroy a fetus’s or newborn’s red blood cells. It can cause anemia, jaundice, widespread swelling, heart failure, or death and is often prevented by blood-group screening and preventive treatment.

\medterm{Erythroblastosis Foetalis} Destruction of fetal or newborn red blood cells by antibodies from the mother. It can cause anemia, jaundice, widespread fluid buildup, heart failure, or death. Pregnancy blood-group testing and preventive treatment can reduce the risk.

\medterm{Erythrocyte} A mature red blood cell that carries oxygen from the lungs to tissues and helps return carbon dioxide to the lungs. Its flexible shape allows it to pass through very small blood vessels.

\medterm{Erythrocyte Casts} Tube-shaped groups of red blood cells seen in urine. They form inside the kidney's small tubes and strongly suggest bleeding or inflammation in the kidney's filtering units, such as glomerulonephritis or vasculitis.
\synonyms
Red Blood Cell Casts; RBC Casts

\medterm{Erythrocyte Enzyme Deficiency} An inherited or acquired lack of a protein that red blood cells need for energy or protection. The cells may break down too early and cause anemia; examples include G6PD and pyruvate-kinase deficiency.

\medterm{Erythrocyte Membrane} The flexible outer layer of a red blood cell. It maintains the cell’s shape, allows it to bend through small vessels, and controls movement of water and substances into and out of the cell.

\medterm{Erythrocyte Sedimentation Rate} A blood test measuring how quickly red blood cells settle in a tube over one hour. A faster rate can occur with inflammation, but it does not identify the cause and is affected by anemia, age, pregnancy, and other conditions.

\medterm{Erythrocytosis} An abnormally high number or volume of red blood cells, reflected by high hemoglobin or hematocrit. It may result from dehydration, low oxygen, excess erythropoietin, or a bone-marrow disorder and can increase clot risk.

\medterm{Erythroderma} Widespread redness and scaling affecting most of the skin. Causes include psoriasis, eczema, medicines, infection, or lymphoma; severe cases can cause fluid loss, temperature problems, infection, and heart strain.

\medterm{Erythrodermic Psoriasis} A rare, potentially life-threatening form of psoriasis causing nearly all the skin to become intensely red, painful, and scaly. Fever, chills, dehydration, infection, and unstable blood pressure may occur and require urgent care.

\medterm{Erythrodysplasia} Abnormal development or maturation of red-blood-cell precursors in the bone marrow. It can produce ineffective blood formation and anemia and may occur with a marrow disorder or other illness.

\medterm{Erythroid Hyperplasia} An increase in immature red-blood-cell precursors in the bone marrow, usually because the body is responding to anemia, blood loss, red-cell destruction, or low oxygen. The underlying cause needs evaluation.

\medterm{Erythroleukoplakia} A mouth lesion containing both red and white areas. It has a substantial risk of advanced precancerous change or cancer, so a dentist or specialist should examine it and usually perform a biopsy.

\medterm{Erythromelalgia} Repeated attacks of burning pain, warmth, and redness in the feet, hands, or sometimes face, often triggered by heat or exercise and eased by cooling. It may be inherited or linked to blood or nerve disorders.

\medterm{Erythromycin} An antibiotic that blocks bacterial protein production. It treats selected infections and can be used when some other antibiotics are unsuitable; nausea, diarrhea, liver injury, and abnormal heart rhythm are possible.

\medterm{Erythromycin Estolate} An oral form of the antibiotic erythromycin. It blocks bacterial protein production and treats selected infections; stomach upset, liver inflammation, allergic reactions, and medicine interactions can occur.

\medterm{Erythromycin Ophthalmic Prophylaxis} Antibiotic ointment placed in a newborn’s eyes soon after birth to reduce the risk of severe eye infection from gonorrhea acquired during delivery. It does not replace testing and treatment of infections during pregnancy.

\medterm{Erythroplakia} A persistent, sharply defined red patch on a moist lining, especially inside the mouth. It may contain advanced precancerous change or cancer and requires prompt examination, usually with biopsy.

\medterm{Erythroplasia} A sharply bordered red patch on skin or a moist lining. Persistent lesions, especially in the mouth or genital area, may represent a serious precancerous change or cancer and need medical assessment.

\medterm{Erythroplasia Of Queyrat} Squamous-cell carcinoma in situ affecting the head of the penis or foreskin. It usually appears as a persistent, well-defined red or velvety patch and can progress to invasive cancer without treatment.

\medterm{Erythropoiesis} The process of making mature red blood cells, mainly in the bone marrow. It is stimulated by erythropoietin when oxygen is low and requires healthy marrow plus adequate iron, vitamin B12, and folate.

\medterm{Erythropoiesis-Stimulating Agents} Medicines that signal the bone marrow to make more red blood cells. They are used for selected anemias, including some caused by kidney disease or cancer treatment, but may increase blood pressure and clot risk.

\medterm{Erythropoietin} A hormone made mainly by the kidneys when body tissues receive too little oxygen. It signals the bone marrow to produce red blood cells; kidney disease can reduce this signal and contribute to anemia.

\medterm{Erythropoietin Deficiency} An insufficient amount of erythropoietin for the body's need, resulting in reduced stimulation of red-blood-cell production in the bone marrow. It commonly contributes to anemia in chronic kidney disease; reduced response to available erythropoietin is a separate phenomenon.

\medterm{Erythropoietin Receptor} A protein on immature red-blood-cell precursors that receives the erythropoietin signal. Activation helps these cells survive, multiply, and mature into red blood cells.

\medterm{Erythropoietin Test} A blood test measuring the kidney-made hormone that stimulates red-blood-cell production. It can help investigate selected anemias or high red-cell counts and is interpreted with the blood count, oxygen status, and kidney function.

\medterm{Erythrosine} A red dye used in some foods, medicines, and laboratory procedures. Exposure is usually small, but regulations vary and people with a known sensitivity should follow product and medical advice.

\medterm{ESBL} An enzyme made by some bacteria that destroys many penicillin and cephalosporin antibiotics. Infections caused by these bacteria can be harder to treat and require laboratory testing to select an effective medicine.

\medterm{ESBL-Producing Organisms} Bacteria that make enzymes capable of breaking down many penicillin and cephalosporin antibiotics. They commonly include certain intestinal bacteria such as E. coli and Klebsiella; treatment depends on susceptibility testing and infection site.

\medterm{Escape Reaction} A body response that restores an organ’s activity after it has been reduced by a medicine, hormone, nerve signal, or feedback system. The response may lessen the effect of continued treatment or stimulation.

\medterm{Eschar} A firm layer of dead tissue, often black or brown, over a burn, pressure injury, wound, infection, or area with poor blood supply. It can hide deeper damage and may need specialist assessment or removal.

\medterm{Escharotomy} A surgical cut through stiff dead tissue around a deep burn to release pressure. It may restore circulation to a limb or allow the chest to expand, and is performed urgently when swelling threatens function.

\medterm{Escherichia} A group of bacteria that includes E. coli and related species found in intestines or the environment. Some strains are harmless, while others cause diarrhea, urinary infection, wound infection, or bloodstream infection.

\medterm{Escherichia Albertii} An uncommon intestinal bacterium that can cause diarrheal illness. It may be mistaken for certain E. coli strains, so specialized laboratory identification may be needed when infection is suspected.

\medterm{Escherichia Blattae} A rarely reported Escherichia species associated mainly with insects and the environment. Its significance in a human sample is uncertain and must be judged from symptoms, the specimen source, and other test results.

\medterm{Escherichia Coli} A bacterium that normally lives in the intestine. Some strains are harmless, while others cause urinary, intestinal, wound, or bloodstream infections, including severe bloody diarrhea and hemolytic uremic syndrome.

\medterm{Escherichia Coli Infection} Infection caused by a harmful strain of E. coli. It may affect the intestine, urinary tract, wounds, or bloodstream; symptoms, contagiousness, and treatment depend on the strain and site.

\medterm{Escherichia Coli O157:H7 Infection} Infection with a toxin-producing E. coli strain, usually acquired from contaminated food, water, animals, or people. It commonly causes severe cramps and bloody diarrhea and can lead to kidney failure; urgent medical advice is important.

\medterm{Escherichia Fergusonii} An uncommon bacterium that can occasionally cause urinary, wound, intestinal, or bloodstream infection, especially in people with serious illness or weakened immunity. Laboratory identification and susceptibility testing guide treatment.

\medterm{Escherichia Hermannii} An uncommon Escherichia species that rarely causes opportunistic infection, including urinary or bloodstream infection. Its importance depends on the patient’s symptoms, immune status, and the specimen in which it is found.

\medterm{Escherichia Vulneris} An uncommon bacterium rarely linked to wound, urinary, or bloodstream infection, usually in people with other serious illnesses. A laboratory result requires clinical correlation to distinguish infection from contamination.

\medterm{Escitalopram} An antidepressant that increases available serotonin in the brain. It treats depression and generalized anxiety disorder; nausea, sleep changes, sexual problems, withdrawal symptoms, bleeding risk, and rarely serotonin toxicity may occur.

\medterm{ESKD} Severe permanent loss of kidney function, usually meaning kidney filtration below about 15 milliliters per minute or a need for dialysis or transplantation. It can cause fluid, electrolyte, anemia, and toxin problems and needs kidney-specialist care.

\medterm{Esomeprazole} A medicine that reduces stomach-acid production by blocking the acid pump in stomach cells. It treats reflux, ulcers, and acid-related disorders; prolonged or inappropriate use may cause infections, low magnesium, or vitamin deficiencies.

\medterm{Esomeprazole Magnesium} A magnesium salt of esomeprazole that reduces stomach-acid production by blocking the stomach’s acid pump. It treats reflux and ulcers; long-term use may cause diarrhea, infections, low magnesium, or nutrient deficiencies.

\medterm{Esophageal} Relating to the esophagus, the muscular tube that carries swallowed food and liquid from the throat to the stomach.

\medterm{Esophageal Achalasia} A disorder in which the lower esophageal valve does not relax properly and the esophagus cannot push food normally into the stomach. It causes difficulty swallowing, regurgitation, chest discomfort, and weight loss.

\medterm{Esophageal Adenocarcinoma} A cancer arising from gland-forming cells, usually in the lower esophagus and often linked to Barrett esophagus. Reflux, obesity, and smoking increase risk; progressive swallowing difficulty or weight loss requires evaluation and biopsy.

\medterm{Esophageal Atresia} A birth defect in which the esophagus ends in a blind pouch instead of connecting normally to the stomach. A connection to the windpipe may also occur, causing drooling, choking, coughing, and breathing problems during feeds.

\medterm{Esophageal Bleeding} Bleeding from the esophagus, which may result from enlarged veins, inflammation, ulcers, tears, injury, or cancer. Vomiting blood, passing black stools, dizziness, or fainting requires emergency medical care.

\medterm{Esophageal Cancer} A malignant tumor of the swallowing tube. The main types are squamous-cell cancer and adenocarcinoma; progressive difficulty swallowing, weight loss, pain, or bleeding may occur, and treatment depends on the stage.

\synonyms
Esophageal Carcinoma; Esophagus Malignancy

\medterm{Esophageal Candidiasis} A yeast infection of the esophagus, more common with weakened immunity, diabetes, or certain medicines. It often causes painful or difficult swallowing and is treated with antifungal medicine, sometimes after endoscopy confirms it.

\medterm{Esophageal Carcinoma} Cancer arising in the esophagus, most often squamous-cell carcinoma or adenocarcinoma. Progressive difficulty swallowing, weight loss, chest pain, or bleeding may occur; endoscopic biopsy and staging guide treatment.

\medterm{Esophageal Culture} Laboratory growth testing on a sample from the esophagus to identify bacteria or fungi when infection is suspected. Results depend on how the sample was collected and must be interpreted with symptoms and endoscopic findings.

\medterm{Esophageal Dilation} A procedure that gently stretches a narrowed area of the esophagus with a tube or balloon passed through an endoscope. It can improve swallowing caused by scarring, inflammation, injury, or other narrowing; tearing and bleeding are uncommon but serious risks.

\medterm{Esophageal Diseases} Disorders affecting the swallowing tube between the throat and stomach. They include reflux, inflammation, narrowing, movement problems, cancer, infection, and birth defects, often causing swallowing difficulty, pain, heartburn, or regurgitation.

\medterm{Esophageal Diverticulum} A pouch that bulges outward from the wall of the esophagus. It may cause difficulty swallowing, food coming back up, coughing, bad breath, or no symptoms; treatment depends on its size, location, and effects.

\medterm{Esophageal Duplication} A rare birth defect consisting of a cyst or extra segment beside the esophagus. It may press on the esophagus or airway, causing swallowing difficulty, breathing problems, infection, or no symptoms.

\medterm{Esophageal Fistula} An abnormal connection between the esophagus and another organ, often the airway. It can allow food or liquid into the lungs, causing coughing during swallowing, repeated pneumonia, poor nutrition, or severe illness.

\medterm{Esophageal Foreign Body} Food or another object stuck in the esophagus, causing pain, difficulty swallowing, drooling, or breathing problems. A battery, sharp object, complete blockage, or suspected tear requires emergency removal.

\medterm{Esophageal Infection} Infection of the esophagus caused by a yeast, virus, or bacterium, usually in people with weakened immunity or damaged lining. It may cause painful swallowing, difficulty swallowing, chest pain, or ulcers.

\medterm{Esophageal Intramural Pseudodiverticulosis} A rare condition in which small drainage ducts of glands within the esophageal wall become widened. It may cause repeated swallowing difficulty, food impaction, or narrowing and is diagnosed with imaging or endoscopy.

\medterm{Esophageal Leukoplakia} A persistent white patch on the esophageal lining caused by thickened surface cells. It is uncommon but may contain precancerous change, so evaluation and biopsy are often needed to determine its significance.

\medterm{Esophageal Manometry} A test that measures pressure and muscle coordination in the esophagus during swallowing. A thin pressure-sensitive tube passes through the nose into the stomach and helps diagnose achalasia, spasm, and other movement disorders.

\medterm{Esophageal Motility Disorder} A disorder of the muscle contractions or nerve control that moves food through the esophagus. It may cause difficulty swallowing, chest pain, regurgitation, or food sticking; testing may include esophageal manometry.

\medterm{Esophageal Necrosis} Death of esophageal lining, sometimes producing a black appearance. It can result from severe low blood flow, shock, acid injury, or caustic substances and may lead to bleeding or perforation.

\medterm{Esophageal Perforation} A hole through the full thickness of the esophageal wall, caused by forceful vomiting, an instrument, injury, or disease. It can release contents into the chest, causing severe infection and requiring emergency treatment.

\medterm{Esophageal pH Monitoring} A test that records how much acid enters the esophagus, usually for 24 hours, using a thin tube or a small wireless sensor. It helps determine whether acid reflux explains persistent symptoms when the diagnosis is uncertain.

\medterm{Esophageal Spasm} A movement disorder in which esophageal muscles contract too strongly or without normal coordination. It can cause intermittent chest pain and difficulty swallowing; pressure testing helps distinguish it from other disorders.

\medterm{Esophageal Spasms} Uncoordinated or forceful contractions of the esophagus that may cause sudden chest pain, food sticking, or difficulty swallowing. Symptoms can resemble heart pain and should be assessed when severe or new.

\medterm{Esophageal Squamous Cell Carcinoma} A cancer arising from the flat cells lining the esophagus, often in its upper or middle part. Smoking, heavy alcohol use, certain infections, and nutritional factors increase risk; diagnosis requires endoscopic biopsy.

\medterm{Esophageal Stenosis} Narrowing of the passage through the esophagus caused by scarring, inflammation, a tumor, radiation, or a birth defect. It commonly causes progressively worsening difficulty swallowing, especially solid foods.

\medterm{Esophageal Stent} A tube placed inside a narrowed or leaking esophagus to keep it open and allow swallowing. It is often used for cancer or injury; movement, blockage, bleeding, pain, or perforation can occur.

\medterm{Esophageal Stricture} A narrowed area of the esophagus, commonly caused by acid reflux, caustic injury, inflammation, radiation, or procedures. It usually causes gradually worsening difficulty swallowing solids and may require stretching during endoscopy.

\medterm{Esophageal Surgery} An operation on the esophagus or the area where it meets the stomach. It may treat cancer, severe narrowing, perforation, reflux, or movement disorders; the procedure and risks depend on the disease and overall health.

\medterm{Esophageal Ulceration} An open sore in the esophageal lining caused by reflux, infection, medicines, caustic injury, inflammation, or cancer. It may cause painful swallowing, chest pain, or bleeding and often requires endoscopic evaluation.

\medterm{Esophageal Varices} Enlarged veins in the esophageal wall, usually caused by high pressure in the portal vein from advanced liver disease. They may rupture without warning and cause life-threatening vomiting of blood.

\synonyms
Esophageal Dilated Veins; Portosystemic Varices

\medterm{Esophageal Varices Cross-Reference} Alternate name; see \textit{Esophageal varices}.

\medterm{Esophageal Web} A thin shelf of tissue projecting into the upper esophagus. It may cause difficulty swallowing solid food and can be associated with iron-deficiency anemia; persistent symptoms may require endoscopy.

\medterm{Esophagectomy} Surgery to remove part or all of the esophagus, usually for cancer or severe irreversible damage. The remaining digestive tract is reconstructed, commonly using the stomach or a segment of intestine.

\medterm{Esophagitis} Inflammation or injury of the esophageal lining. Common causes include acid reflux, infection, medicines that irritate the esophagus, eosinophilic or other immune-mediated disease, swallowed irritants, and radiation. Symptoms may include heartburn, painful or difficult swallowing, chest pain, food sticking, or bleeding; persistent inflammation can cause ulcers or narrowing.

\medterm{Esophagogastric Bypass} Surgery that creates a new route between the esophagus and stomach to restore passage of food, usually after esophageal removal or severe blockage. The method depends on anatomy, disease, and surgical goals.

\medterm{Esophagogastric Junction} The area where the esophagus joins the stomach. Its muscular valve and surrounding structures help food enter the stomach while limiting backward flow of stomach contents.

\medterm{Esophagogastric Tamponade} Temporary control of severe bleeding from esophageal or stomach veins by inflating a balloon device to apply pressure. It is a rescue measure used while definitive treatment is arranged and can obstruct breathing or cause aspiration.

\medterm{Esophagogastroduodenoscopy} An upper endoscopic diagnostic and therapeutic procedure in which a thin, flexible video endoscope is passed through the mouth to examine the lining of the esophagus, stomach, and duodenum.

\synonyms
EGD; Upper Endoscopy; Upper Gastrointestinal Endoscopy

\medterm{Esophagoscope} A flexible or rigid instrument used to look inside the esophagus. It may allow a clinician to take tissue samples, remove a foreign object, treat bleeding, or perform another procedure.

\medterm{Esophagoscopy} Examination of the esophagus with a camera-equipped instrument. It may investigate swallowing difficulty, bleeding, reflux damage, narrowing, foreign objects, infection, or suspected cancer and may allow treatment or biopsy.

\medterm{Esophagram} An X-ray examination performed while a person swallows contrast liquid. It shows the shape and movement of the esophagus and can reveal narrowing, leaks, pouches, blockage, or abnormal swallowing.

\medterm{Esophagus} A muscular tube that carries food and liquid from the throat to the stomach. Coordinated muscle contractions move swallowed material downward, while valves at each end help limit reflux and air entry.

\medterm{Esophagus Disease} A disorder of the esophagus causing swallowing difficulty, pain, reflux, regurgitation, bleeding, obstruction, inflammation, motility disturbance, or cancer.

\medterm{Esophagus Neoplasm} A benign, precancerous, or malignant growth in the esophagus; evaluation may require endoscopy, biopsy, imaging, and staging.

\medterm{Esotropia} A condition in which one eye turns inward toward the nose, causing the eyes to point in different directions. It may begin in infancy or result from focusing problems, illness, or injury; glasses, patching, or surgery may help depending on the cause.

\medterm{ESR} The erythrocyte sedimentation rate, a nonspecific blood test measuring how quickly red cells settle in a tube; it can reflect inflammation but cannot identify its cause.

\synonyms
SED Rate

\medterm{ESRD} Alternate index label; see \textit{End-Stage Renal Disease}.

\medterm{Essential} Necessary or fundamental; in a diagnosis such as essential hypertension, it indicates that no single secondary cause has been identified.

\medterm{Essential Fats} Fatty acids the body cannot make in sufficient amounts and must obtain from food. They support cell membranes, skin, brain function, and production of signaling molecules involved in inflammation.
\synonyms
Essential Fatty Acids

\medterm{Essential Fatty Acid} A fatty acid the body cannot make in sufficient amounts and therefore must obtain from food; the main human examples are linoleic acid and alpha-linolenic acid.

\medterm{Essential Hypertension} Long-term high blood pressure with no single identifiable underlying cause. Many genetic, kidney, blood-vessel, metabolic, and lifestyle factors contribute, and untreated disease raises the risk of stroke, heart disease, kidney damage, and vision problems.

\medterm{Essential Medicines} Medicines considered necessary to meet the most important health needs of a population. They are selected for proven benefit, safety, affordability, and practical availability and should be reliably accessible through healthcare services.

\medterm{Essential Mixed Cryoglobulinemia} An immune disorder in which abnormal proteins in the blood thicken or clump in the cold and inflame small blood vessels. It may cause purple spots, joint pain, weakness, nerve problems, or kidney disease.

\medterm{Essential Nutrient} A nutrient the body cannot make in sufficient amounts and must obtain from food or supplements. Essential nutrients include certain amino acids, fatty acids, vitamins, minerals, and water.

\medterm{Essential Thrombocythemia} A long-term bone-marrow disorder that produces too many platelets. It can cause abnormal clotting or bleeding and may cause headaches, vision changes, burning pain in fingers or toes, or no symptoms.

\medterm{Essential Tremor} A nervous-system disorder causing rhythmic shaking, usually of both hands during movement or when holding a position. It may also affect the head or voice, gradually worsen, and interfere with writing, eating, or other tasks.

\medterm{Essiac} A herbal mixture promoted as a treatment for cancer and other illnesses. Reliable evidence has not shown that it cures cancer, and its ingredients may interact with medicines or cause liver, kidney, or other harm.

\medterm{Estazolam} A sedative medicine sometimes used short-term for insomnia. It slows brain activity but can cause daytime drowsiness, falls, memory problems, slowed breathing, dependence, and withdrawal; alcohol and opioids increase its dangers.

\medterm{Ester} A chemical formed when an acid combines with an alcohol. Esters occur naturally in fats, medicines, and flavorings; their health effects depend on the specific compound, amount, route of exposure, and duration.

\medterm{Esterase} An enzyme that breaks ester chemical bonds. Esterases help process fats, medicines, and other substances; abnormal activity or inhibition can alter drug effects and may contribute to poisoning.

\medterm{Esterified Estrogens} A mixture of estrogen hormones used for selected menopausal symptoms or estrogen deficiency. Treatment can cause bleeding and increase risks of blood clots, stroke, and breast or uterine cancer, so benefits and risks must be individualized.

\medterm{Esthesia} The ability to sense touch, pain, temperature, movement, or body position. Reduced, increased, or altered sensation may result from nerve injury, brain disease, medicines, or changes in consciousness.

\medterm{Esthesiometer} An instrument used to measure sensitivity to touch, pressure, or other sensations. Eye clinicians may use one to assess corneal sensation when nerve injury or an eye-surface problem is suspected.

\medterm{Esthesioneuroblastoma} A rare cancer arising from nerve-related cells high in the nasal cavity. It may cause blocked or bleeding nose, reduced smell, facial pain, or vision symptoms and is diagnosed with imaging and biopsy.

\medterm{Estimated Average Glucose} An estimate of a person’s average blood-glucose level calculated from the hemoglobin A1c result. It is reported in familiar glucose units but does not show daily highs, lows, or sudden changes.

\medterm{Estimated Date Of Confinement} The approximate date when a pregnancy is expected to reach delivery, calculated from the last menstrual period or an early ultrasound. Only a small proportion of births occur on the estimated date.

\medterm{Estimated Due Date} The approximate date of delivery, usually calculated as 40 weeks from the first day of the last menstrual period and adjusted when needed using an early ultrasound. Most births occur before or after this date.

\medterm{Estimated Fetal Weight} An ultrasound-based estimate of a fetus’s weight using measurements of the head, abdomen, and thigh bone. It helps assess growth and plan care, but commonly differs from actual birth weight by about 10–15 percent.

\medterm{Estimated Glomerular Filtration Rate} An estimate of how well the kidneys filter blood, calculated mainly from blood creatinine and other factors. It helps assess kidney disease and adjust medicines but is less reliable during rapidly changing illness or unusual body size.

\medterm{Estradiol} The main active estrogen hormone during most reproductive years. It helps regulate ovulation and menstruation and supports the vagina, uterus, bones, brain, and other tissues; levels fall after menopause.

\synonyms
17$\beta$-Estradiol

\medterm{Estradiol Blood Test} A blood test measuring estradiol. It may help evaluate puberty, menstrual or fertility problems, ovarian or testicular function, menopause, or selected hormone disorders; results vary with age, sex, cycle timing, pregnancy, and medicines.

\medterm{Estramustine} A cancer medicine combining an estrogen-like component with a compound that interferes with cell division. It has been used for selected advanced prostate cancers; blood clots, fluid retention, nausea, and heart effects may occur.

\medterm{Estriol} An estrogen produced mainly by the placenta during pregnancy from substances made by the fetus and mother. Maternal blood estriol may be included in prenatal screening, but an abnormal result does not by itself diagnose a problem.

\medterm{Estrogen} A group of hormones involved in reproductive development, menstruation, fertility, bone strength, and many other body functions. The main forms and amounts vary with age, sex, pregnancy, and ovarian function.

\medterm{Estrogen Overdose} Taking or receiving more estrogen than prescribed. It may cause nausea, vomiting, breast tenderness, headache, or abnormal bleeding and can increase the risk of blood clots; urgent advice is needed after a large exposure.

\medterm{Estrogen Receptor} A cell protein that binds estrogen and changes cell activity. Testing for estrogen receptors helps classify some cancers, especially breast cancer, and can guide hormone-based treatment.

\medterm{Estrogen Replacement Therapy} Treatment supplying estrogen to relieve selected menopausal symptoms or replace deficient hormone. If the uterus is present, a progestogen is usually added to protect its lining; risks and benefits must be reviewed individually.

\medterm{Estrogen-Replacement Therapy} Treatment that supplies estrogen when natural levels are low, such as after menopause or ovarian failure. It may relieve hot flushes and protect bone, but can cause bleeding and increase some health risks.

\synonyms
Estrogen Therapy

\medterm{Estrogens} Hormones, including estradiol, that regulate reproductive development, menstruation, fertility, bones, and other tissues. Before menopause they are produced mainly by the ovaries; after menopause, levels are lower and come from several tissues.

\medterm{Estrogen Test} A blood, urine, or saliva test measuring one or more estrogen hormones. It may help evaluate puberty, menstrual problems, fertility, menopause, pregnancy, or hormone disorders; interpretation depends on the test type and timing.

\synonyms
Estrogen Levels Test
Estrogen Blood Test

\medterm{Estrone} A naturally occurring estrogen made partly by converting other hormones in body tissues. It is the main estrogen present after menopause, although estradiol is usually more active before menopause.

\medterm{Estrous Cycle} The recurring reproductive cycle of many female mammals other than humans. It includes hormonal and physical changes that prepare for ovulation, mating, and possible pregnancy.

\medterm{Eszopiclone} A sleep medicine that enhances inhibitory signaling in the brain. It can help insomnia but may cause drowsiness, impaired coordination, memory problems, dependence, unusual sleep behaviors, or a bitter taste.

\medterm{Etacrynic Acid} A diuretic that makes the kidneys remove more salt and water, reducing swelling and sometimes blood pressure. It can cause dehydration, abnormal electrolytes, kidney problems, or hearing damage, especially at high doses.

\medterm{Etanercept} A biologic medicine that blocks tumor-necrosis factor, an inflammation signal. It treats selected arthritis, psoriasis, and other inflammatory diseases but can increase serious infection risk, so tuberculosis and hepatitis screening may be needed.

\medterm{Ethacrynic Acid} A loop diuretic that increases kidney excretion of salt and water. It can reduce fluid overload but may cause dehydration, low electrolytes, kidney injury, or dose-related hearing damage.

\medterm{Ethambutol} An antibiotic used with other medicines to treat tuberculosis and selected infections caused by related bacteria. It can damage the optic nerve, causing blurred vision or trouble distinguishing colors, so vision testing is important.

\medterm{Ethambutol Toxicity} A dose- and duration-dependent adverse pharmacological effect of the first-line antimycobacterial agent ethambutol, characterized by retrobulbar or axial optic neuritis, presenting with bilateral painless reduction in visual acuity, central scotomas, and loss of red-green color discrimination.

\medterm{Ethane} A colorless, highly flammable gas used mainly to make industrial chemicals. It is not usually poisonous through direct action, but high concentrations can displace oxygen and cause suffocation or fire hazards.

\medterm{Ethanol} The alcohol in alcoholic drinks and some disinfectants, medicines, and fuels. Drinking too much can impair judgment, suppress breathing, lower blood sugar, cause dependence, and damage the liver and other organs.

\medterm{Ethanolamine} A chemical used in some industrial products, medicines, and cosmetic ingredients. Concentrated exposure can irritate or burn the skin, eyes, and airways; health effects depend on the product and amount of exposure.

\medterm{Ethanol Metabolism} The body’s breakdown of alcohol, mainly in the liver, into acetaldehyde and then acetate. The process affects blood-alcohol levels, drug interactions, intoxication, hangovers, and long-term alcohol-related organ damage.

\medterm{Ethanol Poisoning} Harm from drinking a dangerous amount of alcohol, causing confusion, poor coordination, vomiting, slow or irregular breathing, low blood sugar, seizures, coma, or death. An unconscious or poorly breathing person needs emergency care.

\medterm{Ether} A family of chemicals containing an oxygen atom between carbon groups. Some are volatile solvents; diethyl ether was once an anesthetic but is highly flammable, and inhalation can cause irritation, dizziness, or unconsciousness.

\medterm{Ethers} A group of chemicals containing an oxygen atom between two carbon groups. Some are used as solvents or fuels; exposure can cause irritation, dizziness, unconsciousness, or fire and explosion hazards.

\medterm{Ethics} Principles used to decide what is right, fair, respectful, and responsible in healthcare and research. Medical ethics includes informed choice, privacy, avoiding harm, helping patients, and treating people fairly.

\medterm{Ethinyl Estradiol} A synthetic estrogen used in many hormonal contraceptives and some other treatments. It can cause nausea or irregular bleeding and increases blood-clot risk, especially with smoking, migraine with aura, or certain medical conditions.

\medterm{Ethiodized Oil} An iodine-containing oil used as a contrast material in selected imaging or treatment procedures, including some lymphatic and liver procedures. It can affect thyroid function or rarely enter blood vessels and cause an embolism.

\medterm{Ethion} An organophosphate pesticide that can poison the nervous system by blocking acetylcholinesterase. Exposure may cause sweating, vomiting, pinpoint pupils, muscle twitching, weakness, breathing failure, or death and requires urgent treatment.

\medterm{Ethionamide} An antibiotic used with other medicines for drug-resistant or difficult-to-treat tuberculosis. It interferes with mycobacterial cell-wall production and may cause nausea, liver injury, nerve symptoms, or thyroid problems.

\medterm{Ethmocephaly} A severe birth defect in which the eyes are very close together and the nose is replaced or accompanied by a tube-like structure. It usually occurs with major brain-development abnormalities and is often life-limiting.

\medterm{Ethmoidal Air Cells} Multiple small air-filled spaces within the ethmoid bone between the nose and eye sockets. They are part of the sinus system and can become blocked or infected, causing sinus pressure and discharge.

\medterm{Ethmoidal Sinus} A group of small, thin-walled air spaces in the ethmoid bone between the eye sockets and above the nasal cavity. They are divided into front and back groups according to where they drain.

\synonyms
Ethmoid Air Cells; Paranasal Ethmoid Sinus

\medterm{Ethmoid Bone} A light, complex bone between the eye sockets that forms part of the nasal septum, nasal roof, and inner eye socket. It contains air cells and openings for smell nerves, so injury can affect smell or the eye.

\medterm{Ethmoiditis} Inflammation or infection of the ethmoid air spaces between the nose and eyes. It may cause fever, nasal blockage, discharge, facial pain, or swelling and can rarely spread to the eye socket or brain.

\medterm{Ethmoid Sinus} One of several small air spaces in the ethmoid bone between the nose and eye sockets. Inflammation or infection can cause nasal blockage, headache, or eye-area pain and may spread to nearby tissues.

\medterm{Ethmoid Sinusitis} Inflammation or infection of the air spaces between the nose and eyes. It may cause congestion, discharge, fever, headache, or pain around the eyes; eye swelling, vision changes, or severe illness requires urgent care.

\medterm{Ethnicity} A social and cultural identity shaped by shared ancestry, history, language, traditions, or community; it is not a biological diagnosis and should not be used as a substitute for individual risk assessment.

\medterm{Ethoglucid} An older sclerosing substance used historically to treat selected vascular lesions or hemorrhoids; modern use is limited by tissue-injury risks and alternatives.

\medterm{Ethosuximide} An anticonvulsant used mainly for absence seizures; common adverse effects include gastrointestinal symptoms, fatigue, and headache, with rare blood or skin reactions.

\medterm{Ethoxzolamide} A carbonic-anhydrase-inhibiting medicine that reduces formation of certain body fluids. It has been used for selected eye and metabolic conditions; possible effects include tingling, nausea, kidney stones, and electrolyte changes.

\medterm{Ethyl Chloride} A volatile liquid sprayed on skin to produce brief cooling and numbness before minor procedures, but misuse can cause cold injury or toxicity.

\medterm{Ethylene} A colorless, flammable gas used in chemical manufacturing and to ripen fruit. High concentrations can displace oxygen and cause headache, dizziness, drowsiness, unconsciousness, or suffocation. Remove a person from exposure only when it is safe and seek emergency help for breathing or consciousness problems.

\medterm{Ethylene Chlorohydrin} A toxic industrial chemical used in some manufacturing processes. Exposure can irritate the eyes, skin, and airways and may affect the nervous system, liver, or kidneys. It can be absorbed through skin, so contaminated clothing should be removed safely and urgent poison-control advice obtained.

\medterm{Ethylene Dibromide} A toxic chemical formerly used as a soil fumigant and fuel additive. It can irritate the eyes and lungs and may damage the liver, kidneys, or nervous system; long-term exposure is associated with cancer risk. Suspected exposure needs urgent poison-control or occupational-health assessment.

\medterm{Ethylene Glycol} A toxic alcohol found in some antifreeze and industrial products. Ingestion can cause severe metabolic acidosis, neurologic impairment, kidney injury, and death.

\medterm{Ethylene Glycol Blood Test} A blood test measuring ethylene glycol, a toxic alcohol found in some antifreeze products. Detection supports diagnosis of poisoning, but treatment must not wait for the result when the exposure and metabolic findings are strongly suggestive.

\medterm{Ethylene Glycol Poisoning} Poisoning after swallowing ethylene glycol, found in some antifreeze products. Early effects may resemble alcohol intoxication, followed by severe acid buildup, seizures, coma, and kidney failure; emergency antidote and sometimes dialysis are needed.

\medterm{Ethylene Oxide} A flammable gas used to sterilize heat-sensitive medical equipment and make other chemicals. It can irritate the skin and airways, cause nerve injury, and with repeated exposure increase cancer risk. Sterilized equipment must be adequately aerated before use, and workplace exposure requires strict controls.

\medterm{Ethyl Ether} A highly flammable liquid formerly used as an anesthetic and still used as a laboratory solvent. Breathing its vapor can cause irritation, dizziness, drowsiness, or loss of consciousness; swallowing it can injure the lungs if it is aspirated. Exposure requires fresh air, fire precautions, and medical assessment when symptoms are severe.

\medterm{Ethynodiol Diacetate} A synthetic progestogen formerly used in some oral contraceptives. It helps prevent pregnancy mainly by suppressing ovulation and changing cervical mucus and the uterine lining.

\medterm{Etidronic Acid} A medicine related to bisphosphonates that can slow the breakdown of bone by attaching to bone mineral. It has also been used to bind minerals in some nonmedical products; effects and safety depend on the use and dose.

\medterm{Etilefrine} A medicine that stimulates some of the body’s adrenaline-like receptors and can raise blood pressure. It increases blood-vessel tightening and heart activity, so it may cause fast heartbeat, headache, or other effects and requires medical supervision.

\medterm{Etiology} The cause or set of causes of a disease or health condition, including genetic, infectious, environmental, immune, traumatic, and behavioural factors.

\medterm{Etodolac} A nonsteroidal anti-inflammatory drug used for pain and inflammation; gastrointestinal bleeding, kidney injury, cardiovascular events, and hypersensitivity are important risks.

\medterm{Etomidate} A short-acting intravenous anesthetic used for induction, particularly when cardiovascular stability is desired; it can cause myoclonus, nausea, and transient adrenal suppression.

\medterm{Etonogestrel} A man-made form of the hormone progesterone used in a small contraceptive implant placed under the skin of the upper arm. It mainly prevents ovulation and also thickens cervical mucus; irregular bleeding and other hormone-related effects can occur.

\medterm{Etoposide} A chemotherapy medicine that interferes with DNA repair and cell division. It treats several cancers but can cause low blood counts, infection, hair loss, nausea, mouth sores, or a later risk of leukemia.

\medterm{Etoposide Phosphate} A water-soluble form of etoposide that changes into etoposide inside the body. It provides another method of giving the chemotherapy medicine, with similar cancer treatment effects and potential toxicities.

\medterm{Etoricoxib} An anti-inflammatory pain medicine that selectively blocks an enzyme involved in pain and swelling. It may help arthritis or acute pain but can raise blood pressure and increase heart, kidney, fluid-retention, and stomach risks.

\medterm{Eubacterium} A group of bacteria that usually live without oxygen in the mouth and intestine. Most are part of normal microbial communities, but some can rarely cause opportunistic infection in seriously ill people.

\medterm{Eubacterium Limosum} A bacterium that lives without oxygen in the intestine and helps break down nutrients. It is rarely associated with human infection, and its importance in a sample depends on symptoms and the collection site.

\medterm{Eubacterium Rectale} A common intestinal bacterium that helps ferment dietary fiber and produce short-chain fatty acids. These products support the colon and contribute to normal gut microbial activity.

\medterm{Eucalyptus Oil} A concentrated plant oil containing eucalyptol and other chemicals. Swallowing it can poison the nervous system and cause vomiting, drowsiness, seizures, or breathing problems; it should not be taken internally without medical advice.

\medterm{Eucalyptus Oil Overdose} Poisoning after swallowing too much eucalyptus oil. Burning, vomiting, drowsiness, confusion, seizures, and breathing problems may occur; do not induce vomiting and seek emergency poison-control advice.

\medterm{Euchromatin} A relatively open form of DNA packaging inside the cell nucleus. Its structure makes nearby genes more accessible for reading and activity than DNA in tightly packed regions.

\medterm{Eugenicist} A person who supports eugenics, a discredited belief that society should control reproduction to alter inherited traits in a population. Eugenic policies caused coercion, discrimination, and forced sterilization and have no place in ethical healthcare.

\medterm{Eugenics} A discredited ideology that seeks to change a population by encouraging some people to reproduce and preventing others from doing so. Its history includes racism, forced sterilization, institutional abuse, and genocide.

\medterm{Eugenol Oil Overdose} Poisoning after excessive exposure to clove-derived eugenol oil. It may irritate the mouth and stomach and cause liver injury, bleeding, seizures, low blood pressure, slowed breathing, or coma.

\medterm{Euglycemic Diabetic Ketoacidosis} A severe metabolic complication of diabetes mellitus characterized by profound anion gap metabolic acidosis and elevated serum ketones occurring with normal or only mildly elevated blood glucose levels (less than 250 mg/dL), frequently triggered by SGLT2 inhibitor therapy, restricted carbohydrate intake, or acute physiological stress.
\synonyms
EuDKA; Normal Glycemic Ketoacidosis

\medterm{Euphoria} An unusually intense feeling of happiness or well-being. It may occur normally, with mania, brain disease, or some medicines and drugs; sudden extreme euphoria with unsafe behavior needs clinical assessment.

\medterm{Eustachian Tube} A passage connecting the middle ear to the back of the nose. It opens during swallowing or yawning to equalize pressure, drain fluid, and help protect the middle ear from throat secretions.

\medterm{Eustachian Tube Dysfunction} Poor opening or closing of the passage between the middle ear and back of the nose. It can cause ear pressure, popping, muffled hearing, pain, or fluid after colds, allergies, or pressure changes.

\medterm{Eustachian-Tube Dysfunction} Failure of the passage between the middle ear and back of the nose to equalize pressure or clear fluid. It may cause fullness, popping, reduced hearing, or ear pain after infection, allergies, or pressure changes.

\medterm{Eustachian Tube Patency} The ability of the passage between the middle ear and back of the nose to remain open enough to equalize pressure and drain fluid. Poor patency can cause fullness, muffled hearing, or middle-ear fluid.

\medterm{Eustachian Tube Problems} Problems with the ear-to-nose passage can cause fullness, popping, pain, muffled hearing, or repeated middle-ear fluid or infection. Colds, allergies, inflammation, and anatomy may contribute; persistent hearing loss needs examination.

\synonyms
Eustachian Tube Dysfunction (ETD); Auditory Tube Obstruction

\medterm{Eustachian Valve} A crescentic embryological flap of tissue at the junction of the inferior vena cava and the inferior right atrium that in fetal life directs oxygenated umbilical venous blood across the foramen ovale into the left atrium; usually persists as a benign vestigial remnant in adults.

\synonyms
Valve of the Inferior Vena Cava

\medterm{Euthanasia} The deliberate ending of a person’s life by another person to relieve suffering. It is ethically and legally distinct from allowing a natural death, stopping unwanted treatment, or providing comfort-focused palliative care.

\medterm{Euthanasia And Assisted Dying} Terms describing legally regulated practices intended to end life or help a person end life to relieve suffering. Laws and safeguards differ widely; they are distinct from palliative care and withdrawal of unwanted treatment.

\medterm{Euthyroid} Having normal thyroid function, usually shown by thyroid-stimulating hormone and thyroid-hormone levels within the laboratory’s reference ranges. A person can be euthyroid even when another illness temporarily changes some thyroid test results.

\medterm{Euthyroid Goiter} An enlarged thyroid gland with normal thyroid-hormone production. It may result from iodine imbalance, autoimmune disease, nodules, medicines, or inherited factors and should be assessed if persistent, enlarging, or causing swallowing or breathing problems.

\medterm{Euthyroid Sick Syndrome} Temporary changes in thyroid blood tests during severe illness, starvation, or major surgery without a primary thyroid disorder. Treatment usually focuses on the underlying illness rather than thyroid hormone.

\medterm{Eutocia} Normal labor and birth progressing without major obstruction or complication. The term describes the course of delivery, not a separate disease or treatment.

\medterm{Euvolemic Hyponatremia} Low blood sodium with no obvious clinical signs of excess or deficient body fluid volume. Common mechanisms include inappropriate antidiuretic hormone activity, adrenal insufficiency, hypothyroidism, and excess water intake.

\medterm{Evacuation} Removal of people from danger or removal of material from a body cavity, organ, or wound. In clinical care it may refer to clearing a hematoma, bowel, uterus, or other collection; the context determines the meaning.

\medterm{Evacuation Of Uterus} Removal of pregnancy tissue or other retained material from the uterus using medicines, suction, or a surgical procedure. It may treat miscarriage, abortion, retained placenta, or bleeding and carries risks of infection, bleeding, or uterine injury.

\medterm{Evans Syndrome} An autoimmune disorder characterized by the simultaneous or sequential development of autoimmune hemolytic anemia (positive direct antiglobulin test) and immune thrombocytopenic purpura, occasionally accompanied by autoimmune neutropenia.
\synonyms
Evans's Syndrome; AIHA and ITP Syndrome

\medterm{Event-Free Survival} The length of time after treatment begins or a study starts before a defined event occurs, such as disease progression, relapse, a complication, or death. The exact event must be stated for comparisons to be meaningful.

\medterm{Event Monitor} A portable device that records the heart’s electrical rhythm when symptoms occur or at scheduled intervals over days or weeks. It can help detect intermittent abnormal rhythms that a brief ECG may miss.

\medterm{Everolimus} A medicine that suppresses a cell-growth pathway and immune activity. It is used for selected cancers, transplant care, and some genetic disorders; mouth sores, infection, high blood sugar, lung inflammation, and low blood counts may occur.

\medterm{Eversion} Turning the sole of the foot outward, away from the body’s midline. The movement uses several ankle and foot joints and can be limited by injury, arthritis, or nerve and muscle disorders.

\medterm{Evidence-Based Practice} Care that combines the best available research, professional clinical judgment, and the patient’s goals, values, and circumstances. It supports informed decisions while recognizing that evidence and individual needs may differ.

\medterm{Evisceration} Protrusion of internal organs through a surgical abdominal wound that has opened. This is a surgical emergency; the wound should be covered with sterile moist material, and urgent assessment is needed for repair and infection prevention.

\medterm{Evoked Potential} An electrical response recorded from the nervous system after a visual, sound, or touch stimulus. The test assesses whether signals travel normally along sensory pathways and may reveal nerve damage or slowed transmission.

\medterm{Evoked Potentials} Tests that record nervous-system responses to visual, sound, or touch stimuli. They can assess pathways for sight, hearing, or sensation and may help identify slowed signaling from disorders affecting nerves, the brain, or spinal cord.

\synonyms
VEP; BAER; SSEP; Evoked Responses

\medterm{Evolution} Gradual change in inherited traits of a population across generations. In clinical writing, the term can also describe how a disease, symptom, or treatment response changes over time.

\medterm{Evolutionarily Conserved Gene} A gene whose sequence or function remains similar across different species. This similarity suggests that the gene performs an important biological role, although it does not by itself prove a disease association.

\medterm{Ewart Sign} A physical-examination finding of a large fluid collection around the heart: dullness, reduced breath sounds, and bronchial breathing over the lower left back. Modern imaging is more reliable for confirming the cause.

\medterm{Ewing Sarcoma} A rare, aggressive cancer of bone or soft tissue, most often affecting children, teenagers, and young adults. It may cause persistent pain, swelling, fever, or a fracture and requires biopsy and specialist treatment.

\medterm{Ewing's Sarcoma} A malignant tumor of bone or soft tissue, most common in children and young adults. It may cause localized pain, swelling, or fever and is diagnosed with imaging, biopsy, and molecular testing; treatment usually combines chemotherapy with surgery and/or radiation.

\medterm{Ewing's Tumour} A rare aggressive cancer that usually begins in bone or nearby soft tissue in children or young adults. Persistent pain, swelling, fever, or an unexplained fracture may occur; diagnosis requires imaging and biopsy.

\medterm{Exacerbate} To make a disease, symptom, or medical problem worse. Infection, stress, exertion, allergens, medicines, or delayed treatment may exacerbate symptoms, depending on the condition.

\medterm{Exacerbation} A worsening or flare of a disease or its symptoms after a period of relative stability. Triggers may include infection, allergens, treatment interruption, exertion, or another physical or emotional stress.

\medterm{Exacerbation Of COPD} A sudden worsening of breathlessness, cough, or sputum beyond usual daily variation in chronic obstructive pulmonary disease. Infections and air pollution are common triggers; severe breathlessness, confusion, or blue lips requires urgent care.

\medterm{Exaggerated Placental Site} A noncancerous increase in placental cells at the place where the placenta attaches to the uterus. It may resemble a pregnancy-related cancer in a tissue sample, so an experienced pathologist must distinguish it from cancer.

\medterm{Examination} A systematic assessment of health using observation, touch, listening, measurements, and selected tests. It may be general or focused on a complaint and helps identify normal findings, illness, injury, and urgency.

\medterm{Examination Table} A padded, adjustable surface that supports a person during examination, testing, minor procedures, or treatment. It should be stable, cleanable, accessible, and used with appropriate safety measures.

\medterm{Exanthem} A widespread skin rash, often caused by an infection or medicine. Its appearance, timing, fever, mucosal involvement, and other symptoms help clinicians determine whether it may be contagious or serious.

\medterm{Exanthema} A widespread skin eruption associated with an illness affecting the body. Viral infections are common causes, but medicines, bacteria, and immune disorders may also cause it; blisters, skin pain, mouth sores, or severe illness need prompt assessment.

\medterm{Exanthema Subitum} A common viral illness in infants and toddlers, usually caused by human herpesvirus 6 or 7. It often causes three to five days of high fever followed by a pink rash as the fever settles.

\synonyms
Roseola Infantum; Sixth Disease

\medterm{Exanthematous Pustulosis} A sudden medicine-related eruption of many small, noninfectious pustules on red skin, often with fever. It usually improves after the responsible medicine is stopped under medical guidance, but severe reactions require urgent assessment.

\medterm{Exanthem Subitum} A common viral illness in infants and toddlers, usually caused by human herpesvirus 6 or 7. It often causes three to five days of high fever followed by a pink rash as the fever settles.
\synonyms
Roseola Infantum; Sixth Disease

\medterm{Excavation} A hollow, depression, or removal of tissue or material. In medicine it may describe a normal body feature, a disease-related defect, or the surgical creation or clearing of a space.

\medterm{Exceptional Release} Permission to use a donated organ that does not meet usual acceptance criteria when clinicians judge it suitable for a particular recipient. The decision requires careful risk assessment, documentation, and informed consent.

\medterm{Excessive Bruising} Bruises that occur often, are unusually large, appear without clear injury, or last unusually long. Causes include medicines, fragile skin, low platelets, clotting disorders, liver disease, or other illness.

\medterm{Excessive Daytime Sleepiness} An excessive tendency to feel drowsy or fall asleep during waking hours, despite an opportunity for adequate nighttime sleep. It may occur with insufficient sleep, sleep disorders, medicines, medical or neurologic conditions, or mood disorders.

\medterm{Excessive Menstrual Bleeding} Alternate name; see \textit{Heavy Menstrual Bleeding}.

\medterm{Excessive Scarring} Alternate name; see \textit{Keloid}.

\medterm{Excessive Sweating} Sweating more than needed to cool the body, either all over or in specific areas such as the hands, feet, or armpits. It may result from hyperhidrosis, medicines, infection, hormones, or another illness.

\medterm{Excessive Tearing} Tears overflowing onto the face because the eyes make too many tears or drainage is blocked. Irritation, allergy, infection, dry eye, eyelid problems, or a blocked tear duct may be responsible.

\medterm{Excessive Thirst} Polydipsia, an abnormally strong or persistent desire to drink, commonly associated with dehydration, diabetes mellitus, diabetes insipidus, medicines, or psychiatric conditions.

\medterm{Excessive Yawning} Excessive yawning is unusually frequent or difficult-to-control yawning. It may accompany sleep deprivation, excessive daytime sleepiness, medication effects, vasovagal reactions, or less commonly neurologic or cardiopulmonary disease; persistent unexplained episodes merit clinical evaluation.

\medterm{Excess Post-Exercise Oxygen Consumption} Increased oxygen use after exercise as the body returns to its usual metabolic and physiologic state.

\medterm{Exchange Transfusion} A procedure that repeatedly removes small amounts of a baby’s blood and replaces them with compatible donor blood. It can rapidly reduce dangerous bilirubin or antibody-related anemia when other treatment is insufficient.

\medterm{Excise} To remove tissue, a growth, or a lesion by cutting it out. The removed material may be examined in a laboratory to identify the cause and determine whether abnormal cells reach the edges.

\medterm{Excision} Surgical removal of part or all of a lesion, tissue, or organ. The specimen may be examined to establish a diagnosis, and removal may also treat disease when the affected area can be safely removed.

\medterm{Excisional Biopsy} Removal of an entire lump or abnormal area, sometimes with a small margin of nearby tissue, so a laboratory can examine it. It may provide a diagnosis and, for selected small lesions, complete treatment.

\medterm{Excision Biopsy} Alternate name; see \textit{Excisional Biopsy}.

\medterm{Excitation} Activation or increased responsiveness of a cell, nerve, muscle, molecule, or body system after a stimulus. In nerves and muscles, it can start electrical signals that produce movement or other functions.

\medterm{Excitation-Contraction Coupling} The chain of events that links an electrical signal in a muscle cell to contraction. The signal causes calcium release inside the cell, allowing muscle proteins to interact and generate force.

\medterm{Excitatory Neurotransmitter} A chemical messenger that makes a receiving nerve cell more likely to generate an electrical signal. Examples include glutamate and, in some settings, acetylcholine; excessive activity can injure nerve cells.

\medterm{Excitatory Synapse} A junction between nerve cells where chemical signaling makes the receiving cell more likely to produce an electrical impulse. Balanced excitatory and inhibitory signals are needed for normal brain and nerve function.

\medterm{Excitotoxicity} The pathological process of neuronal injury and cell death resulting from excessive or prolonged synaptic stimulation by excitatory neurotransmitters, primarily glutamate, triggering massive intracellular calcium influx, mitochondrial dysfunction, reactive oxygen species generation, and apoptotic protease activation.

\medterm{Excitotoxin} A substance that overstimulates nerve cells and can damage or kill them. Excessive signaling may allow too much calcium into cells, contributing to injury in some brain diseases or poisoning states.

\medterm{Exclamation Point Hair} A short broken hair that is thinner at its outer end and wider near the scalp. It is a characteristic finding of alopecia areata but can occasionally appear in other hair disorders.

\medterm{Exclusion Diet} A planned diet that removes a suspected food and later reintroduces it to see whether symptoms recur. It should be supervised when possible and must not replace emergency allergy evaluation when severe reactions or anaphylaxis are possible.

\medterm{Excoriate} To remove or injure the skin surface by scratching, rubbing, or picking. Repeated excoriation can cause open sores, infection, scarring, and difficulty healing.

\medterm{Excoriation} A superficial skin injury caused by scratching, rubbing, or picking. It may cause pain, bleeding, crusting, infection, or scarring and usually heals when the irritation or behavior stops.

\medterm{Excoriation Disorder} A mental-health disorder involving recurrent skin picking that causes skin lesions and clinically significant distress or impairment. Repeated unsuccessful efforts to reduce or stop the behavior are common.

\synonyms
Skin-Picking Disorder

\medterm{Excreta} Material expelled from the body, including stool, urine, sweat, vomit, and respiratory secretions. Some excreta carry infectious organisms, so handwashing, protective equipment, cleaning, and safe disposal help prevent spread.

\medterm{Excretion} Removal of waste, medicines, or other substances from the body, mainly through urine, bile, breath, stool, sweat, or breast milk. Kidney and liver function strongly affect how long many substances remain in the body.

\medterm{Executive Function} Mental skills used to plan, focus, remember instructions, control impulses, solve problems, make decisions, and switch between tasks. Problems may follow brain injury, dementia, developmental conditions, or mental illness.

\medterm{Executive Functions} Mental abilities that support planning, organization, attention, memory, decision-making, self-control, and flexible behavior. Difficulties can affect work, school, relationships, and self-care and may have neurologic, developmental, or psychiatric causes.

\synonyms
Executive Control

\medterm{Exemestane} A medicine that lowers estrogen production by blocking an enzyme used after menopause. It treats selected hormone-sensitive breast cancers and may cause hot flushes, joint pain, bone loss, fatigue, or abnormal cholesterol.

\medterm{Exenatide} An injectable medicine that imitates a gut hormone, increasing insulin release when glucose is high, reducing glucagon, and slowing stomach emptying. It treats type 2 diabetes but may cause nausea, dehydration, or pancreatitis.

\medterm{Exencephaly} A severe birth defect in which much of the developing brain lies outside the skull because the skull and protective coverings did not form normally. It usually progresses to anencephaly and is not compatible with long-term survival.

\medterm{Exenteration} Extensive surgery removing organs and tissues from a body compartment, such as the pelvis or eye socket. It is reserved for selected advanced cancers or destructive disease and can cause major functional and emotional effects.

\medterm{Exercise} Planned physical activity intended to improve or maintain strength, endurance, flexibility, balance, function, or health. Aerobic, resistance, and balance activities should be chosen and increased according to a person’s abilities and medical risks.

\synonyms
Physical Exercise

\medterm{Exercise And Activity For Weight Loss} Physical activity used with eating changes to increase energy use, preserve muscle, and support gradual weight loss. The safest plan depends on fitness, medical conditions, joint health, and personal ability.

\medterm{Exercise And Age} The way physical activity needs and responses may change with age; safe exercise should match fitness, health, and ability.

\medterm{Exercise And Immunity} The relationship between physical activity and immune function; regular moderate activity generally supports health, while excessive exertion may increase short-term illness risk.

\medterm{Exercise Headaches} Headaches occurring during or after strenuous physical activity. Primary exercise headaches are recurrent and usually not caused by an underlying disorder, but a first, abrupt, severe, or atypical exertional headache requires prompt medical assessment to exclude secondary causes such as subarachnoid hemorrhage, arterial dissection, or reversible cerebral vasoconstriction syndrome.

\medterm{Exercise-Induced Arrhythmias} Abnormal heart rhythms that begin or worsen during physical exertion. They may result from inherited electrical disorders, heart muscle disease, or reduced blood flow; fainting, chest pain, or palpitations during exercise needs urgent evaluation.

\medterm{Exercise-Induced Asthma} Airway narrowing during or soon after exercise, causing cough, wheezing, chest tightness, or breathlessness. It can occur with or without persistent asthma and is managed with an individualized prevention and treatment plan.

\medterm{Exercise-Induced Bronchoconstriction} Temporary narrowing of the airways during or after exercise, causing cough, wheezing, chest tightness, or shortness of breath. Cold, dry air can worsen it; diagnosis and preventive inhaler use should be guided clinically.

\medterm{Exercise Intolerance} An inability to do as much physical activity as expected because of unusual tiredness, breathlessness, chest discomfort, leg pain, weakness, or dizziness. Heart, lung, blood-vessel, muscle, nerve, blood, metabolic, and conditioning problems can cause it.

\medterm{Exercises To Help Prevent Falls} Activities that improve leg strength, balance, flexibility, walking, and confidence. A program should match ability and medical conditions; supervised practice may be safest for people with previous falls or severe weakness.

\medterm{Exercise Stress Test} A test that records heart rate, rhythm, blood pressure, and symptoms while a person walks on a treadmill or pedals a stationary bicycle. It assesses exercise capacity and may suggest reduced heart blood flow or an abnormal rhythm.

\medterm{Exercise Tolerance} The ability to perform physical activity before symptoms such as breathlessness, fatigue, chest discomfort, or pain limit activity.

\medterm{Exfoliation Syndrome} An age-related eye disorder in which abnormal material collects on the lens and other front-eye structures. It increases the risk of glaucoma and can make cataract surgery more complicated, so regular eye checks are important.

\medterm{Exfoliative Cytology} Examination of cells that naturally shed from a body surface or are gently collected. It may help investigate infection, inflammation, or cancer, but abnormal or uncertain results may require tissue biopsy.

\medterm{Exfoliative Dermatitis} Severe redness, scaling, and shedding affecting most of the skin. Medicines, psoriasis, eczema, lymphoma, and other diseases may cause it; fluid loss, infection, temperature problems, and heart strain make urgent medical assessment necessary.

\medterm{Exhalation} Breathing air out of the lungs. It is usually passive at rest as the lungs recoil, but muscles actively assist exhalation during exercise, coughing, or breathing difficulty.

\medterm{Exhaled Nitric Oxide} A breath test that estimates one type of airway inflammation. A raised result can support asthma assessment and suggest likely benefit from inhaled steroids, but smoking, infection, allergies, and current treatment can affect the result.

\medterm{Exhaustion} Extreme physical or mental tiredness that reduces the ability to function. It may follow exertion, poor sleep, illness, dehydration, anemia, hormone disorders, medicines, or stress; severe or persistent symptoms need evaluation.

\medterm{Exhaustion, Heat} Alternate name; see \textit{Heat Exhaustion}.

\medterm{Exhibitionism} Sexual arousal from exposing one’s genitals to an unsuspecting person. Repeated nonconsensual exposure can harm others and becomes a diagnosable disorder when urges or actions cause significant distress, impairment, or risk.

\medterm{Exhibitionistic Disorder} A mental-health disorder involving repeated sexual arousal from exposing the genitals to an unsuspecting person, with action on the urges or clinically significant distress or impairment.

\medterm{Exhumation} Lawful removal of a buried body or remains for identification, further testing, investigation, or repatriation. It requires legal authorization and respectful handling; testing may include toxicology, DNA analysis, or examination of injuries.

\medterm{Existential Distress} Emotional suffering related to meaning, identity, purpose, freedom, mortality, or loss of control. It can occur during serious illness and may improve with empathetic communication, counseling, spiritual care, and symptom relief.

\medterm{Exit Wound} The wound where a projectile leaves the body. Its size and shape vary with the projectile, tissue, clothing, and bone, so appearance alone cannot prove the direction or identify every injury; complete examination is required.

\medterm{Exocervix} The part of the cervix that projects into the vagina and is covered mainly by flat squamous cells. It is examined during pelvic examination and cervical screening.

\medterm{Exocrine} Relating to glands that release substances through ducts onto a body surface or into a cavity. Examples include sweat, salivary, tear, and digestive glands.

\medterm{Exocrine Gland} A gland that sends its secretion through a duct to a body surface or internal cavity. Sweat, salivary, tear, and digestive glands are examples; their secretions help protect, lubricate, or digest.

\medterm{Exocrine Glands} Glands that release substances through ducts onto a surface or into a body cavity. They include sweat, salivary, tear, and digestive glands and support protection, lubrication, digestion, and temperature control.

\medterm{Exocrine Pancreas} The part of the pancreas that releases digestive enzymes and bicarbonate into the small intestine. These secretions help break down fats, proteins, and carbohydrates and neutralize stomach acid.

\medterm{Exocrine Pancreatic Insufficiency} Failure of the pancreas to release enough digestive enzymes, causing poor digestion and absorption. It may lead to greasy stools, bloating, diarrhea, weight loss, and vitamin deficiency.

\medterm{Exocrine Pancreatic Insufficiency EPI} Alternate index label; see \textit{Exocrine Pancreatic Insufficiency}.

\synonyms
EPI; Pancreatic Exocrine Insufficiency

\medterm{Exocytosis} The process by which a cell releases substances outside itself. Small sacs inside the cell join its outer membrane and release hormones, nerve signals, digestive enzymes, or other proteins.

\medterm{Exogenous} Coming from outside the body rather than being produced within it. An exogenous substance may be a medicine, hormone, nutrient, toxin, infection, or other externally introduced material.

\medterm{Exogenous Cushing Syndrome} A condition caused by prolonged or excessive corticosteroid treatment, producing effects similar to too much natural cortisol. Features may include weight gain, easy bruising, high blood pressure, diabetes, weakness, and bone loss.

\medterm{Exomphalos} A birth defect in which abdominal organs protrude through the belly button into the base of the umbilical cord and are covered by a membrane. Newborns need protection of the sac and specialist surgical care.

\medterm{Exon} A segment of a gene that remains in the mature messenger RNA after unused sections are removed. It may contain instructions for making part of a protein or, in some genes, help regulate gene activity.

\medterm{Exopeptidase} An enzyme that removes amino acids from the end of a protein. Digestive exopeptidases help break dietary proteins into smaller units that the intestine can absorb.

\medterm{Exophiala} A group of darkly pigmented environmental fungi that can rarely cause skin, lung, eye, or widespread infection. Disease is more likely in people with weakened immunity or implanted medical devices.

\medterm{Exophthalmic Goitre} An enlarged thyroid associated with Graves disease, in which immune stimulation causes excess thyroid hormone. Graves disease may also cause eye dryness, pain, swelling, or bulging.

\medterm{Exophthalmos} The abnormal forward protrusion or displacement of one or both eyeballs from the orbit, most frequently caused by Graves disease (thyroid eye disease), orbital tumors, or inflammation.

\synonyms
Bulging Eyes; Proptosis; Ocular Proptosis

\medterm{Exosomes} Very small membrane-covered particles released by cells that carry proteins, fats, and genetic material to other cells. They help cells communicate and are being studied for diagnosis and drug delivery, but many medical uses remain experimental.

\medterm{Exostosis} A benign growth of bone. It may be harmless, but location can cause pain, pressure, restricted movement, or blockage; repeated cold-water exposure can produce bony growths in the ear canal.

\synonyms
Surfer's Ear; Osteoma of External Auditory Canal

\medterm{Exotoxins} Poisonous proteins released by some bacteria that damage particular cells or disrupt body functions. They cause diseases such as botulism, tetanus, diphtheria, or cholera and may require urgent treatment.

\medterm{Exotropia} A form of misaligned eyes in which one eye turns outward. It may be occasional or constant and can cause double vision, reduced depth perception, or loss of vision development in children; eye assessment guides treatment.

\medterm{Expectant Treatment} Careful observation with regular review instead of immediate medicine or surgery when immediate intervention is not necessary. It includes clear follow-up and warning signs because treatment may be needed if symptoms worsen or tests change.

\medterm{Expectorants} Medicines intended to loosen or increase airway mucus so it can be coughed up more easily. Benefits vary, and fluids, humidification, and treatment of the underlying illness may be more important than an expectorant.

\medterm{Expectoration} Coughing up and bringing mucus or phlegm out of the respiratory tract. Changes in amount, color, smell, blood, fever, or breathlessness may indicate infection or another lung problem.

\medterm{Experience} A person’s direct knowledge, perception, or emotional response to an event. In healthcare, patient experience includes how care is accessed, communicated, coordinated, and received.

\synonyms
Subjective Experience; Patient-Reported Experience

\medterm{Expiration} Movement of air out of the lungs, usually caused by relaxation of the breathing muscles and natural recoil of the lungs. During forceful breathing, abdominal and chest muscles actively push air out.

\medterm{Expiratory Grunting} A grunting sound made during breathing out, caused by partially closing the voice box. It can briefly help keep small air spaces open but, especially in a baby or child, often signals respiratory distress and needs assessment.

\medterm{Expiratory Reserve Volume} The maximum additional volume of air that can be actively and forcefully exhaled from the lungs at the end of a normal, quiet tidal expiration, normally measuring approximately 1.0 to 1.2 liters in healthy adults.

\synonyms
ERV

\medterm{Explant} To remove an implanted device, prosthesis, or tissue. The term can also mean tissue removed from an organism and maintained in laboratory culture for study.

\medterm{Explantation} Surgical removal of an implanted device or prosthesis. It may be needed because of infection, malfunction, migration, breakage, tissue injury, intolerance, or completion of temporary treatment.

\medterm{Exploratory Laparotomy} Open surgery to enter the abdomen and directly find or treat a serious problem such as bleeding, perforation, blocked bowel, or loss of blood flow when less-invasive tests or treatment are insufficient.

\medterm{Exploratory Operation} Surgery performed to inspect an area directly when scans, tests, or less-invasive procedures cannot establish the cause or extent of disease. Tissue may be sampled and treatment may be performed during the operation.

\medterm{Exposure} Contact with a physical, chemical, biological, or psychological agent. Health risk depends on what the agent is, how much contact occurred, how it entered the body, and how long exposure lasted.

\medterm{Exposure Route} The way a substance enters the body, such as by breathing, swallowing, skin contact, injection, or the eyes. The route affects how quickly and how much of the substance is absorbed and how toxic it may be.

\medterm{Expression} In genetics, the process by which a gene’s information is used to make RNA or protein. In clinical examination, expression can also mean the outward display of emotion or symptoms; altered gene activity can contribute to disease.

\medterm{Expression Hopping} A shift from one form of addictive behavior to another, such as replacing substance use with compulsive gambling. The underlying loss of control and harm may persist even when the specific behavior changes.

\medterm{Expressivity} The degree to which a genetic variant’s effects appear in a person. Variable expressivity means that people with the same variant may have different features or different severity of the associated condition.

\medterm{Expulsion, Stage Of} The second stage of labor, beginning when the cervix is fully open and ending when the baby is born. It includes pushing and may be shortened or assisted when the mother or baby is at risk.

\medterm{Exsanguination} Life-threatening loss of blood from severe internal or external bleeding. It can rapidly cause shock, organ failure, and death; immediate pressure or bleeding control, emergency care, and blood replacement may be needed.

\medterm{Exserohilum} A group of darkly pigmented molds that can rarely cause sinus, eye, skin, or invasive infection. Serious disease is more likely in people with weakened immunity and requires specialist testing and treatment.

\medterm{Extend} To straighten a joint or increase the angle between two bones. For example, extending the elbow moves the forearm away from the upper arm toward a straight position.

\synonyms
Extension

\medterm{Extension} Movement that straightens a joint or increases the angle between bones, such as straightening the knee or elbow. Excessive or forced extension can strain or tear muscles, tendons, ligaments, cartilage, or the joint.

\medterm{Extensively Drug-Resistant Tuberculosis} Tuberculosis caused by bacteria resistant to several key medicines, including isoniazid, rifampin, a fluoroquinolone, and at least one other important medicine. It is difficult to treat and requires specialist, susceptibility-guided therapy.

\synonyms
XDR-TB; XDR Tuberculosis

\medterm{Extensor} A muscle or other structure that straightens a joint or increases the angle between bones. Extensor muscles commonly act opposite flexor muscles and help position limbs and fingers.

\medterm{Extensor Digitorum} A forearm muscle whose tendons straighten the fingers and help extend the wrist. Overuse near its origin can contribute to outer-elbow pain commonly called tennis elbow.

\medterm{Extensor Plantar Response} An abnormal neurological reflex in which firm tactile stroking of the lateral plantar border of the foot produces slow extension (dorsiflexion) of the great toe accompanied by fanning of the other toes, indicating corticospinal (pyramidal) tract impairment.
\synonyms
Positive Babinski Sign; Babinski Reflex

\medterm{Extent} The size, severity, range, or spread of a disease, injury, symptom, or treatment effect. Describing extent helps determine prognosis, further testing, and treatment.

\medterm{External Auditory Canal} The passage from the outer ear to the eardrum. It is lined with skin, hair, and wax-producing glands that help protect the ear, but it can become blocked, irritated, or infected.

\medterm{External Beam Partial Breast Irradiation} An accelerated conformal radiation therapy technique delivered via external beam linac that targets exclusively the lumpectomy cavity margin following breast-conserving surgery, sparing the remainder of the ipsilateral breast tissue in selected early-stage breast cancer patients.

\synonyms
EB-APBI; Accelerated Partial Breast Irradiation; Targeted Breast Radiotherapy

\medterm{External Beam Radiation Therapy} Cancer treatment delivered from a machine outside the body to a carefully planned area. It can damage cancer cells while limiting nearby tissue exposure; fatigue, skin changes, and organ-specific effects may occur.

\medterm{External Capsule} A thin bundle of nerve fibers in the brain near the basal ganglia. It carries signals between the cerebral cortex and deeper brain regions involved in movement, attention, and other functions.

\medterm{External Carotid Artery} A major branch of the neck artery that supplies blood to much of the face, scalp, neck, and parts of the brain coverings through several branches.

\medterm{External Cephalic Version} A procedure in which a clinician applies gentle pressure through the maternal abdomen to turn a fetus from a breech or transverse presentation into a head-down (cephalic) position for vaginal delivery. It is typically performed near term with continuous ultrasound and fetal heart rate monitoring.
\synonyms ECV; Cephalic version.

\medterm{External Ear} The visible part of the ear and the ear canal leading to the eardrum. It collects sound, directs it inward, and helps protect the middle and inner ear.

\medterm{External Fixation} Stabilization of a broken bone or injured limb with pins or wires connected to a frame outside the body. It may be temporary or definitive, especially when swelling, wounds, or contamination prevent internal fixation.

\medterm{External Iliac Artery} A major artery carrying blood from the pelvis toward the leg. It becomes the femoral artery below the groin and gives branches to the lower abdominal wall and pelvic region.

\medterm{External Iliac Vein} A continuation of the femoral vein above the inguinal ligament that joins the internal iliac vein to form the common iliac vein.

\medterm{External Incontinence Devices} Products worn or placed outside the body to collect urine or stool or protect clothing from leakage. Examples include urinary sheaths, absorbent products, and fecal collection systems; skin care and correct fitting help prevent injury.

\medterm{External Jugular Vein} A visible vein running along the side of the neck that drains blood from the scalp and face toward the chest. Its fullness can provide clues about pressure in the right side of the heart.

\medterm{External Nose} The visible nose formed by bone, cartilage, skin, and soft tissue. It shapes the nasal passages, helps warm and filter inhaled air, supports smell, and contributes to facial structure.

\medterm{External Oblique Muscle} The broad, superficial muscle on each side of the abdomen. It compresses abdominal contents and helps bend and rotate the trunk; its lower border contributes to the groin ligament and canal.

\medterm{External Otitis} Inflammation or infection of the ear canal outside the eardrum, often called swimmer’s ear. Moisture, scratching, or skin disease may contribute, causing pain, itching, discharge, swelling, or reduced hearing.

\synonyms
Otitis Externa; Swimmer's Ear

\medterm{External Resorption} Loss of tooth root or other tooth structure caused by a process starting outside the tooth. Injury, pressure, inflammation, or replacement by bone may contribute; treatment depends on location and severity.

\medterm{External Ventricular Drain} A critical neurosurgical bedside procedure in which a flexible plastic catheter is surgically placed through a small burr hole in the skull directly into one of the fluid-filled cavities (lateral ventricles) of the brain. The catheter connects to an external closed drainage and pressure-transducer system. An external ventricular drain serves two vital functions: continuously monitoring intracranial pressure (the pressure inside the skull) and therapeutically draining excess cerebrospinal fluid or blood. It is essential in managing acute hydrocephalus, severe traumatic brain injury, subarachnoid hemorrhage, intraventricular hemorrhage, or intracranial hypertension, preventing brain herniation and death.

\synonyms
EVD; Ventriculostomy; External Ventricular Drainage; Ventricular Catheter

\medterm{Extinction} In learning, the fading of a learned response when reinforcement stops. In neurology, it can mean failing to notice one side of the body or space when both sides are stimulated, suggesting an attention disorder.

\medterm{Extirpation} Complete surgical removal of a tissue, organ, growth, or foreign body. The term describes the extent of removal rather than a particular disease or surgical technique.

\medterm{Extraamniotic} Located outside the amniotic sac surrounding a fetus. The term may describe a procedure, device, or fluid placed between the amnion and nearby tissues during pregnancy care.

\medterm{Extracapsular Cataract Surgery} Cataract surgery that removes the cloudy lens while leaving most of its surrounding capsule in place. An artificial lens is usually inserted into the remaining capsule to restore focusing ability.

\medterm{Extracellular} Located outside cells. Extracellular fluid, proteins, ions, and signals surround cells and help provide nutrients, remove waste, support tissues, and coordinate cell activity.

\medterm{Extracellular Exosome} A tiny membrane-covered vesicle released by a cell that can carry proteins, fats, and genetic material to other cells. Exosomes help cells communicate and are being studied as disease markers and treatment vehicles.

\medterm{Extracellular Fluid} Fluid outside cells, including blood plasma and the fluid between cells. It carries nutrients, hormones, and waste and is regulated mainly by the kidneys, circulation, and hormones.

\medterm{Extracellular Matrix} A network of proteins and other substances outside cells that supports tissues and helps cells attach, communicate, grow, and repair injury. Its composition differs among bone, skin, muscle, and other organs.

\medterm{Extracellular Pathogens} Microorganisms that live and multiply outside human cells. Examples include Streptococcus, Staphylococcus, and some Escherichia coli. They can damage tissues directly or through toxins, and the immune system usually fights them with antibodies and other defenses.

\medterm{Extracellular Region} The area outside cells. It contains fluid, minerals, proteins, and signals that support tissues and influence how nearby cells grow, move, and function.

\medterm{Extracellular Space} The fluid- and matrix-filled area outside cells. Nutrients, waste products, hormones, and signaling molecules move through it, and its composition helps control cell survival, movement, and function.

\medterm{Extracellular Vesicle} A membrane-bound particle released by a cell that transfers molecular cargo between cells and is studied in physiology, disease, biomarkers, and drug delivery.

\medterm{Extracorporeal} Occurring or performed outside the body; extracorporeal circulation and hemodialysis route blood through an external machine.

\medterm{Extracorporeal Circulation} Movement of a person's blood through equipment outside the body and back again. Heart-lung machines, dialysis, and some life-support systems use it to add oxygen, remove waste, or support circulation; bleeding, clotting, infection, and organ injury are possible.

\medterm{Extracorporeal Membrane Oxygenation} An advanced, temporary mechanical life-support technique that drains deoxygenated blood from the body, circulates it through an external artificial membrane lung to add oxygen and remove carbon dioxide, and pumps it back into the bloodstream. It is indicated for severe, life-threatening respiratory failure or cardiogenic shock refractory to standard mechanical ventilation and intensive care. It is utilized in two main configurations: veno-venous (VV-ECMO) for isolated respiratory support, and veno-arterial (VA-ECMO) for combined hemodynamic and respiratory support. Major risks include hemorrhage from systemic anticoagulation, thromboembolism, hemolysis, cannula-site infection, and limb ischemia.

\synonyms
ECMO; Extracorporeal Life Support; ECLS; VV-ECMO; VA-ECMO

\medterm{Extracorporeal Shock Wave Lithotripsy} A treatment that uses focused shock waves from outside the body to break kidney or ureter stones into smaller pieces that can pass in urine. Bruising, bleeding, pain, blockage, or incomplete stone removal can occur.

\medterm{Extracranial} Located or occurring outside the skull or cranial cavity. For example, extracranial arteries supply the scalp and face, while extracranial injury affects tissues outside the skull and brain.

\medterm{Extracranial Hematoma} A collection of blood outside the skull or cranial cavity, usually caused by bleeding into soft tissues after injury.

\medterm{Extract} A concentrated preparation made by separating active substances from a plant, herb, food, or other source. Strength, purity, interactions, and safety vary widely, especially in unregulated supplements.
\synonyms
Botanical Extract; Pharmacological Extract

\medterm{Extraction} Removal of a tooth, tissue, foreign object, or other material from the body using a procedure. The method and risks depend on what is removed, where it is located, and the person’s health.

\medterm{Extra-Dural} Located outside the dura mater, the tough protective membrane around the brain and spinal cord. The term may describe a space, collection of blood, injection, or other structure outside this membrane.

\medterm{Extra-Gastrointestinal Stromal Tumor} A tumor with features of a gastrointestinal stromal tumor that begins outside the digestive tract, often in abdominal supporting tissues. Diagnosis uses imaging, tissue examination, and molecular tests to guide cancer treatment.

\medterm{Extrahepatic Biliary Atresia} Alternate name; see \textit{Biliary Atresia}.

\medterm{Extrahepatic Cholestasis} Reduced or blocked bile flow in the ducts outside the liver. Gallstones, scarring, inflammation, or a tumor may cause yellow skin, dark urine, pale stools, itching, pain, or infection and may need urgent treatment.

\medterm{Extramammary Paget Disease} A rare skin cancer usually affecting the vulva, anus, scrotum, or nearby skin. It often appears as a persistent itchy, sore, red, or scaly patch and requires biopsy to confirm diagnosis and assess for related cancer.

\medterm{Extramammary Paget's Disease} A rare skin cancer involving gland-bearing skin outside the breast, often around the vulva, anus, or penis. It may appear as a persistent itchy, red, scaly, or oozing patch and requires biopsy.

\medterm{Extramedullary Hematopoiesis} The making of blood cells outside the bone marrow, usually in the liver or spleen. It can happen when the bone marrow is damaged, crowded, or unable to make enough blood cells.
\synonyms
EMH; Extramedullary Blood Formation

\medterm{Extranodal} Located outside the lymph nodes. An extranodal disease or lymphoma affects an organ or other tissue, such as the stomach, skin, lung, bone, or brain.

\medterm{Extranodal Marginal Zone Lymphoma} Alternate name; see \textit{MALT Lymphoma}.

\medterm{Extranodal Marginal Zone Lymphoma Of Mucosa-Associated Lymphoid Tissue} A usually slow-growing cancer of B lymphocytes arising in tissues such as the stomach, salivary glands, thyroid, lungs, or eye. Symptoms depend on the organ involved and treatment may include observation, infection treatment, or cancer therapy.

\synonyms
Extranodal Marginal Zone Lymphoma
MALT Lymphoma
Mucosa-Associated Lymphoid Tissue Lymphoma

\medterm{Extranodal Natural Killer T-Cell Lymphoma} An aggressive lymphoma arising from immune cells, often affecting the nose and upper throat. It is strongly associated with Epstein-Barr virus and may cause blockage, bleeding, swelling, fever, or tissue destruction.

\medterm{Extraocular} Outside the eyeball or involving the muscles, nerves, and tissues around it. Extraocular structures move the eye, protect it, support vision, and connect it with the brain and face.

\medterm{Extraocular Muscle} One of the muscles outside the eyeball that moves the eye in different directions. Weakness or nerve injury can cause double vision, eye misalignment, or difficulty looking in a particular direction.

\medterm{Extraocular Muscle Function Testing} Examination of eye movements to assess the muscles and nerves that move the eyes. The person follows targets in different directions; abnormal movement can indicate muscle, nerve, brain, or eye-socket disease.

\medterm{Extraocular Muscles} Six muscles outside each eyeball that move the eyes up, down, inward, outward, and in rotation. Weakness or nerve injury can cause double vision, eye misalignment, or difficulty looking in a particular direction.

\medterm{Extrapulmonary Manifestation Of Tuberculosis} Tuberculosis affecting tissues outside the lungs, such as lymph nodes, bones, kidneys, the brain coverings, abdomen, or reproductive organs. Symptoms vary by site and treatment requires multiple antibiotics.

\medterm{Extrapyramidal Disorder} A disorder of brain circuits that regulate movement, posture, and muscle tone. It may cause stiffness, slowness, tremor, abnormal movements, or restlessness and can result from disease or medicines.

\medterm{Extrapyramidal Symptoms} Movement problems caused by medicines or diseases affecting movement circuits. They include stiffness, tremor, slowed movement, restlessness, sudden spasms, or involuntary movements, commonly after dopamine-blocking medicines.

\medterm{Extrapyramidal System} Interconnected brain structures and pathways, including the basal ganglia, that help control automatic movement, posture, muscle tone, and coordination. Damage or medicine effects can produce abnormal movements.

\medterm{Extrapyramidal Tracts} Nerve pathways outside the main voluntary movement tract that influence posture, muscle tone, automatic movements, and coordination. Disturbance can cause stiffness, tremor, slowness, or involuntary movements.

\medterm{Extraskeletal Chondroma} A benign cartilage-forming lump arising in soft tissue rather than from bone, often in the hands or feet. It may calcify, cause pain or pressure, and sometimes needs removal to confirm the diagnosis.

\medterm{Extraskeletal Mesenchymal Chondrosarcoma} A rare aggressive cancer of soft tissue that contains malignant cartilage. It may arise outside bone and requires specialist biopsy, staging, and combined cancer treatment.

\medterm{Extraskeletal Osteosarcoma} A rare aggressive cancer that produces bone-like material in soft tissue rather than beginning inside bone. It usually affects adults and requires imaging, biopsy, staging, and specialist treatment.

\medterm{Extrasystole} Alternate name; see \textit{Premature Ventricular Contractions}.

\medterm{Extrauterine Pregnancy} A pregnancy implanted outside the normal uterine cavity, most often in a fallopian tube. It can rupture and cause life-threatening internal bleeding; one-sided pain, fainting, or bleeding during pregnancy requires urgent care.

\medterm{Extravasation} Leakage of blood, medicine, urine, or another body fluid from a vessel or duct into nearby tissues. It can cause swelling, pain, inflammation, infection, or tissue death; treatment depends on the substance, location, and severity.

\medterm{Extravascular Hemolysis} The early removal and breakdown of red blood cells mainly in the spleen and liver. It can cause anemia, yellowing of the skin or eyes, and an enlarged spleen.
\synonyms
Reticuloendothelial Hemolysis; Splenic Erythrocyte Clearance

\medterm{Extraventricular Drain} A temporary tube placed through the skull into a fluid-filled space in the brain. It can measure pressure and drain extra brain fluid or blood, especially when fluid buildup is dangerous.

\synonyms
EVD; Ventriculostomy Catheter

\medterm{Extremities} The arms and legs, including their bones, joints, muscles, nerves, blood vessels, and skin. Problems in an extremity may affect movement, sensation, circulation, or function.

\medterm{Extremity} A limb of the body. The upper extremity includes the shoulder, arm, forearm, wrist, and hand; the lower extremity includes the hip, thigh, leg, ankle, and foot.

\medterm{Extremity Angiography} Imaging of blood vessels in an arm or leg, usually after contrast is introduced. It can show narrowing, blockage, aneurysm, abnormal connections, or injury and may sometimes be combined with treatment.

\medterm{Extremity Hardware Removal} Surgical extraction of orthopedic implants such as plates, screws, intramedullary rods, or wires from an arm or leg. Typically indicated for persistent localized pain, implant infection, mechanical irritation, soft-tissue impingement, nonunion, or planned subsequent procedures.

\synonyms
Orthopedic Hardware Removal; Implant Removal; Plate and Screw Removal

\medterm{Extremity X-Ray} An X-ray image of an arm or leg used to look for fractures, joint alignment problems, arthritis, infection, tumors, foreign objects, or other bone abnormalities.

\medterm{Extrinsic} Coming from outside a structure, organ, or person. An extrinsic cause may be an external injury, medicine, infection, toxin, environmental exposure, or other outside influence.

\medterm{Extrinsic Allergic Alveolitis} An immune-mediated lung disease caused by repeated inhalation of organic dusts, molds, bird proteins, or other environmental substances. It can cause cough, breathlessness, fever, and, with ongoing exposure, permanent lung scarring; it is also called hypersensitivity pneumonitis.

\medterm{Extrinsic Factor} An influence from outside a person, tissue, or body process. Examples include medicines, infections, toxins, temperature, nutrition, stress, and environmental exposure; effects depend on the type and amount.

\medterm{Extrinsic Ureteral Obstruction} Blockage of a ureter caused by pressure from outside the ureter, such as a tumor, scar tissue, enlarged lymph nodes, blood-vessel abnormality, or nearby inflammation. It can reduce urine drainage and cause kidney swelling or loss of kidney function.

\synonyms
Extrinsic Ureteral Compression; Secondary Ureteral Obstruction

\medterm{Extubate} To remove a breathing tube after a person can breathe adequately, protect the airway, clear secretions, and maintain safe oxygen levels without the ventilator.

\medterm{Extubation} Removal of a breathing tube from the windpipe after a person can breathe independently, keep the airway open, clear secretions, and maintain adequate oxygen and carbon-dioxide levels.

\medterm{Extubation Criteria} Findings used to judge whether a person can safely have a breathing tube removed. They include adequate breathing and oxygenation, stable circulation, alertness, airway protection, effective cough, and recovery from muscle-paralyzing medicines.

\medterm{Exudate} Protein-rich fluid that leaks from inflamed or injured blood vessels into tissues or a body cavity. Its presence usually indicates inflammation, infection, or tissue injury rather than pressure alone.

\medterm{Exudative Pleural Effusion} A protein-rich collection of fluid between the lung and chest wall caused by local inflammation or injury. Pneumonia, cancer, tuberculosis, and blood clots are possible causes; testing the fluid helps guide treatment.

\medterm{Eye} The organ that detects light and forms vision. It includes the cornea, lens, iris, retina, optic nerve, and supporting tissues within the eye socket; disease can affect sight, movement, comfort, or appearance.

\medterm{Eye Allergies} An allergic reaction affecting the eyes and often the inner eyelids. It usually causes itching, redness, watering, burning, or swollen eyelids and is commonly called allergic conjunctivitis.

\synonyms
Eye Allergy
Ocular Allergy

\medterm{Eye And Orbit Ultrasound} An ultrasound scan of the eye and surrounding socket. It can assess the retina, optic nerve, eye muscles, masses, bleeding, or other structures when the view through the front of the eye is obscured.

\medterm{Eyeball} The roughly spherical organ of vision containing the cornea, sclera, iris, lens, retina, internal fluids, and other tissues. It focuses light onto the retina, which sends visual signals through the optic nerve.

\synonyms
Globe

\medterm{Eye Bank} An organization that recovers, evaluates, stores, and distributes donated corneas or other eye tissue for transplantation, research, or education. Donor screening and tissue testing help protect recipients.

\medterm{Eye, Black} Alternate name; see \textit{Black Eye}.

\medterm{Eye Bleed} Bleeding on, inside, or around the eye. A bright-red painless patch on the white of the eye is often superficial, but bleeding inside the eye, pain, injury, or vision change requires prompt assessment.

\medterm{Eyebrow} A strip of hair above the eye socket that helps divert sweat, water, and particles away from the eye and contributes to facial expression.

\medterm{Eye Cancer} A cancer that begins in the eyelid, surface of the eye, uvea, retina, or tissues around the eye. Warning signs can include a growing or changing spot, persistent redness, vision changes, eye pain, bulging, or a new pupil abnormality; diagnosis usually requires specialist examination and imaging or biopsy.

\medterm{Eye Color} The visible color of the iris, determined mainly by pigment and light scattering. A new, one-sided, uneven, or painful color change may indicate injury, inflammation, medicine effects, or another eye disorder.

\medterm{Eye Development} Formation and maturation of the eyes, optic nerves, and visual pathways before and after birth. Disruption can cause structural abnormalities, reduced vision, misalignment, or delayed development of visual skills.

\medterm{Eye Disease} Any disorder affecting the eye or visual pathways, including refractive errors, cataract, glaucoma, retinal disease, infection, inflammation, and tumors. Symptoms such as sudden vision loss, severe pain, flashes, or a curtain-like shadow require urgent assessment.

\medterm{Eye, Diseases Of} Diseases affecting the eye or visual pathways, including refractive problems, cataracts, glaucoma, retinal disease, infection, inflammation, injury, and tumors. Symptoms may involve vision, pain, redness, discharge, or eye movement.

\medterm{Eye Disorder} A disease or structural problem affecting the eye or visual pathways. It may cause changed vision, pain, redness, discharge, double vision, field loss, or other findings and may require routine or urgent assessment.

\medterm{Eye Disorders} Diseases or structural problems affecting the eye, optic nerve, or visual pathways. They may cause blurred or lost vision, pain, redness, discharge, double vision, or field loss and can require routine or emergency assessment.

\medterm{Eye Drops} Liquid medicines, lubricants, or diagnostic preparations placed on the eye surface. They may relieve dryness or treat disease; contaminated bottles, shared drops, or incorrect use can cause injury or infection.

\medterm{Eye Emergencies} Conditions needing immediate eye assessment include sudden vision loss, chemical exposure, penetrating injury, severe pain, new neurologic symptoms, or rapidly spreading infection. Do not rub an injured eye or press on it.

\medterm{Eye Floaters} Alternate name; see \textit{Floaters}.

\medterm{Eyeglass} A lens mounted in a frame to correct focusing problems such as nearsightedness, farsightedness, astigmatism, or presbyopia. The prescription is selected after measuring vision and the eyes.

\medterm{Eyeglasses and Sunglasses} Alternate name; see \textit{Eyeglasses, Sunglasses, and Magnifiers}.

\medterm{Eyeglasses, Sunglasses, And Magnifiers} Eyeglasses correct focusing errors, sunglasses reduce glare and ultraviolet exposure, and magnifiers enlarge nearby targets. Correct prescription, fit, lens protection, and the underlying eye condition determine the appropriate aid.

\synonyms
Corrective Eyewear; Optical Aids

\medterm{Eye Health} Maintaining vision and preventing avoidable eye disease through appropriate examinations, injury and ultraviolet protection, control of diabetes and blood pressure, safe contact-lens care, and prompt attention to new symptoms.

\medterm{Eye Hemorrhage} Bleeding on, inside, or around the eye. A small surface bleed may be painless, while bleeding in the retina, vitreous, or eye socket can impair vision or cause pressure and needs medical assessment.

\medterm{Eye Infection} Infection of the eyelid, conjunctiva, cornea, internal eye, or surrounding tissues caused by bacteria, viruses, fungi, or parasites. Pain, light sensitivity, reduced vision, or swelling around the eye needs prompt care.

\medterm{Eye Injury} Damage to the eye or nearby tissues from blunt or penetrating trauma, chemicals, foreign objects, radiation, or heat. Vision change, severe pain, bleeding, or an embedded object requires urgent assessment.

\medterm{Eye Irritation} Burning, itching, redness, or watering caused by dryness, allergy, smoke, chemicals, infection, or a foreign particle. Chemical exposure requires immediate prolonged rinsing with clean water and urgent medical advice.

\medterm{Eyelash} A hair growing from the eyelid edge that helps shield the eye from particles and triggers blinking when touched. Misrouted or infected lashes can irritate the eye surface.

\medterm{Eyelashes} Short curved hairs along the eyelid edges that help protect the eye from dust, moisture, and bright light. Inward-growing lashes can rub the cornea and require treatment.

\medterm{Eyelid} A movable fold of skin, muscle, and glands that protects the eye, spreads tears across its surface, and helps regulate light. Swelling, malposition, or weakness can impair protection and vision.

\medterm{Eyelid Bump} A localized eyelid swelling, often a painful infected gland called a stye or a firmer blocked oil gland called a chalazion. Warm compresses may help; persistent, recurrent, or vision-affecting bumps need examination.

\medterm{Eyelid Cancer} Cancer arising in the skin or other tissues of the eyelid, most commonly basal cell carcinoma. A persistent or changing lump, sore, bleeding area, loss of eyelashes, or eyelid distortion needs ophthalmic or dermatologic evaluation.

\medterm{Eyelid Coloboma} A birth defect causing a notch or missing section of the eyelid. It may leave the eye exposed, impair closure, and occur with other eye or facial abnormalities requiring specialist assessment.

\medterm{Eyelid Disorder} A disease or structural problem of the eyelid, such as infection, inflammation, a blocked gland, malposition, weakness, or tumor. It may cause pain, swelling, poor eye protection, irritation, or visual problems.

\medterm{Eyelid Drooping} Alternate name; see \textit{Ptosis}.

\medterm{Eyelid Inflammation} Alternate name; see \textit{Blepharitis}.

\medterm{Eyelid Lift} Alternate name; see \textit{Blepharoplasty}.

\medterm{Eyelid Neoplasm} A benign, precancerous, or cancerous growth of the eyelid. A persistent, enlarging, changing, bleeding, crusted, or ulcerated lesion needs examination by an eye or skin specialist.

\medterm{Eyelids} Movable folds of skin, muscle, and glands that protect the eyes, spread tears, and help control light. Their position and movement are important for keeping the eye surface moist and safe.

\medterm{Eyelid Swelling} Puffiness of one or both eyelids caused by allergy, infection, inflammation, injury, a blocked gland, or fluid accumulation. Fever, severe pain, impaired eye movement, bulging eye, or vision change may indicate an emergency.

\medterm{Eyelid Twitch} Alternate name; see \textit{Eye Twitching}.

\medterm{Eye Melanoma} A cancer arising from pigment-producing cells in the eye, most often in the uvea. It may cause blurred vision, flashes, a dark spot, or no symptoms; treatment depends on size, location, spread, and vision potential.

\synonyms
Uveal Melanoma; Ocular Melanoma; Choroidal Melanoma

\medterm{Eye Movement} Coordinated movement of the eyes that directs gaze and keeps images aligned. It depends on the eye muscles, their nerves, and brain control; abnormal movement may cause double vision or imbalance.

\medterm{Eye Muscle Repair} Treatment to correct a damaged or misaligned eye muscle or its attachment, usually by eye-muscle surgery. It may improve alignment or double vision but does not always restore normal movement and carries surgical risks.

\medterm{Eye Neoplasm} A benign or malignant growth arising in the eye or surrounding tissues. Symptoms and treatment depend on the tissue involved, size, location, effect on vision, and whether the growth has spread.

\medterm{Eye Pain} Pain in or around the eye caused by surface irritation, infection, injury, inflammation, glaucoma, or disease behind the eye. Severe pain with vision loss, light sensitivity, nausea, or trauma requires urgent assessment.

\medterm{Eye Prosthesis} An artificial eye or implant used to restore appearance after an eye is lost or severely damaged. It generally does not restore sight; fitting, cleaning, and periodic specialist care are needed.

\medterm{Eye Redness} Redness caused by enlarged or inflamed blood vessels on or inside the eye. Irritation and conjunctivitis are common causes, but pain, light sensitivity, discharge, injury, or vision change may indicate serious disease.

\medterm{Eyesight} Alternate name; see \textit{Vision}.

\medterm{Eyestrain} Tired, dry, sore, or uncomfortable eyes, sometimes with blurred vision or headache, after prolonged visual effort. Screen use, glare, poor lighting, dry eye, or an uncorrected focusing problem may contribute; rest often helps.

\medterm{Eye Trauma} Injury ranging from a surface scratch to a penetrating wound or ruptured eyeball. Rinse chemical exposure immediately, avoid pressure or rubbing, shield penetrating injuries, and obtain urgent care for vision change or severe pain.

\medterm{Eye Twitching} Repeated involuntary eyelid or eye-area spasms, commonly linked to fatigue, stress, caffeine, irritation, or eye strain. Persistent spasms, facial weakness, eye closure, or other neurologic symptoms need assessment.

\medterm{Ezetimibe} An oral medicine that reduces absorption of cholesterol from the intestine. It lowers blood cholesterol alone or with another medicine and may cause diarrhea, abdominal discomfort, or rarely liver or muscle problems.
